% ucft-final.tex — standalone master manuscript (no \input dependencies).
% All theorems, proofs, postulates, and appendices are inlined.
% Regenerated from .ucft-build/unit-*.tex sources.

\documentclass[11pt]{article}
\usepackage[T1]{fontenc}
\usepackage[utf8]{inputenc}
\usepackage[margin=1in]{geometry}
\usepackage{amsmath,amssymb,amsfonts,amsthm,mathtools}
\usepackage{mathrsfs}
\usepackage{graphicx}
\usepackage{xcolor}
\usepackage{textcomp}
\usepackage{booktabs}
\usepackage{enumitem}
\usepackage{array}
\usepackage{hyperref}

\hypersetup{colorlinks=true,linkcolor=blue!55!black,citecolor=blue!55!black,urlcolor=blue!55!black}
\allowdisplaybreaks
\renewcommand{\arraystretch}{1.18}

\theoremstyle{plain}
\newtheorem{theorem}{Theorem}[section]
\newtheorem{lemma}[theorem]{Lemma}
\newtheorem{corollary}[theorem]{Corollary}
\newtheorem{proposition}[theorem]{Proposition}
\newtheorem{conjecture}[theorem]{Conjecture}
\theoremstyle{definition}
\newtheorem{definition}[theorem]{Definition}
\newtheorem{axiom}[theorem]{Axiom}
\newtheorem{postulate}[theorem]{Postulate}
\newtheorem{construction}[theorem]{Construction}
\newtheorem{indicative}[theorem]{Indicative}
\theoremstyle{remark}
\newtheorem{remark}[theorem]{Remark}
\newtheorem{example}[theorem]{Example}

\providecommand*{\sectionautorefname}{Section}
\providecommand*{\subsectionautorefname}{Section}
\providecommand*{\theoremautorefname}{Theorem}
\providecommand*{\lemmaautorefname}{Lemma}
\providecommand*{\corollaryautorefname}{Corollary}
\providecommand*{\propositionautorefname}{Proposition}
\providecommand*{\conjectureautorefname}{Conjecture}
\providecommand*{\definitionautorefname}{Definition}
\providecommand*{\axiomautorefname}{Axiom}
\providecommand*{\postulateautorefname}{Postulate}
\providecommand*{\constructionautorefname}{Construction}
\providecommand*{\indicativeautorefname}{Indicative}
\providecommand{\cref}[1]{\autoref{#1}}

% Core notation.
\newcommand{\SU}[1]{\mathrm{SU}(#1)}
\newcommand{\SO}[1]{\mathrm{SO}(#1)}
\newcommand{\U}[1]{\mathrm{U}(#1)}
\newcommand{\E}[1]{\mathrm{E}_{#1}}
\newcommand{\SL}{\mathrm{SL}}
\newcommand{\Spin}{\operatorname{Spin}}
\newcommand{\Aut}{\operatorname{Aut}}
\newcommand{\End}{\operatorname{End}}
\newcommand{\Sym}{\operatorname{Sym}}
\newcommand{\rank}{\operatorname{rank}}
\newcommand{\diag}{\operatorname{diag}}
\newcommand{\Hess}{\operatorname{Hess}}
\newcommand{\Tr}{\operatorname{Tr}}
\newcommand{\Str}{\operatorname{STr}}
\newcommand{\dd}{\mathrm{d}}
\newcommand{\R}{\mathbb{R}}
\newcommand{\C}{\mathbb{C}}
\newcommand{\Oct}{\mathbb{O}}
\newcommand{\M}{\mathcal{M}}
\newcommand{\jalg}{J_3(\mathbb{O})}
\newcommand{\SOTen}{\SO{10}}
\newcommand{\UOne}{\U{1}}
\newcommand{\Uone}{\U{1}}
\newcommand{\ESix}{\E{6}}
\newcommand{\EIII}{\mathrm{EIII}}
\newcommand{\SOTenXUOne}{\SOTen \times \UOne}
\newcommand{\MCoset}{\frac{\ESix}{\SOTenXUOne}}
\newcommand{\MCosetInline}{\ESix/(\SOTenXUOne)}
\newcommand{\ESixBreaking}{\ESix \to \SOTenXUOne}
\newcommand{\Herm}[2]{\mathrm{Herm}_{#1}(#2)}
\newcommand{\chr}[1]{\mathrm{char}(#1)}
\newcommand{\Rep}[1]{\mathbf{#1}}
\newcommand{\Repbar}[1]{\overline{\mathbf{#1}}}
\newcommand{\Pone}{\textnormal{P1}}
\newcommand{\Ptwo}{\textnormal{P2}}
\newcommand{\Pthree}{\textnormal{P3}}
\newcommand{\Pfour}{\textnormal{P4}}
\newcommand{\Pfive}{\textnormal{P5}}
\newcommand{\Psix}{\textnormal{P6}}
\newcommand{\MGUT}{M_{\mathrm{GUT}}}
\newcommand{\ghat}{\hat g}
\newcommand{\yhat}{\hat y}
\newcommand{\clock}{\phi^{0}}
\newcommand{\ket}[1]{\lvert #1\rangle}
\newcommand{\dket}[1]{\lvert #1\rangle\!\rangle}
\newcommand{\dbra}[1]{\langle\!\langle #1\rvert}
\newcommand{\sm}{\hspace{.065em}}
\newcommand{\s}{\hspace{.1em}}
\newcommand{\smexp}[1]{^{\sm #1}}
\newcommand{\set}[2]{\{#1\,:\,#2\}}

\title{Universal Coset Field Theory:\\[0.4em]
\large From Axioms to the Standard Model and Gravitation}
\author{Brandon Belna}
\date{June 6, 2026}

\begin{document}
\maketitle

\begin{abstract}
We formulate Universal Coset Field Theory (UCFT) on a foundational graph--Stone
substrate and the exceptional Jordan algebra $J_3(\mathbb O)$ with compact $E_6$
structure. A pre-axiomatic layer (BAMLP: Barrier-Aware Metric--Laguerre
Partitioning) posits a weighted cosmic graph, a division-algebra ladder
$0\to\mathbb{R}\to\mathbb{C}\to\mathbb{H}\to\mathbb{O}$, and an exceptional
cascade $0\to E_8\to E_8\times E_8\to E_6\to\cdots\to 0$. Two axioms specify the
order parameter and quantum kinematics; six postulates \textbf{P1--P6} supply the
$E_6$-invariant potential, vacuum manifold, Lorentzian soldering, chiral matter
content, three generations, and GUT--Yukawa sector. Under these hypotheses we
prove that the global minimum of the potential is the rank-one orbit
$E_6/(\mathrm{Spin}(10)\times U(1))$ (EIII), with $32$ Goldstone modes and
positive-definite coset metric. Postulate~\textbf{P3} yields four-dimensional
Lorentzian spacetime; one-loop fluctuations induce a positive Planck mass
$M_{\mathrm{Pl}}^2=(28/6)\,f^2\mu^2\log(\Lambda^2/\mu^2)/(16\pi^2)$ and an
asymptotically free $\mathrm{Spin}(10)\times U(1)$ sector with
$b_0=12$ for three generations. With \textbf{P4--P6} the theory embeds the
Standard Model, including neutrino masses via the type-I seesaw. Conjectural
results include a gravitational ultraviolet fixed point
$\kappa_\star=96\pi^2/5$ and heterotic string correspondences. Octonionic
algebra, soldering, representation theory, and heat-kernel computations are
collected in appendices.
\end{abstract}

\noindent\textbf{Keywords:}
BAMLP; Stone spaces; graph universe; exceptional ladder; exceptional Jordan
algebra; $E_6$ coset; induced gravity; $\mathrm{Spin}(10)$ grand unification;
Standard Model embedding.

\tableofcontents
\clearpage
%----------------------------------------------------------------------%
% Part: Foundations and the Master Chain
%----------------------------------------------------------------------%
\part{Foundations and the Master Chain}
% --- begin inlined: unit-bamlp.tex ---
%======================================================================%
\section{Foundations: Graphs, Stone Spaces, and the Exceptional Ladder}
\label{sec:bamlp:main}
%======================================================================%

Universal Coset Field Theory (UCFT) is not formulated in a vacuum. The present
section records the \emph{pre-axiomatic} substrate from which the $E_6$/Albert
arena is meant to descend: a combinatorial graph description of configuration
space, a Stone-spectral ladder through the normed division algebras, and an
exceptional-group cascade terminating in the Standard Model. The assignment
mechanism that binds these layers is \textbf{BAMLP} (Barrier-Aware
Metric--Laguerre Partitioning): a capacity-balanced dual on a metric--measure
space that partitions workload into additively weighted cells.

Throughout we classify claims as \textbf{theorem} (standard mathematics),
\textbf{postulate} (added structural identification), or
\textbf{conjecture} (motivated physics correspondence).

%----------------------------------------------------------------------%
\subsection{The cosmic graph as primitive data}
\label{sec:bamlp:graph}
%----------------------------------------------------------------------%

\begin{postulate}[Graph arena]\label{post:bamlp:graph}
Physical configuration space is presented first as a locally finite weighted
graph $\mathcal{G}=(V,E,w)$ together with a finite or $\sigma$-finite workload
measure $\mu$ on $V$ (or on edges/nodes after projection). Operational cost is
a graph-geodesic (or barrier-aware) distance $c_i(x)$ from site $x$ to
territory generator $z_i\in V$. The \emph{graph description of the universe} is
this triple $(\mathcal{G},\mu,\{c_i\})$; continuous geometry and Lie-group
structure are derived, not assumed at this level.
\end{postulate}

In the BAMLP framework~\cite{BAMLP:2026}, territory assignment is the
deterministic rule
\begin{equation}
  x \;\longmapsto\; \arg\min_i\bigl(c_i(x)-\omega_i\bigr),
  \label{eq:bamlp:assign}
\end{equation}
with additive capacity weights $\omega_i$ adjusted until cell masses match
prescribed targets $\tau_i$. On a graph, the cells are additively weighted
graph-Voronoi regions; capacities $\omega_i$ are Lagrange multipliers for the
mass constraints. UCFT interprets this dual not as a political districting
algorithm but as the \emph{selection principle} that balances degrees of
freedom across symmetry channels.

\begin{theorem}[BAMLP capacity dual (finite, bounded cost)]\label{thm:bamlp:dual}
Let $c_i:V\to\mathbb{R}$ be bounded below and $\mu$ a finite measure on $V$.
For target masses $\tau_i>0$ with $\sum_i\tau_i=\mu(V)$, there exist weights
$\omega\in\mathbb{R}^K$ such that the partition induced by
\eqref{eq:bamlp:assign} satisfies $|\mu(\mathrm{cell}_i)-\tau_i|\le\varepsilon$
for any $\varepsilon>0$ under capacity balancing. The map $\omega\mapsto$ cell
masses is concave; in the squared-Euclidean specialization the cells are
Laguerre (power) diagrams.
\end{theorem}

\begin{remark}[UCFT reading]\label{rem:bamlp:ucft}
The $32$ clock fields of UCFT are hypothesized to arise as the \emph{boundary
data} of a BAMLP partition on a configuration graph whose generators label
symmetry-breaking channels. The rank-one EIII vacuum is then the unique
capacity-balanced cell at the $E_6$ level, not an arbitrary minimum of an
unconstrained potential. This identification is \textbf{conjectural}; the
$E_6$-invariant potential of postulate~\textnormal{P1} supplies the proved
dynamics once the arena is fixed.
\end{remark}

%----------------------------------------------------------------------%
\subsection{Stone spaces and the division-algebra ladder}
\label{sec:bamlp:stone}
%----------------------------------------------------------------------%

Stone duality links zero-dimensional compact Hausdorff spaces (Stone spaces)
with Boolean algebras. The profinite spectrum of a graph's clique complex
supplies the combinatorial skeleton on which the normed division-algebra tower
is mounted.

\begin{definition}[Stone spectrum of a graph]\label{def:bamlp:stone}
Let $\mathcal{G}$ be a graph and $\mathcal{B}(\mathcal{G})$ the Boolean algebra
generated by its vertex cuts (sets of vertices separated by removing edge
subsets). The \emph{Stone space} is $S(\mathcal{G})=\mathrm{Spec}(\mathcal{B}
(\mathcal{G}))$, a compact, Hausdorff, totally disconnected space. A morphism of
graphs inducing a Boolean homomorphism gives a continuous map of Stone spaces.
\end{definition}

\begin{theorem}[Stone duality]\label{thm:bamlp:stone}
The category of Stone spaces is contravariantly equivalent to the category of
Boolean algebras. In particular, finite Boolean algebras correspond to finite
discrete Stone spaces; profinite limits correspond to inverse systems of finite
graphs.
\end{theorem}

\begin{theorem}[Hurwitz--Radon division algebra ladder]\label{thm:bamlp:hurwitz}
The only finite-dimensional \emph{normed} division algebras over $\mathbb{R}$
are, up to isomorphism, the tower
\begin{equation}
  0 \;\longrightarrow\; \mathbb{R} \;\longrightarrow\; \mathbb{C}
  \;\longrightarrow\; \mathbb{H} \;\longrightarrow\; \mathbb{O},
  \label{eq:bamlp:divalg}
\end{equation}
of dimensions $1,2,4,8$. The passage $\mathbb{O}\to\mathbb{S}$ (sedenions)
fails: $\mathbb{S}$ is not a division algebra. The Albert algebra
$J_3(\mathbb{O})$ and the cubic norm $N(X)=\det X$ are the canonical
exceptional-Jordan packaging of the terminal algebra $\mathbb{O}$ in this tower.
\end{theorem}

\begin{postulate}[Stone--division lift]\label{post:bamlp:lift}
The profinite Stone limit of the cosmic graph $\mathcal{G}$ carries a
sequential lift through the division-algebra rungs
\eqref{eq:bamlp:divalg}. The zero object $0$ at the base denotes the trivial
Boolean/Stone sector (no observable localisation); each rung adds a
multiplicative structure compatible with the graph's capacity measure:
\begin{center}
\small
\begin{tabular}{@{}clcl@{}}
\toprule
Rung & Algebra & $\dim$ & Automorphism / structure group (typical) \\
\midrule
$0$ & $\{0\}$ & $0$ & trivial \\
$\mathbb{R}$ & reals & $1$ & order, $\mathrm{Aut}=\{1\}$ \\
$\mathbb{C}$ & complexes & $2$ & $\mathrm{U}(1)$ phase \\
$\mathbb{H}$ & quaternions & $4$ & $\mathrm{SU}(2)$ \\
$\mathbb{O}$ & octonions & $8$ & $G_2$ \\
$J_3(\mathbb{O})$ & Albert & $27$ & $F_4\subset E_6$ \\
\bottomrule
\end{tabular}
\end{center}
The lift is functorial with respect to graph refinements: each BAMLP capacity
rebalancing that refines $\mathcal{G}$ induces a compatible map of Stone
spectra and preserves the division-algebra rung labels.
\end{postulate}

%----------------------------------------------------------------------%
\subsection{The exceptional ladder}
\label{sec:bamlp:exceptional}
%----------------------------------------------------------------------%

Above the octonionic/Jordan rung, the exceptional Lie algebras form a second
ladder linked to heterotic string theory and grand unification. UCFT occupies
the $E_6$ rung; the full conjectured cascade is:

\begin{equation}
  0 \;\to\; E_8 \;\to\; E_8\times E_8 \;\to\; E_6
  \;\to\; \mathrm{Spin}(10)\times U(1) \;\to\; SU(5)
  \;\to\; SU(3)_c\times SU(2)_L\times U(1)_Y \;\to\; SU(3)_c\times U(1)_{\mathrm{em}}
  \;\to\; 0.
  \label{eq:bamlp:exc}
\end{equation}

\begin{theorem}[Steps proved or embedded in UCFT]\label{thm:bamlp:proved}
The following steps of \eqref{eq:bamlp:exc} are established in the body of this
manuscript (under postulates~\textnormal{P1--P6} where indicated):
\begin{enumerate}[label=\textup{(\roman*)}]
  \item $E_6\to\mathrm{Spin}(10)\times U(1)$: vacuum selection
    (\autoref{thm:potential:T4}).
  \item $\mathrm{Spin}(10)\times U(1)\to SU(5)\times U(1)_X\to\cdots\to
    SU(3)_c\times U(1)_{\mathrm{em}}$: GUT descent (\autoref{thm:sm:descent}).
  \item $E_8\times E_8$ central charges $(c_L,c_R)=(16,10)$ and
    $h^1(V_{\mathbf{16}})=3$: heterotic \emph{correspondence}
    (\autoref{sec:strings:main}).
\end{enumerate}
\end{theorem}

\begin{conjecture}[BAMLP--exceptional closure]\label{conj:bamlp:closure}
The entire ladder \eqref{eq:bamlp:exc}, including the initial $0\to E_8$ lift
from the Stone--division tower and the final $SU(3)_c\times U(1)_{\mathrm{em}}\to
0$ collapse, is the capacity-balanced partition of a single cosmic graph
$(\mathcal{G},\mu)$ under BAMLP with generators labelled by the groups in
\eqref{eq:bamlp:exc}. In particular:
\begin{itemize}
  \item $E_8$ is the maximal symmetry of the $\mathbb{O}$-valued graph
    measure before Jordan projection;
  \item $E_8\times E_8$ is the product graph for left/right heterotic sectors;
  \item $E_6$ is the rank-one balanced cell selected by
    \eqref{eq:bamlp:assign} with $E_6$-invariant capacity weights;
  \item the terminal $0$ is the trivial gauge group after electroweak breaking.
\end{itemize}
\end{conjecture}

%----------------------------------------------------------------------%
\subsection{Logical position in UCFT}
\label{sec:bamlp:position}
%----------------------------------------------------------------------%

The dependency order becomes:
\begin{center}
\fbox{\parbox{0.92\linewidth}{
\small
\textbf{Graph} $(\mathcal{G},\mu)$ + \textbf{Stone lift}
$\Rightarrow$ \textbf{Division tower}
\eqref{eq:bamlp:divalg}\\
$\Downarrow$ \textbf{BAMLP capacity dual} \eqref{eq:bamlp:assign}\\
$\Downarrow$ \textbf{Exceptional ladder} \eqref{eq:bamlp:exc} (hits $E_6$)\\
$\Downarrow$ \textbf{UCFT axioms I--II} + \textbf{P1--P6}
$\Rightarrow$ Standard Model + gravitation
}}
\end{center}

\begin{remark}[Epistemic status]\label{rem:bamlp:status}
The graph postulate (\autoref{post:bamlp:graph}), Stone--division lift
(\autoref{post:bamlp:lift}), and exceptional closure
(\autoref{conj:bamlp:closure}) are \emph{foundational conjectures}. They are
not used as hypotheses in any theorem of \autoref{sec:potential:main}--%
\autoref{sec:sm:main}. Their role is to supply a \emph{generative narrative}:
if the universe is a capacity-balanced graph partition, the $E_6$/Albert coset
and its proved SM embedding are the natural intermediate rungs, not ad hoc
choices. Deriving~\textnormal{P1--P6} from BAMLP remains open work.
\end{remark}

%----------------------------------------------------------------------%
\subsection{Open problems}
\label{sec:bamlp:open}
%----------------------------------------------------------------------%

\begin{enumerate}
  \item Construct an explicit cosmic graph $\mathcal{G}$ whose BAMLP partition
    at the $E_6$ level reproduces the EIII vacuum manifold
    $\dim_{\mathbb{R}}=32$.
  \item Prove that the Stone profinite limit of $\mathcal{G}$ inherits an
    $F_4$ action descending to the compact $E_6$ structure group on
    $J_3(\mathbb{O})$.
  \item Derive the $E_6$-invariant potential~\eqref{eq:intro:V} as the
    capacity-dual functional of a graph measure, not as a postulate.
  \item Extend the ladder \eqref{eq:bamlp:exc} below $E_8$ with a controlled
    $E_8\to E_6$ breaking mechanism compatible with the heterotic
    $E_8\times E_8$ product.
\end{enumerate}% --- end inlined: unit-bamlp.tex ---
% --- begin inlined: unit-master-chain.tex ---
%======================================================================%
\section{Overview of the Construction}
\label{sec:master:chain}
%======================================================================%

This section summarizes the logical structure of the paper. The foundational
layer (\autoref{sec:bamlp:main}) posits a graph description of the universe,
a Stone-spectral ladder through the division algebras, and BAMLP capacity
balancing on the exceptional cascade
\eqref{eq:bamlp:exc}. The axioms and postulates~\textnormal{P1--P6} then yield a
four-dimensional quantum field theory whose low-energy limit contains the
Standard Model coupled to Einstein--Hilbert gravitation.

\subsection{Axioms}
\label{sec:master:axioms}

\begin{axiom}[I: Order parameter]\label{ax:master:I}
The field $\Phi\in\mathbf{27}=J_3(\mathbb O)$ transforms under compact $E_6$.
The origin $\Phi=0$ is a strict local maximum of the potential of~\textnormal{P1}
(\autoref{thm:potential:T5}); dynamics select the rank-one orbit
(\autoref{thm:potential:T6}).
\end{axiom}

\begin{axiom}[II: Quantum kinematics]\label{ax:master:II}
States admit unitary evolution. In the semiclassical regime this evolution is
realized through a Page--Wootters clock $\clock$ and the Hamiltonian constraint
of induced gravity (\autoref{sec:spacetime:time}).
\end{axiom}

\subsection{Postulates \textnormal{P1--P6}}
\label{sec:master:postulates}

\begin{center}
\small
\begin{tabular}{@{}clp{0.54\linewidth}@{}}
\toprule
 & Postulate & Content \\
\midrule
P1 & Potential & $E_6$-invariant $V(N,X^\#)$; rank-one minimum \\
P2 & Vacuum & $\mathcal M=\mathrm{EIII}=E_6/(\mathrm{Spin}(10)\times U(1))$;
proved from P1 \\
P3 & Soldering & Four clock fields $\to$ Lorentzian $(1,3)$ via $H_2(\mathbb C)$ \\
P4 & Matter & One chiral $\mathbf{16}$ per family \\
P5 & Generations & $SU(3)_F$ on $J_3(\mathbb O)\otimes\mathbb C^3$ \\
P6 & Yukawa & GUT Higgs sector; type-I seesaw \\
\bottomrule
\end{tabular}
\end{center}

\subsection{Logical dependencies}
\label{sec:master:map}

\begin{table}[t]
  \centering
  \small
  \renewcommand{\arraystretch}{1.22}
  \begin{tabular}{@{}clll@{}}
    \toprule
    Step & Input & Output & Reference \\
    \midrule
    $-3$ & Graph $\mathcal{G}$, measure $\mu$ & Stone spectrum, Boolean cuts &
      \autoref{sec:bamlp:graph} \\
    $-2$ & Stone lift & $0\to\mathbb{R}\to\mathbb{C}\to\mathbb{H}\to\mathbb{O}$ &
      \autoref{thm:bamlp:hurwitz} \\
    $-1$ & BAMLP dual & Exceptional ladder to $E_6$ &
      \autoref{conj:bamlp:closure} \\
    0 & Axioms I--II & $J_3(\mathbb O)$, compact $E_6$ & \S\ref{sec:albert:main} \\
    1 & P1 & Rank-one minimum & \autoref{thm:potential:T4} \\
    2 & P2 & $\mathcal M=\mathrm{EIII}$, 32 modes & \autoref{sec:coset:main} \\
    3 & Geometry & $\mathfrak m=\mathbf{16}_{-3}\oplus\overline{\mathbf{16}}_{+3}$ & \autoref{thm:coset:irrep} \\
    4 & Maurer--Cartan & Composite $\mathrm{Spin}(10)\times U(1)$ & \autoref{thm:gauge:composite} \\
    5 & P3 & Lorentzian $M^4$; spectrum $28+4$ & \autoref{sec:spacetime:main} \\
    6 & Axiom II, P3 & Emergent time & \autoref{sec:spacetime:time} \\
    7 & Massive modes & $M_{\mathrm{Pl}}^2>0$ & \autoref{sec:gravity:main} \\
    8 & Matter content & Asymptotic freedom; $b_0(3)=12$ & \autoref{thm:sm:b0} \\
    9 & FRG & $\kappa_\star=96\pi^2/5$ & \autoref{sec:frg:main} \\
    10 & P4--P6 & Standard Model & \autoref{sec:sm:main} \\
    11 & -- & Heterotic correspondence & \autoref{sec:strings:main} \\
    \bottomrule
  \end{tabular}
  \caption{Main logical dependencies.}
  \label{tab:master:chain}
\end{table}

\subsection{Principal parameters}
\label{sec:master:numbers}

\begin{center}
\small
\begin{tabular}{@{}ll@{}}
\toprule
Quantity & Value \\
\midrule
Vacuum & $\mathcal M=E_6/(\mathrm{Spin}(10)\times U(1))$, $\dim_{\mathbb R}=32$ \\
Coset tangent & $\mathfrak m=\mathbf{16}_{-3}\oplus\overline{\mathbf{16}}_{+3}$ \\
Spectrum & $28$ massive scalars $+$ $4$ gauge modes \\
$M_{\mathrm{Pl}}^2$ &
  $\tfrac{28}{6}\,f^2\mu^2\log(\Lambda^2/\mu^2)/(16\pi^2)$ \\
$b_0$ & $26$ (one generation); $12$ (three generations) \\
$\kappa_\star$ & $96\pi^2/5$ (conjectural) \\
\bottomrule
\end{tabular}
\end{center}

\subsection{Summary}
\label{sec:master:summary}

\begin{theorem}[Embedding]\label{thm:master:embedding}
Under Axioms~\autoref{ax:master:I}--\autoref{ax:master:II} and
postulates~\textnormal{P1--P6}, UCFT defines a four-dimensional effective field
theory with proved vacuum structure, induced gravitation, an asymptotically free
gauge sector, and an embedded Standard Model. String-theoretic correspondences
and the gravitational ultraviolet fixed point are conjectural.
\end{theorem}

\begin{proof}
See \autoref{tab:master:chain} and \autoref{thm:intro:embedding}.
\end{proof}

\clearpage% --- end inlined: unit-master-chain.tex ---
% --- begin inlined: unit-intro.tex ---
%======================================================================%
\section{Introduction}
\label{sec:intro:main}
%======================================================================%

\subsection{Motivation}
\label{sec:intro:motivation}

The present work rests on a three-tier foundation. At the deepest level,
\autoref{sec:bamlp:main} posits a \emph{graph description of the universe}---
a weighted configuration graph with a Stone-spectral lift through the division
algebras $\mathbb{R}\to\mathbb{C}\to\mathbb{H}\to\mathbb{O}$---and an
exceptional cascade $0\to E_8\to E_8\times E_8\to E_6\to\cdots\to 0$ selected
by BAMLP capacity balancing. The middle tier is UCFT proper: exceptional Lie
algebras and Jordan algebras provide the arena for unification. The exceptional Jordan (Albert) algebra
$J_{3}(\mathbb{O})$---the algebra of $3\times 3$ Hermitian octonionic
matrices---is the unique finite-dimensional formally real Jordan algebra that is
not special. It has dimension $27$, automorphism group $F_4$, and is preserved
by the exceptional group $E_6$ through its cubic norm $N(X)=\det X$ (up to a
real character). Under the maximal subgroup $H=\mathrm{Spin}(10)\times U(1)$,
the fundamental representation branches as
\begin{equation}
  \mathbf{27} \;=\; \mathbf{1}_{+4}\,\oplus\,\mathbf{10}_{-2}\,\oplus\,\mathbf{16}_{+1},
  \label{eq:intro:27branch}
\end{equation}
which contains a single chiral spinor $\mathbf{16}$ of $\mathrm{Spin}(10)$---one
Standard Model generation together with a right-handed neutrino. The adjoint
$\mathbf{78}$ decomposes as
\begin{equation}
  \mathbf{78} \;=\; \mathbf{45}_{0}\,\oplus\,\mathbf{1}_{0}\,\oplus\,
  \mathbf{16}_{-3}\,\oplus\,\overline{\mathbf{16}}_{+3}
  \qquad (\text{under } H=\mathrm{Spin}(10)\times U(1)),
  \label{eq:intro:78branch}
\end{equation}
exhibiting the coset tangent space $\mathfrak{m}=\mathbf{16}_{-3}\oplus\overline{\mathbf{16}}_{+3}$
of real dimension $32$. The appearance of a viable grand-unified gauge group
together with the chiral matter content of one family motivates the present
construction.

\subsection{Axioms and scope}
\label{sec:intro:axioms}

Universal Coset Field Theory (UCFT) is formulated on the compact real form of
$E_6$, with order parameter $\Phi\in\mathbf{27}=J_3(\mathbb{O})$. The theory is
specified by two axioms and six additional postulates \textnormal{P1--P6}, stated
below. The axioms determine the algebraic arena and the kinematical framework;
the postulates supply the potential, vacuum geometry, spacetime structure,
matter content, and symmetry-breaking data required for contact with four-dimensional
physics.

\begin{axiom}[Order parameter]\label{ax:intro:I}
The order parameter is a field $\Phi$ valued in the matter representation
$\mathbf{27}=J_3(\mathbb{O})$ of $E_6$. The symmetric configuration $\Phi=0$ is
the unique $E_6$-fixed point; the potential of postulate~\textnormal{P1} renders
it a strict local maximum (\autoref{thm:potential:T5}), so that the rank-one
orbit is selected dynamically (\autoref{thm:potential:T6}).
\end{axiom}

\begin{axiom}[Quantum kinematics]\label{ax:intro:II}
Physical states $\psi\in\mathcal{H}$ are acted on by a one-parameter unitary
family $\{U(\tau)\}_{\tau\in\mathbb{R}}$ with $U(0)=\mathrm{Id}$,
$U(\tau+s)=U(\tau)U(s)$, and strong continuity.
\end{axiom}

In the main development, Axiom~II is realized through emergent time: a
Page--Wootters clock field $\varphi^{0}$ together with the Hamiltonian constraint
$\hat{H}_{\mathrm{phys}}\lvert\Psi\rangle\!\rangle=0$ of induced gravity, with
$U(\tau)$ recovered on physical states in the semiclassical regime via a timelike
Killing vector (\autoref{sec:spacetime:time}). The Standard Model and
gravitation are not consequences of the axioms alone; they follow upon imposing
postulates~\textnormal{P1--P6}.

\subsection{Postulates \textnormal{P1--P6}}
\label{sec:intro:postulates}

\begin{description}
  \item[\textnormal{P1} (Potential).] An $E_6$-invariant polynomial potential
    built from the cubic norm $N$ and the sharp map $X^{\#}$, of the form
    \begin{equation}
      V(X)\;=\;\alpha\,\langle X^{\#},X^{\#}\rangle
        \;+\;\beta\,\bigl[\,\mathrm{Tr}(X^{2})-c_2\,\bigr]^{2}
        \;+\;\gamma\,\mathrm{Tr}(X^{2})^{k},
      \qquad \alpha,\beta,\gamma>0,\;\; k\ge 2 \text{ even},
      \label{eq:intro:V}
    \end{equation}
    whose global minimum is the rank-one idempotent orbit. Coercivity is
    established on the compact real form (\autoref{sec:potential:main}).
  \item[\textnormal{P2} (Vacuum).] The vacuum manifold is
    \begin{equation}
      \mathcal{M}\;=\;E_6/\bigl(\mathrm{Spin}(10)\times U(1)\bigr)\;=\;\mathrm{EIII},
      \qquad \dim_{\mathbb{R}}\mathcal{M}=32,
      \label{eq:intro:vacuum}
    \end{equation}
    the rank-one primitive-idempotent orbit (the complex Cayley plane). This is
    proved from~\textnormal{P1} as \autoref{thm:potential:T4}.
  \item[\textnormal{P3} (Soldering).] A soldering map identifies four
    clock-field directions with a tangent frame carrying the Minkowski form
    inherited from $N$ via $H_2(\mathbb{O})\cong\mathbb{R}^{1,9}$. The
    dimension $d=4$ and signature $(1,3)$ are part of this postulate
    (\autoref{sec:spacetime:main}).
  \item[\textnormal{P4} (Matter).] One chiral $\mathbf{16}$ of
    $\mathrm{Spin}(10)$ per family, with charges as in
    \eqref{eq:intro:27branch}.
  \item[\textnormal{P5} (Generations).] Three families via an $SU(3)_F$
    horizontal symmetry on $J_3(\mathbb{O})\otimes\mathbb{C}^{3}$, corroborated
    by the heterotic index $h^{1}(\hat Q,V_{16})=3$.
  \item[\textnormal{P6} (Yukawa sector).] Standard $\mathrm{Spin}(10)$ Higgs and
    Yukawa multiplets ($\mathbf{45}_H/\mathbf{54}_H$, $\mathbf{126}_H$,
    $\mathbf{10}_H$) generating fermion masses and the type-I seesaw.
\end{description}

\subsection{Status of results}
\label{sec:intro:convention}

Results are classified as follows.
\begin{description}
  \item[Theorem.] Proved from explicitly stated hypotheses.
  \item[Postulate.] Additional structure beyond the axioms (P1--P6).
  \item[Conjecture / indicative.] Motivated by explicit but uncontrolled
    calculations, or low-energy correspondences not established as dualities.
\end{description}

\begin{table}[t]
  \centering
  \begin{tabular}{@{}lll@{}}
    \toprule
    Result & Content & Status \\
    \midrule
    Vacuum selection
      & P1$\Rightarrow$P2: global minimum $=$ rank-one EIII orbit
      & Theorem \\
    Origin instability
      & $\mathrm{Hess}\,V(0)\prec 0$, roll-down to $\mathcal{M}$
      & Theorem \\
    Minkowski block
      & $H_2(\mathbb{O})\cong\mathbb{R}^{1,9}$, $SL(2,\mathbb{O})\cong\mathrm{Spin}(1,9)$
      & Theorem \\
    Signature $(1,3)$, $d=4$
      & Soldering postulate P3
      & Postulate \\
    Emergent time
      & Page--Wootters clock; OS reconstruction
      & Theorem / conjecture \\
    Induced Planck mass
      & $M_{\mathrm{Pl}}^{2}=\tfrac{28}{6}\,f^{2}\mu^{2}\log(\Lambda^{2}/\mu^{2})/(16\pi^{2})>0$
      & Theorem \\
    Gauge running
      & Asymptotically free; $b_0(1)=26$, $b_0(3)=12$
      & Theorem \\
    Gravitational fixed point
      & $\kappa_\star=96\pi^{2}/5$
      & Conjecture \\
    Standard Model content
      & One family per $\mathbf{16}$; anomaly cancellation
      & Theorem \\
    Flavour sector
      & Three generations; Yukawa; $m_\nu\sim 0.05\,\mathrm{eV}$
      & Postulate \\
    Heterotic correspondence
      & $(c_L,c_R)=(16,10)$, $h^{1}(V_{16})=3$
      & Correspondence \\
    \bottomrule
  \end{tabular}
  \caption{Classification of the principal results.}
  \label{tab:intro:grades}
\end{table}

\subsection{Main theorem}
\label{sec:intro:claim}

\begin{theorem}[Conditional embedding]\label{thm:intro:embedding}
Assume Axioms~\autoref{ax:intro:I} and~\autoref{ax:intro:II} on the compact real
form of $E_6$, together with postulates~\textnormal{P1--P6}. Then:
\begin{enumerate}[label=\textup{(\roman*)}]
  \item \textnormal{(\autoref{thm:potential:T4}).}
    The global minimum of~\eqref{eq:intro:V} is the orbit
    $\mathcal{M}=E_6/(\mathrm{Spin}(10)\times U(1))=\mathrm{EIII}$,
    with $32$ Goldstone modes spanning
    \begin{equation}
      \mathfrak{m}\;=\;\mathbf{16}_{-3}\,\oplus\,\overline{\mathbf{16}}_{+3},
      \label{eq:intro:coset}
    \end{equation}
    and positive-definite coset metric $K|_{\mathfrak{m}}$.
  \item \textnormal{(P3, \autoref{sec:spacetime:main}).}
    Soldering yields a four-dimensional Lorentzian manifold of signature
    $(-,+,+,+)$ with $28$ physical massive scalars of mass $m^{2}=\mu^{2}$.
  \item \textnormal{(\autoref{sec:gravity:main}).}
    One-loop fluctuations induce
    \begin{equation}
      M_{\mathrm{Pl}}^{2}\;=\;\frac{28}{6}\,
      \frac{f^{2}\mu^{2}\,\log(\Lambda^{2}/\mu^{2})}{16\pi^{2}}\;>\;0.
      \label{eq:intro:MPl}
    \end{equation}
  \item \textnormal{(\autoref{sec:gravity:main}, \autoref{sec:frg:main}).}
    The composite $\mathrm{Spin}(10)\times U(1)$ gauge sector is asymptotically
    free with $\hat g^{2}_\star=0$ and
    \begin{equation}
      b_0=\tfrac{11}{3}C_A-\tfrac{4}{3}T\,n_F-\tfrac{1}{6}T\,n_S,
      \label{eq:intro:betag}
    \end{equation}
    where $b_0=26$ for one generation and $b_0=12$ for three.
  \item \textnormal{(P4--P6, \autoref{sec:sm:main}).}
    The model reduces to the Standard Model with three chiral families and
    neutrino masses via the type-I seesaw.
  \item \textnormal{(\autoref{sec:frg:main}).}
    The gravitational coupling possesses a conjectural ultraviolet fixed point
    $\kappa_\star=96\pi^{2}/5$.
\end{enumerate}
\end{theorem}

\begin{proof}
Parts~(i), (iii), (iv), and the anomaly structure of~(v) are established in the
sections cited. Part~(ii) combines the soldering postulate with the algebraic
theorems of \autoref{sec:spacetime:main}. The flavour content of~(v) is
postulated. Part~(vi) rests on a one-loop truncation.
\end{proof}

\begin{remark}[Scope]\label{rem:intro:caveats}
(i)~Lorentzian signature is not inherited from the coset metric $K|_{\mathfrak{m}}$,
which is positive-definite; it enters through~\textnormal{P3}.
(ii)~The vacuum is the rank-one EIII orbit~\eqref{eq:intro:vacuum}, not the
$26$-dimensional $E_6/F_4$ orbit of $\mathrm{diag}(1,1,1)$.
(iii)~The cosmological constant is a tunable parameter; the fine-tuning problem is
not addressed here.
\end{remark}

\subsection{Organization}
\label{sec:intro:roadmap}

\autoref{sec:master:chain} gives an overview of the logical dependencies.
\autoref{sec:albert:main} reviews the Albert algebra and $E_6$ structure.
\autoref{sec:potential:main} proves vacuum selection; \autoref{sec:coset:main}
develops the coset $\sigma$-model. \autoref{sec:spacetime:main} treats soldering
and emergent time. \autoref{sec:gravity:main} and \autoref{sec:frg:main} compute
induced gravitation and renormalization-group flow. \autoref{sec:sm:main}
carries out $\mathrm{Spin}(10)$ descent to the Standard Model;
\autoref{sec:strings:main} discusses string-theoretic correspondences. Appendices
collect the underlying octonionic and representation-theoretic results.

We use signature $(-,+,+,+)$ throughout; $[f]=\mathrm{mass}$, and
$\hat g^{2}$, $\xi$, $\hat y^{2}$, $\kappa$, $\Lambda_{\mathrm{cc}}$ are
dimensionless. The $32$ coset Goldstone coordinates are termed \emph{clock fields}.% --- end inlined: unit-intro.tex ---
% --- begin inlined: unit-axioms.tex ---
%======================================================================%
% UNIT: axioms  --  Axioms and Added Postulates
%======================================================================%
\section{Axioms and Added Postulates}
\label{sec:axioms:main}

This section fixes the foundation of Universal Coset Field Theory (UCFT).
We state the single dynamical \emph{axiom} that supplies the arena --- an
order parameter valued in the matter representation of $\ESix$, together
with a quantization prescription --- and we then enumerate, once and up
front, the added postulates $\Pone$--$\Psix$ that convert this arena into a
four-dimensional Lorentzian quantum field theory. UCFT is formulated as a
\emph{conditional embedding}: the axiom supplies the kinematical arena,
$\Pone$--$\Psix$ supply the remaining structure, and every nontrivial claim
below is classified as a \textbf{theorem} (proved from stated hypotheses),
\textbf{postulate} (assumed: $\Pone$--$\Psix$), or
\textbf{conjecture} (motivated, not controlled).

Conventions used throughout the manuscript are fixed here: the spacetime
metric signature is $(-,+,+,+)$; the decay constant $f$ carries mass
dimension $[f]=\text{mass}$; the breaking/mass scale is $\mu$, the UV cutoff
is $\Lambda$, and the breaking-VEV (GUT) scale tied to $f$ is $\MGUT$; the
couplings $\ghat^2$, $\xi$, $\yhat^2$, $\kappa$, $\Lambda_{\mathrm{cc}}$ are
dimensionless. The $32$ Goldstone directions of the symmetry-breaking coset
are called the \emph{clock fields}.

%%%%%%%%%%%%%%%%%%%%%%%%%%%%%%%%%%%%%%%%%%%%%%%%%%%%%%%%%%%%%%%%%%%%%%
\subsection{Axiom I: the order parameter and the \texorpdfstring{$\ESix$}{E6}/Albert arena}
\label{sec:axioms:axiomI}
%%%%%%%%%%%%%%%%%%%%%%%%%%%%%%%%%%%%%%%%%%%%%%%%%%%%%%%%%%%%%%%%%%%%%%

The arena of UCFT is the exceptional Jordan (Albert) algebra
$\jalg$, the $27$-dimensional formally real Jordan algebra of
$3\times3$ Hermitian octonionic matrices,
\[
  \jalg=\bigl\{X\in M_3(\Oct): X=X^{\dagger}\bigr\},
  \qquad X\circ Y=\tfrac12(XY+YX),
  \qquad \dim_{\mathbb R}\jalg=27 .
\]
Its reduced structure group --- the group preserving the cubic norm
$N(X)=\det_{\Oct}X$ up to a real character --- is the exceptional Lie group
$\ESix$, and the $27$ is its defining (matter) representation. We commit to
\emph{two real forms} (see \autoref{sec:axioms:realform}): the
\textbf{compact} $\ESix$ (equivalently $E_{6(-14)}$) for all dynamics, and
the split-signature $E_{6(-26)}$ used \emph{only} as the rigid classifier of
the real orbits.

\begin{axiom}[Order parameter]\label{ax:axioms:I}
The fundamental order parameter of UCFT is a field $\Phi$ that is a section
of an $\ESix$ vector bundle valued in the matter representation
$\Rep{27}\cong\jalg$. Locally, $\Phi:\,U\to\jalg$.
\end{axiom}

\paragraph{Quantization prescription (theorem-grade kinematics).}
The quantum theory is defined by the Euclidean functional integral over
configurations of $\Phi$ with the $\ESix$-invariant weight $e^{-S_E[\Phi]}$,
$S_E=\int\!\bigl(\tfrac12\,K_{\!AB}\,\partial\Phi^A\partial\Phi^B+V(\Phi)\bigr)$,
treated as the generating functional of connected correlators. Its physical
(Lorentzian, unitary, positive-energy) content is recovered not by any naive
``$e^0\mapsto i\,e^0$'' Wick trick but by Osterwalder--Schrader (OS)
reconstruction along the soldered timelike clock field $\clock$
(\autoref{sec:axioms:time}, and \autoref{sec:soldering:os}). The OS
reconstruction theorem is a \textbf{theorem} given its hypotheses; that the
UCFT $\sigma$-model \emph{satisfies} reflection positivity is
\textbf{indicative}, a property to be verified, not assumed.

\begin{remark}[Axiom~I and emergent time]\label{rem:axioms:honest}
Earlier formulations posited a second axiom asserting a global unitary
$U(t)=e^{-iHt}$ with an external time parameter. In the present foundation
that statement is \textbf{realized as a theorem}: time is emergent, recovered
as a Page--Wootters clock on physical states subject to the induced-gravity
Hamiltonian constraint (\autoref{thm:axioms:time}). UCFT therefore rests on the
\emph{single} dynamical Axiom~\ref{ax:axioms:I} together with the quantization
prescription; the axiom supplies the arena and $\Pone$--$\Psix$ supply the
added structure required for contact with four-dimensional physics.
\end{remark}

%%%%%%%%%%%%%%%%%%%%%%%%%%%%%%%%%%%%%%%%%%%%%%%%%%%%%%%%%%%%%%%%%%%%%%
\subsection{Real form, invariants, and the vacuum manifold}
\label{sec:axioms:realform}
%%%%%%%%%%%%%%%%%%%%%%%%%%%%%%%%%%%%%%%%%%%%%%%%%%%%%%%%%%%%%%%%%%%%%%

Two facts must not be conflated. The automorphism group of $\jalg$ (the group
preserving the Jordan product, hence \emph{both} $Q(X)=\Tr(X^2)$ and $N(X)$)
is the compact $F_4$, $\dim 52$. The reduced structure group preserving
\emph{only} $N$ (up to a character) is the larger $\ESix$, $\dim 78$; it does
\textbf{not} preserve $Q=\Tr(X^2)$. Consequently:

\begin{theorem}[Invariants and orbit classification]\label{thm:axioms:orbits}
Only the cubic norm $N$ is $\ESix$ structure-group invariant, up to a real
character $\lambda(g)$; the quadratic form $Q(X)=\Tr(X^2)$ is
$F_4$-invariant only. Real $\ESix$-orbits in $\jalg$ are classified by
$\mathrm{rank}\,X\in\{0,1,2,3\}$ together with $N(X)$ ---
\emph{not} by the pair $(Q,N)$. Moreover
\[
  \mathrm{rank}\,X\le1\iff X^{\#}=0,\qquad
  \mathrm{rank}\,X\le2\iff N(X)=0,\qquad
  \mathrm{rank}\,X=3\iff N(X)\neq0,
\]
where $X^{\#}=\nabla N(X)$ is the $\ESix$-covariant sharp (adjoint) map.
\end{theorem}

\begin{remark}[Real form committed to]\label{rem:axioms:compact}
For the dynamics we use the \textbf{compact} $\ESix$ (equivalently
$E_{6(-14)}$, whose maximal compact subgroup is exactly
$H=\SOTen\times\UOne$). On the compact form the $27$ is unitary, the trace
form $Q$ is the invariant metric, the norm character is identically $1$, and
the coset metric is \emph{positive-definite} --- giving a healthy kinetic
term and a unitary gauge sector. Coercivity and boundedness-below of $V$ are
statements about the compact form (non-compact orbits cannot be coercive).
The split form $E_{6(-26)}$ is invoked \emph{only} in
\autoref{thm:axioms:orbits} as the rigid classifier of the real $27$ by rank
and $N$.
\end{remark}

\begin{theorem}[Vacuum manifold $=$ EIII]\label{thm:axioms:eiii}
The symmetry-breaking vacuum manifold of UCFT is the rank-$1$
primitive-idempotent orbit
\begin{equation}\label{eq:axioms:eiii}
  \M \;=\; \frac{\ESix}{\SOTen\times\UOne} \;=\; \mathrm{EIII},
  \qquad \dim_{\mathbb R}\M = 78-46 = 32,
\end{equation}
the closed minimal $\ESix$-orbit in $\mathbb P(27_{\mathbb C})$ --- the
complex Cayley plane. This is the orbit of a rank-$1$ primitive idempotent,
\textbf{not} the orbit of $\diag(1,1,1)$ (whose stabiliser is $F_4$ and whose
orbit $\ESix/F_4$ has real dimension $26$).
\end{theorem}

The defining datum of the breaking is the branching of the adjoint and the
fundamental under $H=\SOTen\times\UOne$:
\begin{align}
  \Rep{78} &= \Rep{45}_{0}\oplus\Rep{1}_{0}
              \oplus\Rep{16}_{-3}\oplus\Repbar{16}_{+3}, \label{eq:axioms:adj}\\
  \Rep{27} &= \Rep{1}_{+4}\oplus\Rep{10}_{-2}\oplus\Rep{16}_{+1}.
              \label{eq:axioms:fund}
\end{align}
The coset tangent space is therefore
\begin{equation}\label{eq:axioms:tangent}
  \mathfrak m \;=\; \Rep{16}_{-3}\oplus\Repbar{16}_{+3}
  \qquad (32\ \text{real}),
\end{equation}
a real-irreducible module of complex type (commutant $=\mathbb C$). The
$\UOne$ charge forbids the symmetric $\Rep{16}\times\Rep{16}$ and
$\Repbar{16}\times\Repbar{16}$ invariants, leaving the unique charge-zero
cross-pairing; by Schur the invariant coset metric $K|_{\mathfrak m}$ is
unique and \emph{positive-definite}. Crucially, $\mathfrak m$ contains
\textbf{no} $\SOTen$ singlet: it is \emph{not} $2\times\Rep{10}+4$, nor
$\Rep{10}+\Repbar{10}+4$, nor anything with a singlet. The $32$ Goldstones
spanning $\mathfrak m$ are the clock fields.

%%%%%%%%%%%%%%%%%%%%%%%%%%%%%%%%%%%%%%%%%%%%%%%%%%%%%%%%%%%%%%%%%%%%%%
\subsection{The added postulates \texorpdfstring{$\Pone$--$\Psix$}{P1--P6}}
\label{sec:axioms:postulates}
%%%%%%%%%%%%%%%%%%%%%%%%%%%%%%%%%%%%%%%%%%%%%%%%%%%%%%%%%%%%%%%%%%%%%%

We now state, once and explicitly, the added structure. These are
\textbf{postulates}, not theorems; the manuscript cites them as
$\Pone$--$\Psix$ wherever they are used.

\begin{postulate}[Potential class]\label{post:axioms:P1}
($\Pone$) The dynamics for $\Phi$ are governed by an $\ESix$-invariant
polynomial potential built from the cubic norm $N$ and the sharp map
$X^{\#}$,
\begin{equation}\label{eq:axioms:V}
  V(X)\;=\;\alpha\,\langle X^{\#},X^{\#}\rangle
          \;+\;\beta\,\bigl[\Tr(X^2)-c_2\bigr]^2
          \;+\;\gamma\,\Tr(X^2)^{k},
  \qquad \alpha,\beta,\gamma>0,\ \ k\ge2\ \text{even},
\end{equation}
whose global minimum is the rank-$1$ idempotent orbit. The covariant term
$\langle X^{\#},X^{\#}\rangle\ge0$ vanishes exactly on $\mathrm{rank}\le1$;
the $\Tr(X^2)$ terms are $F_4$-invariant selectors fixing the nonzero scale
and representative. On the compact real form every term of
\eqref{eq:axioms:V} is a genuine $\ESix$ invariant.
\end{postulate}

\begin{postulate}[Vacuum selection]\label{post:axioms:P2}
($\Ptwo$) The vacuum is the minimal orbit
$\M=\ESix/(\SOTen\times\UOne)=\mathrm{EIII}$, $\dim_{\mathbb R}\M=32$. This is
\emph{proved} from $\Pone$ as the vacuum-selection theorem
(\autoref{thm:potential:T4}); see also \autoref{thm:axioms:eiii}.
\end{postulate}

\begin{postulate}[Soldering and signature]\label{post:axioms:P3}
($\Pthree$) A soldering map identifies four clock-field frame directions
$e^{a}{}_{\mu}$ ($a=0,1,2,3$) with a tangent $4$-plane of an emergent
$4$-manifold $\mathcal M^4$, with vierbein
$e^{a}{}_{\mu}(x)=\langle\tau^a,\theta_{\mathfrak m}(\partial_\mu)\rangle$
the pull-back of the Maurer--Cartan form, and induced metric
$g_{\mu\nu}=e^{a}{}_{\mu}e^{b}{}_{\nu}\,\eta_{ab}$,
$\eta=\diag(-1,+1,+1,+1)$. The Minkowski form $\eta$ is the restriction of
the cubic norm $N$ to the sub-block $H_2(\Oct)\cong\mathbb R^{1,9}$, reduced
to $H_2(\mathbb C)\cong\mathbb R^{1,3}$ by a vierbein condensate. The
spacetime dimension $d=4$ and signature $(1,3)=(-,+,+,+)$ enter \emph{here},
as added structure --- not as Killing-form, coset, or rank-lemma
consequences. See \autoref{sec:soldering:main}.
\end{postulate}

\begin{postulate}[Matter content and chirality]\label{post:axioms:P4}
($\Pfour$) Each family contributes one chiral $\Rep{16}$ of $\SOTen$, with
the $\UOne$ charges of \eqref{eq:axioms:fund} (the family lives in
$\Rep{16}_{+1}\subset\Rep{27}$); there are no mirror partners.
\end{postulate}

\begin{postulate}[Family symmetry]\label{post:axioms:P5}
($\Pfive$) There are three generations, realized (primary route) by an
$SU(3)_F$ horizontal symmetry acting on $\jalg\otimes\mathbb C^3$,
yielding three chiral $\Rep{16}$s with no mirrors. This is corroborated
(a correspondence, not a second mechanism) by the heterotic count
$h^1(\widehat Q,V_{16})=3$. See \autoref{sec:sm:generations}.
\end{postulate}

\begin{postulate}[GUT-Higgs / Yukawa sector]\label{post:axioms:P6}
($\Psix$) Fermion masses arise from a real $\SOTen$ Higgs/Yukawa sector
($\Rep{45}_H$ or $\Rep{54}_H$, and $\Rep{126}_H$, with electroweak
$\Rep{10}_H$), generating Dirac masses and the type-I seesaw. The naive
coupling $-y\,\bar\psi\,\phi^a T_a\,\psi$ with $\phi\in\{\Rep{16},\Repbar{16}\}$
is forbidden: $\Rep{16}\times\Rep{16}=\Rep{10}+\Rep{120}+\Rep{126}$ and
$\Rep{16}\times\Repbar{16}=\Rep{1}+\Rep{45}+\Rep{210}$ contain no $H$-singlet
from such a scalar, and $H$ is unbroken. See \autoref{sec:sm:yukawa}.
\end{postulate}

\begin{remark}[Epistemic status of the postulates]\label{rem:axioms:tagging}
Of these, $\Ptwo$ is \emph{proved} from $\Pone$ (\autoref{thm:potential:T4});
$\Pthree$--$\Psix$ are \emph{added structure}. The gravitational
fixed point of \autoref{sec:gravity:as} is \emph{conjectural}, and the
Standard Model is \emph{embedded} under $\Pfour$--$\Psix$.
\end{remark}

%%%%%%%%%%%%%%%%%%%%%%%%%%%%%%%%%%%%%%%%%%%%%%%%%%%%%%%%%%%%%%%%%%%%%%
\subsection{From Axiom I to a broken vacuum}
\label{sec:axioms:broken}
%%%%%%%%%%%%%%%%%%%%%%%%%%%%%%%%%%%%%%%%%%%%%%%%%%%%%%%%%%%%%%%%%%%%%%

The $27$ is a non-trivial irreducible $\ESix$ representation, so the
\emph{only} vector fixed by the whole structure group is the origin
$\Phi=0$. Maximal symmetry is therefore not a stable ground state: the
potential $\Pone$ makes $\Phi=0$ a strict local maximum.

\begin{theorem}[The symmetric origin is a strict maximum]\label{thm:axioms:origin}
For $V$ as in \eqref{eq:axioms:V},
\[
  \Hess V(0)=-4\beta c_2\,\mathbb 1\prec0,
\]
i.e.\ all $27$ eigenvalues equal $-4\beta c_2<0$, so $\Phi=0$ is a strict
local maximum. Since $V(0)=\beta c_2^2$ strictly exceeds the value
$W(q_\star)$ attained on the rank-$1$ orbit, the broken vacuum
$\M=\mathrm{EIII}$ beats the origin both locally and globally.
\end{theorem}

\begin{proof}[Proof sketch]
The covariant term contributes $\langle X^{\#},X^{\#}\rangle=O(\|X\|^4)$, so
its Hessian vanishes at the origin; only the selector $\beta[\Tr(X^2)-c_2]^2$
contributes at quadratic order, giving $\Hess V(0)=-4\beta c_2\,\mathbb 1$.
The global statement is the vacuum-selection theorem
\autoref{thm:potential:T4}: $V=\alpha\langle X^{\#},X^{\#}\rangle+W(\Tr(X^2))$
with $W$ strictly convex and $W(q_\star)<W(0)=\beta c_2^2$, so every global
minimizer has $X^{\#}=0$ (rank $\le1$) and $\Tr(X^2)=q_\star>0$ (rank
exactly $1$).
\end{proof}

\begin{remark}[Forward reference to roll-down]\label{rem:axioms:rolldown}
Because $\Phi=0$ is linearly unstable (growth rate $4\beta c_2$) and $V$ is a
strict Lyapunov function that is coercive on the compact real form, generic
initial data flow under the equations of motion to the broken vacuum $\M$.
The roll-down dynamics, the tree-level vanishing of $\Hess V|_{\mathfrak m}$
(the $32$ exact Goldstone/clock fields), and the single common radiative mass
fixed by Schur on the irreducible $\mathfrak m$, are established in
\autoref{thm:potential:T5} and \autoref{thm:potential:T6}, with the spectrum
analysed in \autoref{sec:gravity:spectrum}.
\end{remark}

%%%%%%%%%%%%%%%%%%%%%%%%%%%%%%%%%%%%%%%%%%%%%%%%%%%%%%%%%%%%%%%%%%%%%%
\subsection{Reconciliation of axiomatic and emergent time}
\label{sec:axioms:time}
%%%%%%%%%%%%%%%%%%%%%%%%%%%%%%%%%%%%%%%%%%%%%%%%%%%%%%%%%%%%%%%%%%%%%%

The legacy ``Axiom~II'' posited a global unitary $U(t)=e^{-iHt}$ acting along
an external time. UCFT has no external time: time is one of the soldered
clock fields ($\Pthree$). We reconcile the two pictures by realizing the
unitary-evolution statement as a theorem on physical states.

\begin{theorem}[Emergent time as a Page--Wootters clock]\label{thm:axioms:time}
Induced gravity supplies a diffeomorphism/Hamiltonian constraint
$\widehat H_{\mathrm{phys}}\,\dket{\Psi}=0$ on physical states. Designate the
soldered timelike clock field $\clock$ a Page--Wootters clock. Conditioning,
\[
  \ket{\psi(\tau)}=\dbra{\clock=\tau}\Psi\rangle\!\rangle,
  \qquad
  i\,\partial_\tau\ket{\psi(\tau)}=\widehat H_{\mathrm{sys}}\ket{\psi(\tau)},
\]
so the one-parameter unitary group $U(\tau)=e^{-i\widehat H_{\mathrm{sys}}\tau}$
is recovered on physical states in the semiclassical regime, with the
identification $t\leftrightarrow x^0$ via
$H=\int T_{00}$ along the timelike Killing direction. Global hyperbolicity is
an assumption.
\end{theorem}

\begin{remark}[Euclidean to Lorentzian]\label{rem:axioms:os}
The passage to Lorentzian, unitary, positive-energy dynamics is via
Osterwalder--Schrader reconstruction along $\clock$ (reflection positivity),
\emph{not} by the ill-defined ``multiply $e^0$ by $i$'' Wick prescription.
The metric $\eta^{(1,3)}$ is already the soldered ($\Pthree$) Lorentzian form
before any continuation; OS reconstruction then supplies unitary,
positive-energy evolution. As noted under the quantization prescription, OS
reconstruction is a \textbf{theorem} given its hypotheses, while reflection
positivity of the UCFT $\sigma$-model is \textbf{indicative}. See
\autoref{sec:soldering:os}.
\end{remark}
% --- end inlined: unit-axioms.tex ---
%----------------------------------------------------------------------%
% Part: Albert Algebra, Vacuum Selection, and Coset Geometry
%----------------------------------------------------------------------%
\part{Albert Algebra, Vacuum Selection, and Coset Geometry}
% --- begin inlined: unit-albert.tex ---
%%%%%%%%%%%%%%%%%%%%%%%%%%%%%%%%%%%%%%%%%%%%%%%%%%%%%%%%%%%%%%%%%%%%%%
\section{The Albert Algebra and Exceptional Structure}
\label{sec:albert:main}
%%%%%%%%%%%%%%%%%%%%%%%%%%%%%%%%%%%%%%%%%%%%%%%%%%%%%%%%%%%%%%%%%%%%%%

This section assembles the algebraic foundation on which the whole
construction rests: the exceptional Jordan (Albert) algebra
$J_{3}(\mathbb{O})$, its trace and cubic invariants, the sharp map, the
notion of rank, and the four real forms of the exceptional groups that
act on it. We isolate, with care, \emph{which} invariant is preserved by
\emph{which} group, since the entire vacuum-selection argument
(\autoref{sec:potential:main}) and the Lorentzian soldering
(\autoref{sec:soldering:main}) depend on getting this exactly right. The
heavy octonionic identities (associator gymnastics, the explicit cubic
norm in components, the Freudenthal cross product) are deferred to
\autoref{app:albert:main}; here we state the structural facts and prove
the ones that are short. Throughout, the metric signature is
$(-,+,+,+)$, $[f]=\mathrm{mass}$, and the dimensionless couplings are
$\hat g^{2},\xi,\hat y^{2},\kappa$.

%%%%%%%%%%%%%%%%%%%%%%%%%%%%%%%%%%%%%%%%%%%%%%%%%%%%%%%%%%%%%%%%%%%%%%
\subsection{The exceptional Jordan algebra \texorpdfstring{$J_{3}(\mathbb{O})$}{J3(O)}}
\label{sec:albert:def}
%%%%%%%%%%%%%%%%%%%%%%%%%%%%%%%%%%%%%%%%%%%%%%%%%%%%%%%%%%%%%%%%%%%%%%

Let $\mathbb{O}$ denote the (real, division) octonions, with conjugation
$x\mapsto \bar x$, norm $n(x)=x\bar x\in\mathbb{R}_{\ge 0}$, and real part
$\operatorname{Re} x=\tfrac12(x+\bar x)$. Octonionic multiplication is
nonassociative; this nonassociativity is exactly what forces the
\emph{exceptional} structure below and obstructs any matrix-algebra
shortcut.

\begin{definition}[Albert algebra]\label{def:albert:J3O}
The \emph{Albert algebra} is the space of $3\times 3$ octonionic
Hermitian matrices
\begin{equation}\label{eq:albert:J3Odef}
  J_{3}(\mathbb{O})
  =\Bigl\{\,X=X^{\dagger}\in M_{3}(\mathbb{O})\,\Bigr\}
  =\left\{
  \begin{pmatrix}
   \xi_{1} & z_{3} & \bar z_{2}\\
   \bar z_{3} & \xi_{2} & z_{1}\\
   z_{2} & \bar z_{1} & \xi_{3}
  \end{pmatrix}
  : \xi_{i}\in\mathbb{R},\; z_{i}\in\mathbb{O}\right\},
\end{equation}
equipped with the commutative \emph{Jordan product}
\begin{equation}\label{eq:albert:jordanprod}
  X\circ Y=\tfrac12\,(XY+YX).
\end{equation}
\end{definition}

The three real diagonal entries and three octonionic off-diagonal
entries give
\begin{equation}\label{eq:albert:dim}
  \dim_{\mathbb{R}} J_{3}(\mathbb{O})
  = 3\cdot 1 + 3\cdot 8 = 27 .
\end{equation}

\begin{theorem}[Formally real Jordan algebra]\label{thm:albert:formallyreal}
$\bigl(J_{3}(\mathbb{O}),\circ\bigr)$ is a $27$-dimensional, commutative,
formally real (Euclidean) Jordan algebra: the Jordan identity
$X\circ\bigl(Y\circ X^{\circ 2}\bigr)=\bigl(X\circ Y\bigr)\circ X^{\circ 2}$
holds, and $\sum_{i} X_{i}\circ X_{i}=0$ implies every $X_{i}=0$. It is
the unique exceptional simple formally real Jordan algebra (it is not
special, i.e.\ not a Jordan subalgebra of any associative algebra under
the symmetrised product).
\end{theorem}

\begin{proof}[Proof sketch]
Commutativity is manifest from \eqref{eq:albert:jordanprod}. The Jordan
identity is a power-associativity statement that, by linearisation,
reduces to relations involving at most two octonions at a time; since any
two octonions generate an associative (indeed commutative) subalgebra
$\cong\mathbb{C}$ or $\mathbb{H}$ (Artin's theorem), the identity holds
despite global nonassociativity. Formal reality follows from positivity
of the trace form \eqref{eq:albert:traceform} below. Exceptionality
(non-speciality) is the Albert--Jacobson theorem; see
\autoref{app:albert:main} and \cite{Albert:1934,Jacobson:1968,McCrimmon:2004}.
\end{proof}

%%%%%%%%%%%%%%%%%%%%%%%%%%%%%%%%%%%%%%%%%%%%%%%%%%%%%%%%%%%%%%%%%%%%%%
\subsection{Trace form, cubic norm, and the sharp map}
\label{sec:albert:invariants}
%%%%%%%%%%%%%%%%%%%%%%%%%%%%%%%%%%%%%%%%%%%%%%%%%%%%%%%%%%%%%%%%%%%%%%

Three invariant objects organise the algebra: a quadratic trace form, a
cubic norm, and a quadratic-in-$X$ \emph{sharp} (adjoint) map. Their
distinct invariance groups are the crux of the whole construction.

\paragraph{Linear and quadratic traces.}
Define the linear trace and the symmetric \emph{trace form}
\begin{align}
  \operatorname{tr}(X) &= \xi_{1}+\xi_{2}+\xi_{3},
  \label{eq:albert:lintrace}\\[2pt]
  \langle X,Y\rangle &= \operatorname{Tr}(X\circ Y),
  \qquad
  Q(X)\equiv \operatorname{Tr}(X^{2})=\langle X,X\rangle .
  \label{eq:albert:traceform}
\end{align}
In components,
\begin{equation}\label{eq:albert:Qcomp}
  Q(X)=\operatorname{Tr}(X^{2})
  =\sum_{i}\xi_{i}^{2}+2\sum_{i} n(z_{i}) \;\ge 0 ,
\end{equation}
so $Q$ is positive-definite: $J_{3}(\mathbb{O})$ is Euclidean, and $Q$ is
the natural $\mathrm{O}(27)$-type inner product. This positivity is what
makes the dynamical coset metric definite later
(\autoref{sec:coset:main}).

\paragraph{Cubic norm.}
The \emph{cubic norm} (the octonionic ``determinant'') is
\begin{equation}\label{eq:albert:Ndef}
  N(X)=\det\nolimits_{\mathbb{O}}(X)
  =\xi_{1}\xi_{2}\xi_{3}
   -\sum_{i}\xi_{i}\,n(z_{i})
   +2\,\operatorname{Re}\!\bigl(z_{1}z_{2}z_{3}\bigr).
\end{equation}
The bracketing of the cubic term $z_{1}z_{2}z_{3}$ is unambiguous under
$\operatorname{Re}(\cdot)$ despite nonassociativity (a standard
octonionic identity; \autoref{app:albert:main}). Unlike $Q$, the cubic
norm is \emph{odd} under $X\mapsto -X$ and is genuinely indefinite; its
restriction to the $H_{2}(\mathbb{O})$ sub-block supplies the Minkowski
form used in soldering (\autoref{sec:soldering:main}).

\paragraph{Sharp map.}
The \emph{sharp} (adjoint) map $X\mapsto X^{\#}$ is the gradient of $N$
with respect to the trace form,
\begin{equation}\label{eq:albert:sharpdef}
  \langle X^{\#}, Y\rangle = \operatorname{D}N(X)[Y]
  \quad\Longleftrightarrow\quad
  X^{\#}=\nabla N(X),
\end{equation}
equivalently the quadratic map characterised by
\begin{equation}\label{eq:albert:sharpid}
  X\circ X^{\#}=N(X)\,\mathbf 1,
  \qquad
  (X^{\#})^{\#}=N(X)\,X .
\end{equation}
In a \emph{Jordan frame} $X=\operatorname{diag}(\lambda_{1},\lambda_{2},\lambda_{3})$,
\begin{equation}\label{eq:albert:sharpframe}
  X^{\#}=\operatorname{diag}\!\bigl(\lambda_{2}\lambda_{3},\,
  \lambda_{3}\lambda_{1},\,\lambda_{1}\lambda_{2}\bigr),
  \qquad
  \langle X^{\#},X^{\#}\rangle
  =(\lambda_{2}\lambda_{3})^{2}+(\lambda_{3}\lambda_{1})^{2}+(\lambda_{1}\lambda_{2})^{2}\ge 0 .
\end{equation}
The combination $\langle X^{\#},X^{\#}\rangle$ is a sum of squares,
hence non-negative; this is the quantity whose vanishing characterises
the vacuum orbit.

%%%%%%%%%%%%%%%%%%%%%%%%%%%%%%%%%%%%%%%%%%%%%%%%%%%%%%%%%%%%%%%%%%%%%%
\subsection{Rank, idempotents, and the rank-one criterion}
\label{sec:albert:rank}
%%%%%%%%%%%%%%%%%%%%%%%%%%%%%%%%%%%%%%%%%%%%%%%%%%%%%%%%%%%%%%%%%%%%%%

\begin{definition}[Rank]\label{def:albert:rank}
The \emph{rank} of $X\in J_{3}(\mathbb{O})$ is the degree of its minimal
polynomial; equivalently, in a Jordan frame it is the number of nonzero
eigenvalues $\lambda_{i}$. Thus $\operatorname{rank}(X)\in\{0,1,2,3\}$,
with $\operatorname{rank}(0)=0$. A \emph{primitive idempotent} is a
rank-one $E$ with $E\circ E=E$ (equivalently $E^{2}=E$, $\operatorname{tr}E=1$).
\end{definition}

\begin{theorem}[Rank-one criterion]\label{thm:albert:rankone}
For all $X\in J_{3}(\mathbb{O})$,
\begin{equation}\label{eq:albert:rankone}
  \operatorname{rank}(X)\le 1
  \quad\Longleftrightarrow\quad
  X^{\#}=0 .
\end{equation}
Moreover $\operatorname{rank}(X)\le 2\Leftrightarrow N(X)=0$ (with
$X^{\#}\neq 0$), and $\operatorname{rank}(X)=3\Leftrightarrow N(X)\neq 0$.
\end{theorem}

\begin{proof}
Every $X$ is conjugate by $F_{4}=\operatorname{Aut}(J_{3}(\mathbb{O}))$
(\autoref{thm:albert:fourforms}) to a Jordan frame
$\operatorname{diag}(\lambda_{1},\lambda_{2},\lambda_{3})$
(spectral theorem for Euclidean Jordan algebras). By
\eqref{eq:albert:sharpframe}, $X^{\#}=0$ iff all three pairwise products
$\lambda_{i}\lambda_{j}$ vanish, i.e.\ at most one $\lambda_{i}$ is
nonzero, i.e.\ $\operatorname{rank}(X)\le1$. The norm statements follow
from $N(X)=\lambda_{1}\lambda_{2}\lambda_{3}$ in a frame: $N=0$ iff some
$\lambda_{i}=0$ (rank $\le 2$), and $N\neq0$ iff all nonzero (rank $3$).
Since $F_{4}$ preserves both $N$ and $X^{\#}$, the criteria are
frame-independent.
\end{proof}

\begin{remark}[Scope]\label{rmk:albert:scope}
\autoref{thm:albert:rankone} is a \emph{theorem} of Jordan algebra, used
verbatim in the vacuum-selection argument: the $E_{6}$-covariant term
$\langle X^{\#},X^{\#}\rangle$ in the potential $V$ vanishes
\emph{exactly} on the rank-$\le1$ cone (\autoref{sec:potential:main}).
The selection of the rank-one orbit as the global minimum is proved
there; here we only supply the algebraic criterion.
\end{remark}

%%%%%%%%%%%%%%%%%%%%%%%%%%%%%%%%%%%%%%%%%%%%%%%%%%%%%%%%%%%%%%%%%%%%%%
\subsection{The four real forms and what each preserves}
\label{sec:albert:groups}
%%%%%%%%%%%%%%%%%%%%%%%%%%%%%%%%%%%%%%%%%%%%%%%%%%%%%%%%%%%%%%%%%%%%%%

The single most consequential fact of this section is that
\emph{different groups preserve different invariants}. Conflating them
(in particular, asserting that the structure group preserves the trace
form $Q$, or that $(Q,N)$ separates orbits) is the root error against
which the entire construction is hardened.

\begin{theorem}[Four real forms]\label{thm:albert:fourforms}
The following groups act on $J_{3}(\mathbb{O})$ and the real $27$:
\begin{enumerate}
\item\label{it:albert:F4}
  $F_{4}=\operatorname{Aut}\bigl(J_{3}(\mathbb{O}),\circ\bigr)$ is the
  \textbf{compact} exceptional group, $\dim F_{4}=52$. It preserves the
  full Jordan product, hence \emph{both} the trace form
  $Q=\operatorname{Tr}(X^{2})$ \emph{and} the cubic norm $N$, and acts as
  $F_{4}\subset\mathrm{SO}(27)$ (it is the stabiliser of the identity
  $\mathbf 1$).
\item\label{it:albert:E6m26}
  The \textbf{reduced structure group}
  $E_{6(-26)}$ is the \textbf{non-compact} real form,
  $\dim E_{6(-26)}=78$. It preserves the cubic norm $N$ \emph{up to a
  real positive character} $\lambda(g)$,
  \begin{equation}\label{eq:albert:Ncharacter}
    N(g\cdot X)=\lambda(g)\,N(X),\qquad \lambda:E_{6(-26)}\to\mathbb{R}_{>0},
  \end{equation}
  and does \emph{not} preserve $Q$. (If it did, it would lie in
  $\mathrm{SO}(27)$ and be compact, a contradiction.)
\item\label{it:albert:E6c}
  The \textbf{compact} real form, here denoted compact $E_{6}$,
  $\dim E_{6}=78$, acts unitarily on the complexified $27_{\mathbb{C}}$;
  on it the norm character is identically $1$ (a continuous homomorphism
  from a compact connected group to $\mathbb{R}_{>0}$ is trivial), and
  $Q$ is the invariant Hermitian form.
\item\label{it:albert:E6m14}
  The real form $E_{6(-14)}$, $\dim=78$, has \textbf{maximal compact
  subgroup} $\operatorname{Spin}(10)\times\mathrm{U}(1)$.
\end{enumerate}
\end{theorem}

\begin{proof}[Proof sketch]
\ref{it:albert:F4} An automorphism of the Jordan product preserves
powers, hence the trace form (via $Q(X)=\operatorname{Tr}X^{2}$) and the
norm (via Cayley--Hamilton in the algebra); compactness and $\dim 52$ are
classical \cite{Springer:2000}. \ref{it:albert:E6m26} The reduced
structure group is by definition the norm-similarity group;
\eqref{eq:albert:Ncharacter} is its defining property, and non-compactness
follows because a non-trivial similarity character exists. \ref{it:albert:E6c}
Triviality of the character on a compact group is the homomorphism
argument given. \ref{it:albert:E6m14} is the standard real-form
classification; the maximal compact of $E_{6(-14)}$ is
$\operatorname{Spin}(10)\times\mathrm{U}(1)$
\cite{Springer:2000,Slansky:1981,Helgason:1978}.
\end{proof}

\begin{postulate}[Real form for dynamics]\label{post:albert:realform}
For the \emph{dynamics} of UCFT we commit to the \textbf{compact} real
form (equivalently $E_{6(-14)}$, whose maximal compact is exactly
$H=\operatorname{Spin}(10)\times\mathrm{U}(1)$). Then $H$ is compact, the
coset metric on $E_{6}/H$ is positive-definite, the kinetic term is
healthy, and the gauge sector is unitary. The non-compact form
$E_{6(-26)}$ is used \textbf{only} as the rigid orbit-classifier of the
real $27$. This choice is part of the postulate set (it is the form on
which the potential $V$ is coercive; non-compact orbits cannot be
coercive, since the norm character \eqref{eq:albert:Ncharacter} is then
unbounded). Cf.\ \hyperref[post:potential:P1]{P1} and the vacuum theorem
in \autoref{sec:potential:main}.
\end{postulate}

\begin{theorem}[Orbit classification by rank and norm]\label{thm:albert:orbits}
The orbits of $E_{6(-26)}$ on the real $27=J_{3}(\mathbb{O})$ are
classified by the pair $\bigl(\operatorname{rank}(X),\,N(X)\bigr)$, and
\emph{not} by $\bigl(Q(X),N(X)\bigr)$. Explicitly, the nonzero orbit
strata are
\begin{equation}\label{eq:albert:orbitstrata}
  \underbrace{\operatorname{rank}=1}_{X^{\#}=0}\;\subset\;
  \underbrace{\operatorname{rank}=2}_{X^{\#}\neq 0,\,N=0}\;\subset\;
  \underbrace{\operatorname{rank}=3}_{N\neq 0}\,,
\end{equation}
the rank-$3$ stratum further refined by the sign / value of $N$. The
trace form $Q$ is \textbf{not} an $E_{6}$-invariant (only an
$F_{4}$-invariant), so any ``$(Q,N)$ separates orbits'' statement is
false.
\end{theorem}

\begin{proof}
$E_{6(-26)}$ preserves $N$ up to the character
\eqref{eq:albert:Ncharacter} and acts transitively on each rank stratum
of fixed $N$-class (Springer; \autoref{app:albert:main}). The rank is an
algebraic invariant (\autoref{def:albert:rank}), preserved because the
sharp map is $E_{6}$-covariant: $(g\cdot X)^{\#}=\lambda(g)\,
(g^{\#})\!\cdot X^{\#}$ in the appropriate contragredient action, so
$X^{\#}=0$ is a group-invariant condition (and likewise $N=0$). That $Q$
is not invariant is \autoref{thm:albert:fourforms}\ref{it:albert:E6m26}:
a structure-group element with $\lambda(g)\neq1$ generically rescales
$\operatorname{Tr}(X^{2})$.
\end{proof}

\begin{corollary}[Closed minimal orbit]\label{cor:albert:minorbit}
The unique closed nonzero $E_{6}$-orbit in $\mathbb{P}(27_{\mathbb{C}})$
is the rank-one (primitive-idempotent) orbit, the \emph{complex Cayley
plane}
\begin{equation}\label{eq:albert:EIII}
  \mathcal{M}
  = E_{6}/\bigl(\operatorname{Spin}(10)\times\mathrm{U}(1)\bigr)
  = \mathrm{EIII},
  \qquad
  \dim_{\mathbb{R}}\mathcal{M}=32 .
\end{equation}
This is the orbit of a rank-one primitive idempotent, characterised by
$X^{\#}=0$ (\autoref{thm:albert:rankone}); it is \textbf{not} the orbit
of $\operatorname{diag}(1,1,1)$ (which is rank $3$, has
$\langle X^{\#},X^{\#}\rangle=3>0$, stabiliser $F_{4}$, orbit $E_{6}/F_{4}$
of real dimension $26$).
\end{corollary}

\begin{proof}
The closed orbit of a reductive group in a projective space is the orbit
of a highest-weight line; for the $27$ of $E_{6}$ this is the rank-one
line, with stabiliser the parabolic whose reductive part fixes the
$H$-decomposition \eqref{eq:albert:adjbranch}. Over the compact form the
isotropy is $H=\operatorname{Spin}(10)\times\mathrm{U}(1)$, and
$\dim_{\mathbb{R}}E_{6}-\dim_{\mathbb{R}}H=78-46=32$. The real rank-one
idempotents alone form $F_{4}/\operatorname{Spin}(9)\cong
\mathbb{OP}^{2}$ of dimension $16$; the full complex line under $E_{6}$
sweeps the larger $32$-dimensional $\mathrm{EIII}$
\cite{Helgason:1978,Springer:2000}.
\end{proof}

\begin{remark}[Consistency constraints]\label{rmk:albert:forbidden}
Three statements are \emph{false} and are never used:
(i) that the structure group preserves $Q=\operatorname{Tr}(X^{2})$
(only $F_{4}$ does);
(ii) that $(Q,N)$ separates $E_{6}$-orbits (orbits are classified by rank
and $N$, \autoref{thm:albert:orbits});
(iii) that the $32$-dimensional vacuum is the orbit of
$\operatorname{diag}(1,1,1)$ with stabiliser $F_{4}$ (that orbit is
$E_{6}/F_{4}$, dimension $26$; the vacuum is the rank-one EIII orbit,
\autoref{cor:albert:minorbit}).
\end{remark}

%%%%%%%%%%%%%%%%%%%%%%%%%%%%%%%%%%%%%%%%%%%%%%%%%%%%%%%%%%%%%%%%%%%%%%
\subsection{Branchings under \texorpdfstring{$H=\operatorname{Spin}(10)\times\mathrm{U}(1)$}{H = Spin(10) x U(1)}}
\label{sec:albert:branchings}
%%%%%%%%%%%%%%%%%%%%%%%%%%%%%%%%%%%%%%%%%%%%%%%%%%%%%%%%%%%%%%%%%%%%%%

The isotropy group $H=\operatorname{Spin}(10)\times\mathrm{U}(1)$ of the
vacuum orbit \eqref{eq:albert:EIII} is exactly the gauge group of the
construction. The relevant representations branch as follows; these are
the canonical decompositions used everywhere downstream.

\begin{theorem}[Defining branchings]\label{thm:albert:branchings}
Under $H=\operatorname{Spin}(10)\times\mathrm{U}(1)$ (subscripts are the
$\mathrm{U}(1)$ charges),
\begin{align}
  \mathbf{78}
   &= \mathbf{45}_{0}\oplus\mathbf{1}_{0}
      \oplus\mathbf{16}_{-3}\oplus\overline{\mathbf{16}}_{+3}
      \qquad\text{(adjoint),}
   \label{eq:albert:adjbranch}\\[2pt]
  \mathbf{27}
   &= \mathbf{1}_{+4}\oplus\mathbf{10}_{-2}\oplus\mathbf{16}_{+1}
      \qquad\text{(fundamental / matter).}
   \label{eq:albert:fundbranch}
\end{align}
Consequently the coset tangent space is
\begin{equation}\label{eq:albert:cosettangent}
  \mathfrak{m}
   = \mathbf{16}_{-3}\oplus\overline{\mathbf{16}}_{+3}
   \qquad(\,32\text{ real dimensions}\,),
\end{equation}
which contains \textbf{no} $\operatorname{Spin}(10)$ singlet.
\end{theorem}

\begin{proof}
The decompositions \eqref{eq:albert:adjbranch}--\eqref{eq:albert:fundbranch}
are the standard $E_{6}\downarrow\operatorname{Spin}(10)\times\mathrm{U}(1)$
branchings \cite{Slansky:1981}; the $\mathrm{U}(1)$ charges are fixed by
requiring the cubic $E_{6}$-invariant $\mathbf{27}^{3}$ and the adjoint
commutators to be charge-neutral. The coset directions are the
non-$H$ generators in \eqref{eq:albert:adjbranch}, namely
$\mathbf{16}_{-3}\oplus\overline{\mathbf{16}}_{+3}$, of real dimension
$16+16=32=\dim_{\mathbb{R}}\mathcal{M}$ (\autoref{cor:albert:minorbit}).
Dimension and charge counting exclude any singlet in $\mathfrak{m}$.
\end{proof}

\begin{theorem}[Irreducibility and the definite coset metric]\label{thm:albert:irrep}
The coset tangent module $\mathfrak{m}=\mathbf{16}_{-3}\oplus
\overline{\mathbf{16}}_{+3}$ is real-irreducible \emph{of complex type}
(its commutant is $\mathbb{C}$). The $H$-invariant metric on
$\mathfrak{m}$ is therefore \textbf{unique up to scale} and, on the
compact real form (\autoref{post:albert:realform}),
\textbf{positive-definite}.
\end{theorem}

\begin{proof}
The $\mathrm{U}(1)$ charge $\pm3$ forbids the symmetric self-pairings
$\mathbf{16}_{-3}\!\otimes\!\mathbf{16}_{-3}$ and
$\overline{\mathbf{16}}_{+3}\!\otimes\!\overline{\mathbf{16}}_{+3}$
(charges $\mp6\neq0$), leaving only the charge-$0$ cross-pairing
$\mathbf{16}_{-3}\!\otimes\!\overline{\mathbf{16}}_{+3}$, which is the
restriction of the Killing form. Hence the invariant bilinear is
one-dimensional (Schur, complex type), so unique up to scale. On the
compact form the Killing form is definite; restricted to $\mathfrak{m}$
it is positive-definite, giving the Euclidean target-space metric
$K|_{\mathfrak{m}}$ used in \autoref{sec:coset:main}.
\end{proof}

\begin{remark}[No Lorentzian content in the coset metric]\label{rmk:albert:nolorentz}
The coset metric $K|_{\mathfrak{m}}$ is positive-definite, signature
$(32,0)$: it carries \emph{zero} Lorentzian content. In particular there
is no $(4,28)$ Killing/coset signature, and the physical $(1,3)$ Minkowski
signature does \emph{not} fall out of the coset geometry. Lorentzian
signature enters only through the soldering postulate
\hyperref[post:soldering:P3]{P3}, which uses the indefinite cubic norm
$N$ restricted to $H_{2}(\mathbb{O})\cong\mathbb{R}^{1,9}$ on a
\emph{different} module (\autoref{sec:soldering:main}). The two forms
(definite $K|_{\mathfrak{m}}$, indefinite $N$) live on distinct spaces;
there is no tension.
\end{remark}

%%%%%%%%%%%%%%%%%%%%%%%%%%%%%%%%%%%%%%%%%%%%%%%%%%%%%%%%%%%%%%%%%%%%%%
\subsection{Summary of canonical structure}
\label{sec:albert:summary}
%%%%%%%%%%%%%%%%%%%%%%%%%%%%%%%%%%%%%%%%%%%%%%%%%%%%%%%%%%%%%%%%%%%%%%

For reference, \autoref{tab:albert:summary} collects the canonical
algebraic data established above; every downstream section uses these
values.

\begin{table}[ht]
\centering
\renewcommand{\arraystretch}{1.3}
\begin{tabular}{ll}
\hline
Object & Canonical value / property \\
\hline
$\dim_{\mathbb{R}} J_{3}(\mathbb{O})$ & $27$ \\
Trace form $Q=\operatorname{Tr}(X^{2})$ & positive-definite; $F_{4}$-invariant \emph{only} \\
Cubic norm $N=\det_{\mathbb{O}}X$ & indefinite; $E_{6}$-invariant up to character $\lambda(g)$ \\
Sharp map $X^{\#}=\nabla N$ & $E_{6}$-covariant; $X^{\#}=0\Leftrightarrow\operatorname{rank}\le1$ \\
$F_{4}=\operatorname{Aut}(J_{3}(\mathbb{O}))$ & compact, $\dim 52$; preserves $Q$ and $N$ \\
$E_{6(-26)}$ (orbit classifier) & non-compact, $\dim 78$; preserves $N$ only \\
compact $E_{6}$ / $E_{6(-14)}$ (dynamics) & $\dim 78$; max.\ compact $\operatorname{Spin}(10)\times\mathrm{U}(1)$ \\
Vacuum manifold $\mathcal{M}$ & $E_{6}/(\operatorname{Spin}(10)\times\mathrm{U}(1))=\mathrm{EIII}$, $\dim 32$ \\
Adjoint branching & $\mathbf{78}=\mathbf{45}_{0}\oplus\mathbf{1}_{0}\oplus\mathbf{16}_{-3}\oplus\overline{\mathbf{16}}_{+3}$ \\
Fundamental branching & $\mathbf{27}=\mathbf{1}_{+4}\oplus\mathbf{10}_{-2}\oplus\mathbf{16}_{+1}$ \\
Coset tangent $\mathfrak{m}$ & $\mathbf{16}_{-3}\oplus\overline{\mathbf{16}}_{+3}$ (32 real); no singlet \\
Coset metric $K|_{\mathfrak{m}}$ & positive-definite (Euclidean), unique up to scale \\
\hline
\end{tabular}
\caption{Canonical algebraic data of the Albert algebra and its
exceptional structure, as used throughout UCFT.}
\label{tab:albert:summary}
\end{table}

These facts are purely algebraic theorems, with the single
\emph{postulate} being the commitment to the compact real form for the
dynamics (\autoref{post:albert:realform}, part of
\hyperref[post:potential:P1]{P1}). The vacuum-selection theorem that
promotes \eqref{eq:albert:EIII} from ``a distinguished orbit'' to ``the
global minimum of $V$'' is proved in \autoref{sec:potential:main}; the
Lorentzian soldering that uses $N$ to install signature $(-,+,+,+)$ is in
\autoref{sec:soldering:main}; the gauge/clock-field dynamics on
$\mathfrak{m}$ are in \autoref{sec:coset:main}.
% --- end inlined: unit-albert.tex ---
% --- begin inlined: unit-potential.tex ---
\section{The Invariant Potential and Vacuum Selection}
\label{sec:potential:main}

The order parameter of UCFT is a field $\Phi$ valued in the matter representation
$27=J_3(\mathbb O)$ of $E_6$ (Axiom~I). This section presents the $E_6$-invariant potential
$V$ of postulate~\textbf{P1} and proves the vacuum-selection statement~\textbf{P2}: the
global-minimum locus of $V$ is exactly the rank-$1$ primitive-idempotent orbit
$\mathcal M=\mathrm{EIII}=E_6/(\mathrm{Spin}(10)\times U(1))$ of real dimension $32$. We also
show that the symmetric origin $\Phi=0$ is a strict local \emph{maximum}, which supplies the
roll-down dynamics. The potential class is the added postulate; that the chosen member of the
class selects the rank-$1$ orbit is a \textbf{theorem}. Throughout, the dynamics live on the
\emph{compact} real form $E_6$ (equivalently $E_{6(-14)}$, maximal compact
$H=\mathrm{Spin}(10)\times U(1)$); the non-compact $E_{6(-26)}$ enters only as the rigid
orbit classifier (rank and $N$ only). The vacuum geometry and its tangent space are used in
\autoref{sec:soldering:main} (signature) and \autoref{sec:gravity:main} (spectrum, induced
gravity).

\subsection{Albert-algebra setup and the two real forms}
\label{sec:potential:setup}

Let $\mathbb O$ be the real octonions with conjugation $x\mapsto\bar x$ and norm
$n(x)=x\bar x$. The real Albert algebra is the $27$-dimensional Euclidean (formally real)
Jordan algebra
\begin{equation}
\label{eq:potential:albert}
J_3(\mathbb O)=\left\{
X=\begin{pmatrix}\xi_1 & x_3 & \bar x_2\\ \bar x_3 & \xi_2 & x_1\\ x_2 & \bar x_1 & \xi_3\end{pmatrix}
:\ \xi_i\in\mathbb R,\ x_i\in\mathbb O\right\},
\qquad X\circ Y=\tfrac12(XY+YX).
\end{equation}
Its two basic invariants are the \emph{trace form} (a Euclidean, positive-definite bilinear
form) and the \emph{cubic norm}:
\begin{equation}
\label{eq:potential:invariants}
\langle X,Y\rangle=\operatorname{Tr}(X\circ Y),\qquad
Q(X):=\operatorname{Tr}(X^2)=\langle X,X\rangle\ge0,\qquad
N(X)=\det X\ \ (\text{cubic}).
\end{equation}
With the linear trace $\operatorname{tr}(X)=\xi_1+\xi_2+\xi_3$ and identity
$E=\operatorname{diag}(1,1,1)$, every $X$ obeys the cubic characteristic identity
$X^{\circ3}-\operatorname{tr}(X)\,X^{\circ2}+\sigma_2(X)\,X-N(X)\,E=0$ with
$\sigma_2(X)=\tfrac12(\operatorname{tr}(X)^2-Q(X))$. The Jordan eigenvalues
$\lambda_1,\lambda_2,\lambda_3\in\mathbb R$ satisfy
$\operatorname{tr}(X)=\sum_i\lambda_i$, $Q(X)=\sum_i\lambda_i^2$,
$N(X)=\lambda_1\lambda_2\lambda_3$.

The \emph{sharp map} (Jordan adjoint) $X\mapsto X^\#$ is the unique quadratic map with
\begin{equation}
\label{eq:potential:sharp}
X^\#\circ X=N(X)\,E,\qquad (X^\#)^\#=N(X)\,X,\qquad X^\#=\nabla N(X),
\end{equation}
and is the matrix $2\times2$-cofactor (adjugate): in a Jordan frame
$\operatorname{diag}(\lambda_1,\lambda_2,\lambda_3)^\#
=\operatorname{diag}(\lambda_2\lambda_3,\lambda_3\lambda_1,\lambda_1\lambda_2)$.

\paragraph{The two $E_6$'s (a locked distinction).}
The reduced structure group is the non-compact $E_{6(-26)}$ (dim $78$); it preserves $N$ only
up to a scalar character, $N(gX)=\lambda(g)N(X)$, and covariantly transforms the sharp map.
It does \emph{not} preserve $Q=\operatorname{Tr}(X^2)$; indeed $\operatorname{Aut}(J_3(\mathbb O))=F_4$
(compact, dim $52$) is exactly the subgroup preserving both $Q$ and $N$. We use $E_{6(-26)}$
solely to classify orbits of the real $27$ by \textbf{rank and $N$}. The \emph{dynamics}, by
contrast, live on the \textbf{compact} $E_6$: there $H=\mathrm{Spin}(10)\times U(1)$ is
compact, so the coset metric is positive-definite and the gauge sector unitary; on the
compact form $Q=\langle X,X\rangle$ is the invariant trace form and the norm character is
identically $1$. Orbits are classified by rank and $N$ (\emph{not} by $(Q,N)$):
\begin{equation}
\label{eq:potential:rank}
\operatorname{rank}(X)\le1\iff X^\#=0,\qquad
\operatorname{rank}(X)\le2\iff N(X)=0,\qquad
\operatorname{rank}(X)=3\iff N(X)\neq0.
\end{equation}

\subsection{The potential class (Postulate P1)}
\label{sec:potential:class}

\begin{postulate}[Potential class, P1]
\label{post:potential:P1}
The order parameter $\Phi$ is governed by the $E_6$-invariant polynomial potential
\begin{equation}
\label{eq:potential:v}
V(X)=\alpha\,\langle X^\#,X^\#\rangle
+\beta\,\bigl(\operatorname{Tr}(X^2)-c_2\bigr)^2
+\gamma\,\bigl(\operatorname{Tr}(X^2)\bigr)^k,
\qquad \alpha,\beta,\gamma>0,\ \ k\ge2\ \text{even},
\end{equation}
built from the cubic norm (through the sharp map $X^\#=\nabla N$) and the trace form $Q$.
On the compact real form every term is a genuine $E_6$-invariant.
\end{postulate}

The first term is genuinely $E_6$-covariant and vanishes precisely on rank $\le1$; the
$Q$-terms are $F_4$-invariant selectors that fix the scale and the representative. Postulate
P1 asserts that one is \emph{allowed} to write such a potential; the theorems below establish
\emph{what it selects}.

\subsection{The vacuum-selection theorem (P1 \texorpdfstring{$\Rightarrow$}{=>} P2)}
\label{sec:potential:theorem}

\begin{lemma}[Zero locus of the sharp map]
\label{lem:potential:L1}
For $X\in J_3(\mathbb O)$, \ $X^\#=0\iff\operatorname{rank}(X)\le1$. The normalized
($\operatorname{tr}=1$) rank-$1$ elements are exactly the primitive idempotents, forming the
single $F_4$-orbit $F_4/\mathrm{Spin}(9)\cong\mathbb{OP}^2$ (real dim $16$); the
projective $E_6$-orbit of a rank-$1$ line in $\mathbb P(27_{\mathbb C})$ is the complex
Cayley plane $\mathrm{EIII}=E_6/(\mathrm{Spin}(10)\times U(1))$, real dim $32$.
\end{lemma}

\begin{proof}
In a Jordan frame $X=\operatorname{diag}(\lambda_1,\lambda_2,\lambda_3)$, the sharp map gives
$X^\#=\operatorname{diag}(\lambda_2\lambda_3,\lambda_3\lambda_1,\lambda_1\lambda_2)$, which
vanishes iff at most one $\lambda_i\neq0$, i.e. $\operatorname{rank}(X)\le1$; $F_4$-equivariance
transports this to all of $J_3(\mathbb O)$. The orbit statements are the standard
identifications of the real rank-$1$ idempotents ($\mathbb{OP}^2$, dim $16$) and the minimal
$E_6$-orbit / complex Cayley plane (EIII, dim $32$).
\end{proof}

\begin{theorem}[Positivity of the sharp term]
\label{thm:potential:T1}
For all $X\in J_3(\mathbb O)$, \ $\langle X^\#,X^\#\rangle\ge0$, with equality iff $X^\#=0$
iff $\operatorname{rank}(X)\le1$. Explicitly
\begin{equation}
\label{eq:potential:sharpsq}
\langle X^\#,X^\#\rangle=(\lambda_2\lambda_3)^2+(\lambda_3\lambda_1)^2+(\lambda_1\lambda_2)^2
=\sigma_2(X)^2-2\,\operatorname{tr}(X)\,N(X)\ \ (\text{a sum of squares}).
\end{equation}
\end{theorem}

\begin{proof}
The trace form is positive-definite on the Euclidean Jordan algebra, so $\langle Y,Y\rangle\ge0$
with equality iff $Y=0$; take $Y=X^\#$ and apply \autoref{lem:potential:L1}. The frame
expression \eqref{eq:potential:sharpsq} is manifestly a sum of squares.
\end{proof}

\begin{theorem}[Compact-$E_6$ invariance of the sharp term]
\label{thm:potential:T2}
The scalar $\langle X^\#,X^\#\rangle$ is invariant under the compact real form $E_6$ (and
$E_{6(-14)}$): \ $\langle(gX)^\#,(gX)^\#\rangle=\langle X^\#,X^\#\rangle$ for all
$g\in E_6^{\,\mathrm{compact}}$.
\end{theorem}

\begin{proof}
On the compact form the $27$ is unitary, with invariant positive-definite form whose real
part is the trace form; hence the compact $E_6$ preserves $\langle\cdot,\cdot\rangle$. As a
degree-$4$ relative invariant, $\langle X^\#,X^\#\rangle$ transforms by a power of the norm
character $\lambda(g)$; but a continuous homomorphism from a compact connected group to
$\mathbb R_{>0}$ is identically $1$, so $\lambda\equiv1$ and the relative invariant is a
genuine invariant. (Under $E_{6(-26)}$, $\lambda$ is unbounded and $\langle\cdot,\cdot\rangle$
is not preserved; this is precisely why the dynamics use the compact form.)
\end{proof}

\paragraph{Status of $Q$.}
$Q=\operatorname{Tr}(X^2)$ is $F_4$-invariant only and is \emph{not} preserved by the
non-compact structure group; we never claim otherwise, nor that $(Q,N)$ separates $E_6$-orbits.
On the compact form, however, $Q=\langle X,X\rangle$ \emph{is} the invariant trace form, so
$V$ in \eqref{eq:potential:v} is a genuine compact-$E_6$ invariant potential.

\begin{lemma}[The selector has a unique positive minimizer]
\label{lem:potential:L2}
Set $W(q)=\beta(q-c_2)^2+\gamma q^k$ for $q\ge0$, $\beta,\gamma>0$, $k\ge2$ even. There is a
$c_2^{\,\star}>0$ such that for $c_2>c_2^{\,\star}$, $W$ is strictly convex on $[0,\infty)$ and
attains its global minimum at a unique interior point $q_\star\in(0,c_2)$ with $W'(q_\star)=0$,
$W''(q_\star)>0$, and $W(q_\star)<W(0)=\beta c_2^2$. Moreover
\begin{equation}
\label{eq:potential:qstar}
q_\star=c_2-\frac{\gamma k}{2\beta}\,c_2^{\,k-1}+O(c_2^{\,2k-3}),
\qquad q_\star\xrightarrow[\gamma\to0^+]{}c_2.
\end{equation}
\end{lemma}

\begin{proof}
$W(0)=\beta c_2^2>0$ and $W(q)\to+\infty$, so a minimum exists on a compact set. Since
$W''(q)=2\beta+\gamma k(k-1)q^{k-2}>0$ for all $q\ge0$ ($k\ge2$ even), $W$ is strictly convex;
with $W'(0)=-2\beta c_2<0$ it has a unique stationary point $q_\star>0$, the global minimizer,
and $W(q_\star)<W(0)$. As $W'(c_2)=\gamma k c_2^{k-1}>0$, the root lies below $c_2$; the
expansion follows by solving $2\beta(q_\star-c_2)=-\gamma k q_\star^{k-1}$ perturbatively.
\end{proof}

The selector thus \emph{fixes the scale}: among rank-$1$ elements $X=\lambda P$ (with
$Q=\lambda^2$, $N=0$, $X^\#=0$), the minimizers have $\lambda^2=q_\star$, a fixed nonzero
norm; in particular the origin ($q=0$, $W(0)=\beta c_2^2>W(q_\star)$) is excluded.

\begin{lemma}[Coercivity requires the compact form]
\label{lem:potential:L3}
On $J_3(\mathbb O)$ with the compact-$E_6$ Euclidean norm $\|X\|^2=Q(X)$, the potential $V$ is
coercive and bounded below: $V(X)\ge\gamma\|X\|^{2k}\to+\infty$. On a non-compact real-form
orbit, where $Q$ is not preserved and not a proper exhaustion, $V$ need not be coercive.
\end{lemma}

\begin{proof}
$V\ge\gamma Q^k=\gamma\|X\|^{2k}$ from \eqref{eq:potential:v} ($\alpha$-term and the square
are $\ge0$), giving coercivity on the compact form. Under $E_{6(-26)}$ the character
$\lambda(g)$ is unbounded, so a fixed $X$ with $N(X)\neq0$ can be driven to $Q\to0$ along a
fixed $N$-level set; $Q$ is not proper there and sublevel sets are unbounded.
\end{proof}

\begin{theorem}[Global-minimum locus is the rank-$1$ EIII orbit; P2]
\label{thm:potential:T4}
With $\alpha,\beta,\gamma>0$, $k\ge2$ even, and $c_2>c_2^{\,\star}$
(\autoref{lem:potential:L2}), every global minimizer $X_\star$ of $V$ on $J_3(\mathbb O)$
satisfies $X_\star^\#=0$ and $\operatorname{Tr}(X_\star^2)=q_\star$; equivalently
\begin{equation}
\label{eq:potential:argmin}
\arg\min_{X\in J_3(\mathbb O)}V(X)
=\{\,X=\lambda P:\ P\ \text{primitive idempotent},\ \lambda^2=q_\star\,\},
\end{equation}
whose projective $E_6$-orbit is the vacuum manifold
\begin{equation}
\label{eq:potential:vacuum}
\mathcal M=\mathrm{EIII}=\frac{E_6}{\mathrm{Spin}(10)\times U(1)},
\qquad \dim_{\mathbb R}\mathcal M=32.
\end{equation}
This proves postulate~\textbf{P2} from \autoref{post:potential:P1}.
\end{theorem}

\begin{proof}
By \autoref{lem:potential:L3} a global minimum is attained. For any $X$,
\begin{equation*}
V(X)=\underbrace{\alpha\,\langle X^\#,X^\#\rangle}_{\ge0}+W(Q(X))\ge W(Q(X))\ge W(q_\star),
\end{equation*}
using \autoref{thm:potential:T1} and \autoref{lem:potential:L2}. Equality in the first step
forces $X^\#=0$, i.e. $\operatorname{rank}(X)\le1$ (\autoref{lem:potential:L1}); equality in
the second forces $Q(X)=q_\star>0$, excluding $X=0$ and forcing rank exactly $1$. Both hold
iff $X=\lambda P$ with $\lambda^2=q_\star$, and every such $X$ achieves $V=W(q_\star)$. The
projective orbit of these rank-$1$ idempotents is EIII (\autoref{lem:potential:L1}).
\end{proof}

\paragraph{Why rank $2$ and rank $3$ strictly lose.}
For $\operatorname{rank}(X)\ge2$, \autoref{lem:potential:L1} gives $X^\#\neq0$, so
$\langle X^\#,X^\#\rangle>0$ strictly (\autoref{thm:potential:T1}) and
$V(X)>W(q_\star)$ by a deficit at least $\alpha\langle X^\#,X^\#\rangle>0$. In particular
$\operatorname{diag}(1,1,1)$ (rank $3$, $\langle X^\#,X^\#\rangle=3>0$, orbit $E_6/F_4$ of
real dimension $26$) is \emph{not} a minimizer. The vacuum is the rank-$1$ idempotent EIII
orbit (dim $32$), \emph{not} the $E_6/F_4$ orbit of $\operatorname{diag}(1,1,1)$ (dim $26$).

\begin{remark}[Coset tangent and metric]
\label{rmk:potential:tangent}
The tangent space to $\mathcal M$ at the base idempotent is the coset module
\begin{equation}
\label{eq:potential:cosettangent}
\mathfrak m=16_{-3}\oplus\overline{16}_{+3}\quad(32\ \text{real}),
\end{equation}
read off from the adjoint branching $78=45_0\oplus1_0\oplus16_{-3}\oplus\overline{16}_{+3}$
under $H=\mathrm{Spin}(10)\times U(1)$. It is real-irreducible of complex type (commutant
$\mathbb C$) and contains \emph{no} $\mathrm{SO}(10)$ singlet; the $U(1)$ charge forbids the
symmetric $16\times16$ and $\overline{16}\times\overline{16}$ invariants, leaving a unique
charge-$0$ cross-pairing. The induced coset metric $K|_{\mathfrak m}$ is therefore unique and
positive-definite (Schur), and serves as the kinetic/target-space metric for the $32$ clock
fields (Goldstones). It carries \emph{no} Lorentzian content: the $(1,3)$ signature is added
structure (P3), soldered from the cubic norm via $H_2(\mathbb O)\cong\mathbb R^{1,9}$, not a
Killing-form consequence; see \autoref{sec:soldering:main}.
\end{remark}

\subsection{The symmetric origin is a strict local maximum}
\label{sec:potential:origin}

\begin{theorem}[Origin is a strict local maximum]
\label{thm:potential:T5}
At $X=0$,
\begin{equation}
\label{eq:potential:hess0}
\nabla V(0)=0,\qquad
\operatorname{Hess}V(0)=-4\beta c_2\,\mathbb 1\ \prec\ 0
\quad(\text{negative-definite, }\beta,c_2>0),
\end{equation}
so $X=0$ is a strict local maximum, all $27$ Hessian eigenvalues equal to $-4\beta c_2<0$.
Since $V(0)=\beta c_2^2>W(q_\star)$, the rank-$1$ orbit beats the origin locally and globally.
\end{theorem}

\begin{proof}
$V$ has no linear part. The $\alpha\langle X^\#,X^\#\rangle$ term is homogeneous of degree $4$
and the $\gamma Q^k$ term of degree $2k\ge4$, so both have vanishing Hessian at $0$. Only the
quadratic piece $-2\beta c_2 Q$ of $\beta(Q-c_2)^2$ contributes; with $Q$ having Hessian
$2\,\mathbb 1$ in the trace metric this gives $\operatorname{Hess}V(0)=-4\beta c_2\,\mathbb 1$,
negative-definite.
\end{proof}

\begin{theorem}[Roll-down to $\mathcal M$]
\label{thm:potential:T6}
Under gradient flow $\dot X=-\nabla V(X)$ (or the EOM with damping), the origin is linearly
unstable with growth rate $4\beta c_2>0$ in all $27$ directions; $V$ is a strict Lyapunov
function, coercivity (\autoref{lem:potential:L3}) confines trajectories to compact sublevel
sets, and by LaSalle the $\omega$-limit set lies in $\{\nabla V=0\}$. The unique stable
critical set is the global minimum $\mathcal M$ (\autoref{thm:potential:T4}); the origin
(maximum, \autoref{thm:potential:T5}) and the rank-$2$/rank-$3$ saddles ($X^\#\neq0$) are all
unstable. Hence generic initial data roll down to $\mathcal M=\mathrm{EIII}$.
\end{theorem}

\begin{proof}
Linearization at $0$ uses $\operatorname{Hess}V(0)\prec0$ (\autoref{thm:potential:T5}), giving
positive growth in every direction. $\dot V=-\|\nabla V\|^2\le0$ vanishes only at critical
points; LaSalle plus coercivity confine $\omega$-limits to the critical set, of which only
$\mathcal M$ is Lyapunov-stable. The basin of $\mathcal M$ is open and dense.
\end{proof}

\paragraph{Tree-level Goldstone consistency.}
On $\mathcal M$, $V$ is $E_6$-invariant and constant along the orbit, so
$\operatorname{Hess}V|_{\mathfrak m}=0$ at tree level: the $32$ clock fields are exact
Goldstones. Radiative (Coleman--Weinberg) corrections are $H$-invariant; by Schur on the
real-irreducible module $\mathfrak m$ they give \emph{one} common mass on all $32$ directions.
Of these, $28$ are physical massive scalars (single common mass) and $4$ are the
gauge/soldering modes eaten by diffeomorphism and local Lorentz; this $28+4$ split and the
common radiative mass are carried into the spectrum and induced-gravity analysis of
\autoref{sec:gravity:main}. The origin's $27$ negative modes and the vacuum's flat coset
directions are not in tension: they live at different points of field space.

\subsection{Summary of the logical chain}
\label{sec:potential:summary}

\begin{center}
\begin{tabular}{lll}
\toprule
step & statement & type \\
\midrule
\autoref{lem:potential:L1} & $X^\#=0\iff\operatorname{rank}(X)\le1$; rank-$1$ orbit $=$ EIII (dim $32$) & Theorem \\
\autoref{thm:potential:T1} & $\langle X^\#,X^\#\rangle\ge0$, $=0\iff\operatorname{rank}\le1$ & Theorem \\
\autoref{thm:potential:T2} & $\langle X^\#,X^\#\rangle$ is compact-$E_6$ invariant & Theorem \\
\autoref{lem:potential:L2} & selector $W$ strictly convex, unique $q_\star>0$, $W(q_\star)<W(0)$ & Theorem \\
\autoref{lem:potential:L3} & coercivity holds on the compact form, fails on non-compact orbits & Theorem \\
\autoref{thm:potential:T4} & \textbf{global-min locus $=$ rank-$1$ EIII orbit, dim $32$ (P2)} & \textbf{Theorem} \\
\autoref{thm:potential:T5} & $\operatorname{Hess}V(0)=-4\beta c_2\,\mathbb1\prec0$: origin strict max & Theorem \\
\autoref{thm:potential:T6} & roll-down: generic flows attract to $\mathcal M$ & Theorem \\
\autoref{post:potential:P1} & the potential class is added structure & Postulate \\
\bottomrule
\end{tabular}
\end{center}

Given the postulated potential class P1, \autoref{thm:potential:T4} proves P2: the unique
global-minimum locus of $V$ is the rank-$1$ primitive-idempotent orbit
$\mathcal M=\mathrm{EIII}=E_6/(\mathrm{Spin}(10)\times U(1))$, real dimension $32$. The sharp
term enforces rank $\le1$; the $Q$-selector fixes the nonzero scale and excludes the origin;
the compact real form supplies coercivity and a definite invariant metric; and the origin is
a strict local maximum driving the roll-down. No rank-$2$/rank-$3$ point --- in particular
none on the $E_6/F_4$ orbit of $\operatorname{diag}(1,1,1)$ (dim $26$) --- can be a minimizer.
% --- end inlined: unit-potential.tex ---
% --- begin inlined: unit-coset.tex ---
%%%%%%%%%%%%%%%%%%%%%%%%%%%%%%%%%%%%%%%%%%%%%%%%%%%%%%%%%%%%%%%%%%%%%%%%%%
%% UNIT: coset  --  Coset Geometry and the Clock Fields
%% Self-contained section fragment (no preamble / documentclass).
%% Cross-references to other units use stable labels via \autoref/\cref.
%%%%%%%%%%%%%%%%%%%%%%%%%%%%%%%%%%%%%%%%%%%%%%%%%%%%%%%%%%%%%%%%%%%%%%%%%%
\section{Coset Geometry and the Clock Fields}
\label{sec:coset:main}

The vacuum-selection theorem (\autoref{thm:vacuum:T4}) establishes that the
global minimum of the $E_{6}$-invariant potential of Postulate~P1 is the closed
rank-$1$ primitive-idempotent orbit in $\mathbb P(\mathbf{27}_{\mathbb C})$, namely
the Hermitian symmetric space
\begin{equation}\label{eq:coset:vacuum}
  \mathcal M \;=\; \frac{E_{6}}{\Spin(10)\times U(1)} \;=\; \mathrm{EIII},
  \qquad \dim_{\mathbb R}\mathcal M = 32 .
\end{equation}
Here and throughout the \emph{dynamics} are carried by the \emph{compact} real
form of $E_{6}$ (equivalently $E_{6(-14)}$, whose maximal compact subgroup is
exactly $H = \Spin(10)\times U(1)$); the non-compact $E_{6(-26)}$ enters only as
the rigid classifier of orbits in the real $\mathbf{27}$ by rank and cubic norm
$N$ (\autoref{sec:vacuum}). Because $H$ is compact, the isotropy
representation carries an invariant \emph{positive-definite} metric, which we
will identify as the kinetic/target-space metric of the $\sigma$-model.

This section is geometric and tree-level. We (i) identify the tangent space
$\mathfrak m$ of $\mathcal M$ and prove, from the $U(1)$ charge structure of a
Hermitian symmetric space, that the invariant metric $K|_{\mathfrak m}$ is unique
and positive-definite; (ii) build the nonlinear $\sigma$-model from the
Maurer--Cartan form, with its $32$ Goldstone \emph{clock fields}; and (iii)
record that at tree level all $32$ clock fields are exact massless Goldstone
modes. The Lorentzian frame of four of them (Postulate~P3) is supplied in
\autoref{sec:soldering}; their one-loop spectrum (Schur, $28$ massive $+\,4$
gauge/soldering) in \autoref{sec:gravity}.

%%%%%%%%%%%%%%%%%%%%%%%%%%%%%%%%%%%%%%%%%%%%%%%%%%%%%%%%%%%%%%%%%%%%%%%%%%
\subsection{The tangent space \texorpdfstring{$\mathfrak m$}{m}}
\label{sec:coset:tangent}
%%%%%%%%%%%%%%%%%%%%%%%%%%%%%%%%%%%%%%%%%%%%%%%%%%%%%%%%%%%%%%%%%%%%%%%%%%

The orbit \eqref{eq:coset:vacuum} is defined by the Cartan involution
$\vartheta$ whose $(+1)$-eigenspace is $\mathfrak h = \mathfrak{so}(10)\oplus
\mathfrak u(1)$ and whose $(-1)$-eigenspace $\mathfrak m$ is the tangent space at
the base point. The branching of the adjoint $\mathbf{78}$ of $E_{6}$ under
$H = \Spin(10)\times U(1)$ is the involution datum
\begin{equation}\label{eq:coset:adjoint}
  \mathbf{78}
  \;=\; \underbrace{\mathbf{45}_{0}\oplus\mathbf 1_{0}}_{\mathfrak h}
       \;\oplus\;
        \underbrace{\mathbf{16}_{-3}\oplus\overline{\mathbf{16}}_{+3}}_{\mathfrak m},
\end{equation}
so that
\begin{equation}\label{eq:coset:tangent}
  \mathfrak m \;=\; \mathbf{16}_{-3}\oplus\overline{\mathbf{16}}_{+3},
  \qquad \dim_{\mathbb R}\mathfrak m = 32 .
\end{equation}
The two summands are the chiral spinor and antispinor of $\Spin(10)$, carrying
opposite $U(1)$ charges $\mp 3$. For comparison, the matter representation
branches as
\begin{equation}\label{eq:coset:fund}
  \mathbf{27} \;=\; \mathbf 1_{+4}\oplus\mathbf{10}_{-2}\oplus\mathbf{16}_{+1},
\end{equation}
the $\mathbf{16}_{+1}$ of which will house one chiral family
(\autoref{sec:sm}).

\begin{remark}\label{rem:coset:nosinglet}
$\mathfrak m$ contains \emph{no} $\Spin(10)$ singlet. In particular $\mathcal M$
is \emph{not} the orbit $E_{6}/F_{4}$ of $\mathrm{diag}(1,1,1)$ (real dimension
$26$), and $\mathfrak m$ admits no decomposition of the spurious form
$2\times\mathbf{10}\oplus\mathbf 4$ or $\mathbf{10}\oplus\overline{\mathbf{10}}
\oplus\mathbf 4$. The four Lorentzian frame directions of
\autoref{sec:soldering} are \emph{not} an $H$-invariant subspace of
$\mathfrak m$; they are selected by the soldering map $\sigma$ (Postulate~P3),
not by the branching \eqref{eq:coset:tangent}.
\end{remark}

The decisive structural fact is the type of $\mathfrak m$ as a real
$H$-representation.

\begin{theorem}[Real-irreducibility of $\mathfrak m$]\label{thm:coset:irrep}
As a real representation of $H = \Spin(10)\times U(1)$, the tangent space
$\mathfrak m = \mathbf{16}_{-3}\oplus\overline{\mathbf{16}}_{+3}$ is irreducible
of \emph{complex type}: its commutant algebra is $\mathbb C$.
\end{theorem}

\begin{proof}
The complexified module $\mathfrak m_{\mathbb C}$ is the sum of the two distinct
complex-irreducible $H$-modules $\mathbf{16}_{-3}$ and
$\overline{\mathbf{16}}_{+3}$, which are interchanged by complex conjugation
(the $U(1)$ charge flips sign and $\mathbf{16}\leftrightarrow\overline{\mathbf{16}}$).
A real-irreducible module whose complexification is $V\oplus\overline V$ with
$V\not\cong\overline V$ is irreducible of \emph{complex type}; its commutant is
$\mathbb C$, acting as the complex structure $J$ that rotates the
$\mathbf{16}_{-3}$ phase by $e^{i\theta}$ and the $\overline{\mathbf{16}}_{+3}$
phase by $e^{-i\theta}$. (This $J$ is precisely the action of the $\mathfrak u(1)$
generator $\mathbf 1_{0}$ in \eqref{eq:coset:adjoint}; it is the integrable
almost-complex structure that makes $\mathcal M=\mathrm{EIII}$ Hermitian
symmetric.) No real $H$-invariant subspace of $\mathfrak m$ other than $0$ and
$\mathfrak m$ exists, since any such subspace would complexify to a proper
nonzero $H$-submodule of $\mathbf{16}_{-3}\oplus\overline{\mathbf{16}}_{+3}$
stable under conjugation, hence to $\mathbf{16}_{-3}\oplus\overline{\mathbf{16}}_{+3}$
itself.
\end{proof}

\begin{remark}\label{rem:coset:hermitian}
\autoref{thm:coset:irrep} replaces, with a cleaner argument, the older
``single weight orbit'' claim of $\mathrm{lem}{:}\mathrm{Irrep}$: irreducibility
follows directly from the Hermitian-symmetric-space structure (a single
$\mathbb C$-line $\mathfrak m^{(1,0)} = \mathbf{16}_{-3}$ and its conjugate),
without inspecting the root diagram.
\end{remark}

%%%%%%%%%%%%%%%%%%%%%%%%%%%%%%%%%%%%%%%%%%%%%%%%%%%%%%%%%%%%%%%%%%%%%%%%%%
\subsection{The invariant metric}
\label{sec:coset:metric}
%%%%%%%%%%%%%%%%%%%%%%%%%%%%%%%%%%%%%%%%%%%%%%%%%%%%%%%%%%%%%%%%%%%%%%%%%%

\begin{theorem}[Unique positive-definite coset metric]\label{thm:coset:metric}
Up to an overall positive constant $f^{2}$, there is a unique $E_{6}$-invariant
symmetric bilinear form on $\mathfrak m$,
\begin{equation}\label{eq:coset:metric}
  \mathfrak g_{\alpha\beta} \;=\; f^{2}\,K_{\alpha\beta},
  \qquad
  K_{\alpha\beta} \;=\; \bigl(K_{E_{6}}\bigr)\big|_{\mathfrak m},
\end{equation}
the restriction to $\mathfrak m$ of the Killing form of (compact) $E_{6}$. For
$f^{2}>0$ this metric is positive-definite (Euclidean signature
$(32,0)$); it is the kinetic/target-space metric of the $\sigma$-model. Here
$[f] = \text{mass}$.
\end{theorem}

\begin{proof}
\emph{Existence and the charge selection rule.} An $H$-invariant symmetric form
on $\mathfrak m$ is an element of $\bigl(\Sym^{2}\mathfrak m^{*}\bigr)^{H}$.
Complexifying, $\Sym^{2}\mathfrak m_{\mathbb C}^{*}$ decomposes through
\begin{equation}\label{eq:coset:sym2}
  \Sym^{2}\!\bigl(\mathbf{16}_{-3}\oplus\overline{\mathbf{16}}_{+3}\bigr)
  = \underbrace{\Sym^{2}\mathbf{16}_{-3}}_{U(1)\text{ charge }-6}
    \;\oplus\;
    \underbrace{\bigl(\mathbf{16}_{-3}\otimes\overline{\mathbf{16}}_{+3}\bigr)}_{\text{charge }0}
    \;\oplus\;
    \underbrace{\Sym^{2}\overline{\mathbf{16}}_{+3}}_{\text{charge }+6}.
\end{equation}
The first and third summands carry nonzero $U(1)$ charge ($\mp 6$) and therefore
contain \emph{no} $H$-singlet: the symmetric $\mathbf{16}\times\mathbf{16}$ and
$\overline{\mathbf{16}}\times\overline{\mathbf{16}}$ invariants are forbidden by
the $U(1)$ charge. Only the charge-$0$ cross-pairing $\mathbf{16}_{-3}\otimes
\overline{\mathbf{16}}_{+3}\supset\mathbf 1_{0}$ contributes an $H$-singlet, and
it does so exactly once (the $\Spin(10)$ contraction
$\langle\,\cdot\,,\,\overline{\,\cdot\,}\,\rangle$ is unique up to scale). Hence
$\dim\bigl(\Sym^{2}\mathfrak m^{*}\bigr)^{H} = 1$. This unique invariant is the
Killing restriction $K|_{\mathfrak m}$.

\emph{Uniqueness (Schur).} By \autoref{thm:coset:irrep}, $\mathfrak m$ is
real-irreducible, so Schur's lemma gives at most a one-dimensional space of
invariant symmetric forms; the count above shows it is exactly one-dimensional.
Any two are proportional, $\mathfrak g_{\alpha\beta} = c\,K_{\alpha\beta}$.

\emph{Definiteness.} On the compact real form of $E_{6}$ the Killing form is
negative-definite, so $-K_{E_{6}}$ is a positive-definite invariant inner
product on $\mathfrak e_{6}$; its restriction to $\mathfrak m$ is therefore
positive-definite. Fixing the sign and scale by the kinetic normalization,
$c = f^{2}>0$, gives \eqref{eq:coset:metric}.
\end{proof}

\begin{remark}\label{rem:coset:nolorentz}
\autoref{thm:coset:metric} forbids any $(4,28)$ or $(1,3)$ ``Killing/coset
signature.'' The compact-$E_{6}$ Killing form on $\mathfrak m$ has signature
$(32,0)$ and zero Lorentzian content. Lorentzian signature enters \emph{only}
through the soldering Postulate~P3 (\autoref{sec:soldering}), rooted in the cubic
norm $N$ on $J_{3}(\mathbb O)$ via $H_{2}(\mathbb O)\cong\mathbb R^{1,9}$, not in
the coset geometry.
\end{remark}

%%%%%%%%%%%%%%%%%%%%%%%%%%%%%%%%%%%%%%%%%%%%%%%%%%%%%%%%%%%%%%%%%%%%%%%%%%
\subsection{The Maurer--Cartan form and the clock-field \texorpdfstring{$\sigma$}{sigma}-model}
\label{sec:coset:sigma}
%%%%%%%%%%%%%%%%%%%%%%%%%%%%%%%%%%%%%%%%%%%%%%%%%%%%%%%%%%%%%%%%%%%%%%%%%%

Following the coset construction of nonlinear realizations, choose a local
representative $g(x)\in E_{6}$ of the map $x\mapsto\pi(x)\in\mathcal M$. The
$32$ broken generators $X_{\alpha}\in\mathfrak m$ furnish local coordinates
$\phi^{\alpha}(x)$ on $\mathcal M$; these are the \emph{clock fields}. The
left-invariant Maurer--Cartan one-form
\begin{equation}\label{eq:coset:MC}
  \theta_{\mu}(x) \;=\; g^{-1}(x)\,\partial_{\mu}g(x)\;\in\;\mathfrak e_{6}
\end{equation}
decomposes into unbroken ($\mathfrak h$) and broken ($\mathfrak m$) parts,
\begin{equation}\label{eq:coset:MCsplit}
  \theta_{\mu}
   \;=\; \underbrace{A_{\mu}{}^{i}\,H_{i}}_{\mathfrak h\text{-connection}}
        \;+\;
         \underbrace{e_{\mu}{}^{\alpha}\,X_{\alpha}}_{\mathfrak m\text{-vielbein}},
  \qquad i = 1,\dots,46,\;\; \alpha = 1,\dots,32 .
\end{equation}
Under a local $H$ transformation $g(x)\mapsto g(x)h(x)^{-1}$, the $\mathfrak
h$-part $A_{\mu}$ transforms as a $\Spin(10)\times U(1)$ gauge connection and the
$\mathfrak m$-part $e_{\mu}{}^{\alpha}$ transforms covariantly; the composite
gauge sector built from $A_{\mu}$ is developed in \autoref{sec:gauge}, and the
soldering of four components of $e_{\mu}{}^{\alpha}$ to a spacetime vierbein in
\autoref{sec:soldering}.

The $H$-covariant derivative of the clock fields and the field strength of the
composite connection are
\begin{equation}\label{eq:coset:covD}
  D_{\mu}\phi^{\alpha}
    = \partial_{\mu}\phi^{\alpha}
      + A_{\mu}{}^{i}\,(t_{i})^{\alpha}{}_{\beta}\,\phi^{\beta},
  \qquad
  F_{\mu\nu}
    = \partial_{\mu}A_{\nu} - \partial_{\nu}A_{\mu} + [A_{\mu},A_{\nu}],
\end{equation}
with $(t_{i})^{\alpha}{}_{\beta}$ the matrices of $\mathfrak h$ acting on
$\mathfrak m = \mathbf{16}_{-3}\oplus\overline{\mathbf{16}}_{+3}$. The most
general two-derivative action consistent with global $E_{6}$, local $H$, and
(after Postulate~P3) local Lorentz invariance is fixed up to the couplings to be
the gauged nonlinear $\sigma$-model
\begin{equation}\label{eq:coset:action}
  S_{\sigma}[\phi,A]
   = \int \dd^{4}x\,\sqrt{-g}\;
     \frac12\, g^{\mu\nu}\,\mathfrak g_{\alpha\beta}\,
       D_{\mu}\phi^{\alpha}\,D_{\nu}\phi^{\beta}
     \;-\;\frac14\,\hat g^{-2}\,\Tr\!\bigl(F_{\mu\nu}F^{\mu\nu}\bigr),
\end{equation}
with the target metric $\mathfrak g_{\alpha\beta} = f^{2}K_{\alpha\beta}$ of
\autoref{thm:coset:metric} and signature convention $(-,+,+,+)$. The matter
(Weyl-fermion) and Yukawa sectors are added in
\autoref{sec:sm} (Postulates~P4, P6); the gauge coupling $\hat g^{2}$,
the dimensionless $\xi,\hat y^{2},\kappa$, and the running of $f$ are studied in
\autoref{sec:gravity}.

\begin{lemma}[Healthy kinetic term]\label{lem:coset:positivity}
For $f^{2}>0$ the clock-field kinetic term in \eqref{eq:coset:action} is
positive-definite, and the classical $\sigma$-model energy is bounded below.
\end{lemma}

\begin{proof}
Immediate from \autoref{thm:coset:metric}: $\mathfrak g_{\alpha\beta}$ is
positive-definite on $\mathfrak m$, so for the spatial gradient and time
derivative the energy density
$\tfrac12\mathfrak g_{\alpha\beta}\bigl(\dot\phi^{\alpha}\dot\phi^{\beta}
+\partial_{i}\phi^{\alpha}\partial_{i}\phi^{\beta}\bigr)\ge 0$, and the gauge
term enters at most quadratically. (This is where the \emph{compact} real form
is essential; a non-compact form would render $K|_{\mathfrak m}$ indefinite.)
\end{proof}

%%%%%%%%%%%%%%%%%%%%%%%%%%%%%%%%%%%%%%%%%%%%%%%%%%%%%%%%%%%%%%%%%%%%%%%%%%
\subsection{Tree-level massless Goldstones}
\label{sec:coset:goldstone}
%%%%%%%%%%%%%%%%%%%%%%%%%%%%%%%%%%%%%%%%%%%%%%%%%%%%%%%%%%%%%%%%%%%%%%%%%%

\begin{theorem}[Tree-level Goldstone theorem]\label{thm:coset:goldstone}
At tree level the potential $V$ of Postulate~P1 is $E_{6}$-invariant, and
$\mathcal M$ is an orbit of broken vacua. Hence the Hessian of $V$ vanishes
along the coset directions,
\begin{equation}\label{eq:coset:hess}
  \mathrm{Hess}\,V_{\mathrm{tree}}\big|_{\mathfrak m} = 0,
\end{equation}
and all $32$ clock fields $\phi^{\alpha}$ are exact massless Goldstone bosons.
\end{theorem}

\begin{proof}
The vacuum $X_{0}\in\mathcal M$ minimizes $V$ (\autoref{thm:vacuum:T4}). For any
broken generator $X_{\alpha}\in\mathfrak m$ the curve
$t\mapsto e^{tX_{\alpha}}\!\cdot X_{0}$ lies in the minimizing orbit, so $V$ is
constant along it; differentiating twice gives
$X_{\alpha}{}^{\!\top}\,\mathrm{Hess}\,V(X_{0})\,X_{\beta}=0$ for all
$X_{\alpha},X_{\beta}\in\mathfrak m$. Thus the $32$ coset directions are exact
null eigenvectors of the tree Hessian, i.e.\ massless Goldstone modes
(Goldstone's theorem). The $46$ unbroken directions in $\mathfrak h$ produce the
massless $\Spin(10)\times U(1)$ gauge bosons of \autoref{sec:gauge}.
\end{proof}

\begin{remark}[What lifts the masses, and what stays massless]\label{rem:coset:loop}
The tree degeneracy \eqref{eq:coset:hess} is lifted radiatively. Because the
one-loop Coleman--Weinberg potential is $H$-invariant and $\mathfrak m$ is a
\emph{single} real-irreducible $H$-module (\autoref{thm:coset:irrep}), Schur's
lemma forces a \emph{single common eigenvalue} on all of $\mathfrak m$: one
common radiative mass, with no $\Spin(10)$-singlet subspace to split off. Of the
$32$ clock fields, four are then identified by the soldering map (Postulate~P3,
\autoref{sec:soldering}) with the spacetime frame and are protected by gauge
symmetry (diffeomorphism $+$ local Lorentz, Ward identity
$\nabla_{\mu}(\delta\Gamma/\delta e^{a}{}_{\mu})=0$); they are eaten, not Hessian
zeros. This yields the locked count $32 = 28\;(\text{massive, common mass})
+ 4\;(\text{gauge/soldering})$, with the canonical mass $m^{2}=\mu^{2}$
($[f]=\text{mass}$; the dim-$4$ Hessian eigenvalue is $f^{2}\mu^{2}$). The
spectrum is established in \autoref{thm:gravity:mass}; we do \emph{not} claim a
four-dimensional $H$-invariant Goldstone subspace here.
\end{remark}

\begin{remark}[Topology of the clock-field manifold]\label{rem:coset:topology}
$\mathcal M = \mathrm{EIII}$ is simply connected as a compact Hermitian
symmetric space, so $\pi_{1}(\mathcal M)$ is trivial; in particular
$\pi_{1}(\mathcal M)=\mathbb Z$ is \emph{false} of $\mathrm{EIII}$ itself. The
vortex/string configurations of \autoref{sec:strings} therefore arise from the
\emph{gauged} (U(1)-quotient) configuration space rather than from
$\pi_{1}(\mathcal M)$, and are presented there as a low-energy correspondence.
\end{remark}

\paragraph{Status of the claims.}
\autoref{thm:coset:irrep}, \autoref{thm:coset:metric},
\autoref{lem:coset:positivity}, and \autoref{thm:coset:goldstone} are
\textbf{theorems}, proved from the locked branching \eqref{eq:coset:adjoint} and
the compact real form. The four-direction soldering invoked in
\autoref{rem:coset:nosinglet} and \autoref{rem:coset:loop} is
\textbf{Postulate~P3} (added structure, \autoref{sec:soldering}); the matter and
Yukawa content is \textbf{Postulates~P4, P6} (\autoref{sec:sm}). The coset geometry and its unique positive-definite
metric are derived; $d=4$, signature $(1,3)$, and the Standard-Model
content are added postulates.
% --- end inlined: unit-coset.tex ---
%----------------------------------------------------------------------%
% Part: Spacetime, Gauge Sector, and Mass Spectrum
%----------------------------------------------------------------------%
\part{Spacetime, Gauge Sector, and Mass Spectrum}
% --- begin inlined: unit-spacetime.tex ---
%======================================================================%
\section{Emergent Spacetime, Soldering, and Lorentzian Signature}
\label{sec:spacetime:main}
%======================================================================%

The dynamics of UCFT live on the compact real form of $E_6$, whose vacuum
manifold is the rank-1 primitive-idempotent orbit
$\mathcal M = \EIII = E_6/(\Spin(10)\times\Uone)$, $\dim_{\mathbb R}\mathcal M = 32$
(\autoref{thm:potential:T4}). The induced metric on $\mathcal M$ --- the
kinetic/target-space metric of the coset $\sigma$-model --- is
\emph{positive-definite}: the maximal compact subgroup
$H = \Spin(10)\times\Uone$ acts on the reductive complement
$\mathfrak m = \mathbf{16}_{-3}\oplus\overline{\mathbf{16}}_{+3}$, a real-irreducible
module of complex type whose Killing restriction $K|_{\mathfrak m}$ is the unique
(Schur) $H$-invariant form, of signature $(32,0)$. There is \emph{no} Lorentzian
content in the coset geometry.

This section establishes how a four-dimensional \emph{Lorentzian} spacetime is
nonetheless \emph{soldered} onto this Euclidean arena. The logical separation, which
we maintain throughout, is the following. The positive-definite object is the coset
metric $K|_{\mathfrak m}$ on the $32$-dimensional internal tangent space $\mathfrak m$.
The indefinite object is the cubic norm $N$ on the matter representation
$\mathbf{27}=J_3(\mathbb O)$, whose $2\times2$ octonion-Hermitian sub-block
$H_2(\mathbb O)\cong\mathbb R^{1,9}$ carries a Minkowski quadratic form. These are
two distinct forms on two distinct spaces; no single form is asked to be both
definite and indefinite. The soldering map (Postulate~P3) is the bridge that imports
the indefinite $(1,3)$ structure onto a tangent 4-plane of the world manifold,
leaving $K|_{\mathfrak m}$ untouched.

We classify every nontrivial claim as a \textbf{theorem} (proved from stated
hypotheses), \textbf{postulate} (added structure, P1--P6), or
\textbf{conjecture} (motivated, not controlled). The structural constraint
recorded in \autoref{thm:spacetime:antikilling} is that \emph{$d=4$ and signature
$(1,3)$ are added structure (P3), not consequences of the Killing form.}

%----------------------------------------------------------------------%
\subsection{The Minkowski block inside the Albert algebra}
\label{sec:spacetime:block}
%----------------------------------------------------------------------%

Let $\mathbb O$ be the real division octonions, with conjugation $x\mapsto\bar x$,
norm $n(x)=x\bar x\in\mathbb R_{\ge0}$, and Euclidean inner product
$\langle x,y\rangle=\tfrac12(x\bar y+y\bar x)$ on $\mathbb O\cong\mathbb R^8$. The
space of $2\times2$ octonion-Hermitian matrices is
\begin{equation}\label{eq:spacetime:H2O}
  H_2(\mathbb O)=\Bigl\{\,X=\begin{pmatrix} a & x\\ \bar x & b\end{pmatrix}
  \;:\; a,b\in\mathbb R,\; x\in\mathbb O \,\Bigr\},\qquad
  \dim_{\mathbb R}H_2(\mathbb O)=10,
\end{equation}
with the (well-defined, despite non-associativity) determinant
$\det X = ab - x\bar x = ab - n(x)$.

\begin{theorem}[Minkowski identification of $H_2(\mathbb O)$]
\label{thm:spacetime:minkowski}
In light-cone coordinates $a=X^0+X^9$, $b=X^0-X^9$, and $x=(X^1,\dots,X^8)$ in an
orthonormal octonion basis,
\begin{equation}\label{eq:spacetime:det19}
  \det X = (X^0)^2-(X^9)^2-\sum_{i=1}^{8}(X^i)^2 = -\,\eta_{MN}\,X^M X^N,\qquad
  \eta=\diag(-1,\underbrace{+1,\dots,+1}_{9}),
\end{equation}
so that $\bigl(H_2(\mathbb O),\,-\det\bigr)\cong\mathbb R^{1,9}$.
\end{theorem}
\begin{proof}
$\det X = ab-n(x)=(X^0+X^9)(X^0-X^9)-\sum_i(X^i)^2=(X^0)^2-(X^9)^2-\sum_i(X^i)^2$.
By Sylvester's law of inertia the diagonal entries contribute one $+$ ($X^0$) and one
$-$ ($X^9$), and $-n(x)$ contributes eight $-$'s; hence $-\det$ has one timelike and
nine spacelike directions in the convention $(-,+,\dots,+)$.
\end{proof}

\begin{theorem}[Octonionic exceptional isomorphisms]
\label{thm:spacetime:exc}
There are Lie-group isomorphisms, the four members of the classical family
$\SL(2,\mathbb A)\cong\Spin(n{+}1,1)$ for the normed division algebras
$\mathbb A=\mathbb R,\mathbb C,\mathbb H,\mathbb O$ of real dimension $n=1,2,4,8$:
\begin{equation}\label{eq:spacetime:sl2}
  \SL(2,\mathbb R)\cong\Spin(1,2),\quad
  \SL(2,\mathbb C)\cong\Spin(1,3),\quad
  \SL(2,\mathbb H)\cong\Spin(1,5),\quad
  \SL(2,\mathbb O)\cong\Spin(1,9).
\end{equation}
The defining $10$-dimensional representation of $\Spin(1,9)$ on $\mathbb R^{1,9}$ is
realized on $(H_2(\mathbb O),-\det)$; the two chiral $\mathbf{16}$ spinors are
realized on $\mathbb O^2\cong\mathbb R^{16}$.
\end{theorem}
\begin{proof}
Classical (Sudbery; Kugo--Townsend; Manogue--Schray; Baez). For $\mathbb A=\mathbb O$
the group $\SL(2,\mathbb O)$ is \emph{defined} as $\Spin(1,9)$ with its
$\mathbb R^{1,9}$ action presented by \eqref{eq:spacetime:H2O}--\eqref{eq:spacetime:det19},
non-associativity precluding a literal matrix group.
\end{proof}

\begin{theorem}[The physical Lorentz 4-plane]
\label{thm:spacetime:lorentz4}
The chain of composition-algebra inclusions $\mathbb R\subset\mathbb C\subset\mathbb H\subset\mathbb O$
induces $\Spin(1,1)\subset\Spin(1,3)\subset\Spin(1,5)\subset\Spin(1,9)$. In particular
$\SL(2,\mathbb C)=\Spin(1,3)\subset\SL(2,\mathbb O)=\Spin(1,9)$ via the sub-block
$H_2(\mathbb C)=\{[\,[a,z],[\bar z,b]\,]:a,b\in\mathbb R,\,z\in\mathbb C\}\subset H_2(\mathbb O)$,
$\dim_{\mathbb R}H_2(\mathbb C)=4$, with $\det|_{H_2(\mathbb C)}=ab-|z|^2$ the $(1,3)$
Minkowski form in the convention $(-,+,+,+)$.
\end{theorem}
\begin{proof}
The same light-cone change of variables as \autoref{thm:spacetime:minkowski},
restricted to $\mathbb C$, gives $\det=(X^0)^2-(X^1)^2-(X^2)^2-(X^3)^2$. The
$\SL(2,\mathbb C)$ determinant-covariance is the standard four-dimensional statement.
\end{proof}

The ambient indefinite structure is the octonionic $\mathbb R^{1,9}$; the
\emph{physical} Lorentz signature $(1,3)$ is a soldered 4-plane inside it
(\autoref{sec:spacetime:P3}).

%----------------------------------------------------------------------%
\subsection{Source of indefiniteness: the cubic norm, not the trace form}
\label{sec:spacetime:norm}
%----------------------------------------------------------------------%

The Albert algebra $J_3(\mathbb O)$ of $3\times3$ octonion-Hermitian matrices,
$\dim_{\mathbb R}J_3(\mathbb O)=27$, carries two basic polynomial invariants: the
quadratic trace form $Q(X)=\Tr(X^2)$ (a sum of real squares and positive-definite
octonion norms, hence \emph{positive-definite}) and the cubic norm
$N(X)=\det X$ (a homogeneous cubic). Their invariance groups differ crucially.

\begin{theorem}[Separation of invariance groups]
\label{thm:spacetime:groups}
$F_4=\Aut(J_3(\mathbb O))$ (compact, $\dim 52$) preserves the Jordan product, hence
\emph{both} $Q$ and $N$; since $Q$ is positive-definite, $F_4\subset\SO(27)$. The
reduced structure group $E_{6(-26)}$ ($\dim 78$) preserves $N$ only up to a real
character $\lambda(g)$ and does \emph{not} preserve $Q$: were it to, it would lie in
$\SO(27)$ and be compact, contradicting its non-compactness (maximal compact $F_4$,
$26$ non-compact directions move $Q$).
\end{theorem}
\begin{proof}
Automorphisms fix $\mathbf 1$, commute with $\circ$, hence fix the characteristic
coefficients $\Tr$, $Q=\Tr(X^2)$, $N=\det$; Euclidean signature forces
$F_4\subset\SO(27)$. The reduced structure group is by definition the linear group
fixing $N$ up to scalar; $\dim E_{6(-26)}=78>52=\dim F_4$ with non-compactness gives
the second claim.
\end{proof}

This excludes the incorrect claim that the structure group preserves $Q$ or that
$(Q,N)$ separates orbits: only $N$ is structure-group invariant (up to a character),
and real $E_6$-orbits on the $\mathbf{27}$ are classified by \emph{rank}
$\in\{0,1,2,3\}$ and the value of $N$, via the sharp map $X^\#$
($\rank X\le1\Leftrightarrow X^\#=0$; $\rank X\le2\Leftrightarrow N(X)=0$), cf.
\autoref{thm:potential:L1}.

\begin{theorem}[$N$ is genuinely indefinite, and contains the Minkowski block]
\label{thm:spacetime:Nindef}
$N$ takes strictly positive, strictly negative, and zero values on
$J_3(\mathbb O)\cong\mathbb R^{27}$. Restricting $N$ to the upper-left $2\times2$ block
$Y=[\,[\xi_1,z_3],[\bar z_3,\xi_2]\,]\in H_2(\mathbb O)$ with the third pivot set to
$\xi_3=1$ and $z_1=z_2=0$,
\begin{equation}\label{eq:spacetime:Nrestrict}
  N(X)\big|_{\mathrm{block}}=\xi_1\xi_2-n(z_3)=\det Y,
\end{equation}
so the Minkowski form $-\det=-\eta^{(1,9)}$ of \autoref{thm:spacetime:minkowski} is the
restriction of the cubic norm $N$ to an $H_2(\mathbb O)\cong\mathbb R^{1,9}$ sub-block.
\end{theorem}
\begin{proof}
On $X=\diag(\xi_1,\xi_2,\xi_3)$, $N=\xi_1\xi_2\xi_3$ takes values $+1,-1,0$ at
$(1,1,1),(-1,1,1),(0,1,1)$; an odd-degree form is necessarily indefinite. The
restriction \eqref{eq:spacetime:Nrestrict} is the explicit evaluation of
$N(X)=\xi_1\xi_2\xi_3+2\Re(z_1z_2z_3)-\sum_i\xi_i\,n(z_i)$ at the stated pivot.
\end{proof}

\begin{theorem}[Definite coset metric and indefinite norm are consistent]
\label{thm:spacetime:consistent}
The positive-definite coset metric $K|_{\mathfrak m}$ and the indefinite spacetime
form are mutually consistent because they live on different modules:
\begin{center}\renewcommand{\arraystretch}{1.25}
\begin{tabular}{@{}llll@{}}
\toprule
Object & Lives on & Signature & Invariance \\
\midrule
Coset metric $K|_{\mathfrak m}$ & $\mathfrak m=\mathbf{16}_{-3}\oplus\overline{\mathbf{16}}_{+3}$ ($32$) & $(32,0)$ definite & $H$ (unique, Schur)\\
Trace form $Q=\Tr(X^2)$ & $J_3(\mathbb O)=\mathbb R^{27}$ & $(27,0)$ definite & $F_4$ (\emph{not} $E_6$)\\
Cubic norm $N=\det$ & $J_3(\mathbb O)=\mathbb R^{27}$ & indefinite & $E_6$ up to character\\
Minkowski $-\det$ & $H_2(\mathbb O)\subset J_3(\mathbb O)$ & $(1,9)$ & $\SL(2,\mathbb O)=\Spin(1,9)$\\
Soldered $\eta$ & 4-plane of $T_x\mathcal M^4$ & $(1,3)$, $(-,+,+,+)$ & $\SO(1,3)$ (postulated, P3)\\
\bottomrule
\end{tabular}
\end{center}
\end{theorem}
\begin{proof}
$K|_{\mathfrak m}$ is positive-definite by the compact-real-form choice (compact $E_6$
has negative-definite Killing form; the induced metric on $\mathfrak m$ is its
sign-flip), unique by Schur since $\mathfrak m$ is real-irreducible of complex type with
commutant $\mathbb C$ (\autoref{cor:VacuumOrbit}). The indefinite $N$ lives on the
$\mathbf{27}$, a representation distinct from the adjoint complement $\mathfrak m$;
hence ``$K|_{\mathfrak m}$ Euclidean'' and ``$N$ indefinite'' are statements about
different data on different spaces. The spacetime signature is imported from
$N|_{H_2(\mathbb O)}$ via P3, and is never asked of $K|_{\mathfrak m}$.
\end{proof}

This excludes the incorrect $(4,28)$ Killing/coset signature: Lorentzian structure does
\emph{not} follow from the coset geometry alone.

%----------------------------------------------------------------------%
\subsection{The soldering postulate P3}
\label{sec:spacetime:P3}
%----------------------------------------------------------------------%

We now state P3. It is a \textbf{Postulate} throughout: \emph{added} structure, not a
theorem of the coset geometry.

\paragraph{Theorem-level inputs.}
From the dynamics we have: (T-a) the vacuum manifold
$\mathcal M=\EIII=E_6/(\Spin(10)\times\Uone)$, $\dim_{\mathbb R}=32$, with reductive
decomposition $\mathfrak e_6=\mathfrak h\oplus\mathfrak m$,
$\mathfrak h=\mathfrak{so}(10)\oplus\mathfrak u(1)$,
$\mathfrak m=\mathbf{16}_{-3}\oplus\overline{\mathbf{16}}_{+3}$, and Maurer--Cartan
form $\theta=\theta_{\mathfrak h}+\theta_{\mathfrak m}$ valued in $\mathfrak e_6$;
(T-b) the matter rep $\mathbf{27}=J_3(\mathbb O)$ carrying the indefinite $N$ with its
$H_2(\mathbb O)\cong\mathbb R^{1,9}$ and $H_2(\mathbb C)\cong\mathbb R^{1,3}$ sub-blocks
(\autoref{thm:spacetime:minkowski}--\autoref{thm:spacetime:lorentz4},
\autoref{thm:spacetime:Nindef}); (T-c) the branching
$\mathbf{27}=\mathbf 1_{+4}\oplus\mathbf{10}_{-2}\oplus\mathbf{16}_{+1}$ under $H$.

The $32$ clock fields are the coordinates along $\mathfrak m$ (Goldstone directions of
$\EIII$). Of these, $28$ are physical massive scalars (single common radiative mass by
Schur on the irreducible $\mathfrak m$, \autoref{thm:gravity:mass}) and $4$ are
gauge/soldering modes eaten by world-volume diffeomorphism plus local Lorentz. The $4$
frame modes are \emph{gauge-protected}, not Hessian zeros: $\mathfrak m$ is
irreducible and contains no $\SO(10)$ singlet, so no $4$-dimensional $H$-invariant
subspace of $\mathfrak m$ exists. The old rank-four lemma (\texttt{lem:Rank4}) is
hereby relabeled: $d=4$ is fixed by P3, and its circular $r>4$ branch is deleted.

\begin{postulate}[P3a, soldering / frame identification]\label{post:spacetime:P3a}
There exists a soldering map $\sigma$ identifying four distinguished clock-field frame
directions, written $e^a_\mu$ ($a=0,1,2,3$; $\mu$ a world index), with a tangent
4-plane $\Pi_x\subset T_x\mathcal M^4$ of an emergent four-dimensional world manifold
$\mathcal M^4$. The vierbein is the pull-back of the Maurer--Cartan form along these
four directions,
\begin{equation}\label{eq:spacetime:vierbein}
  e^a_\mu(x)=\bigl\langle\,\tau^a,\,\theta_{\mathfrak m}(\partial_\mu)\,\bigr\rangle,
\end{equation}
where $\{\tau^a\}_{a=0}^{3}$ is the soldering frame --- the image of the
$H_2(\mathbb C)\cong\mathbb R^{1,3}$ basis under the embedding of
\autoref{post:spacetime:P3d}. The vierbein \emph{emerges dynamically}.
\end{postulate}

\begin{postulate}[P3b, induced Minkowski form]\label{post:spacetime:P3b}
The induced metric on $\Pi_x$ is the Minkowski form inherited from the cubic-norm /
$H_2(\mathbb O)\cong\mathbb R^{1,9}$ structure:
\begin{equation}\label{eq:spacetime:metric}
  g_{\mu\nu}(x)=e^a_\mu(x)\,e^b_\nu(x)\,\eta_{ab},\qquad
  \eta_{ab}=\diag(-1,+1,+1,+1),
\end{equation}
with $\eta$ the restriction of $-\det=-N$ to the $H_2(\mathbb C)$ sub-block
(\autoref{thm:spacetime:lorentz4}, \autoref{thm:spacetime:Nindef}).
\end{postulate}

\begin{postulate}[P3c, dimension and signature are added]\label{post:spacetime:P3c}
The numbers $d=4$ and signature $(1,3)=(-,+,+,+)$ are \emph{part of P3}: added
structure, \emph{not} consequences of the Killing form, the coset signature, or any
rank lemma.
\end{postulate}

\begin{postulate}[P3d, vierbein condensate / frame reduction]\label{post:spacetime:P3d}
A vierbein condensate (a VEV of the soldering bilinear) selects the
$H_2(\mathbb C)\cong\mathbb R^{1,3}$ sub-block of $H_2(\mathbb O)\cong\mathbb R^{1,9}$ and
breaks the ambient frame symmetry
\begin{equation}\label{eq:spacetime:framebreak}
  \Spin(1,9)=\SL(2,\mathbb O)\;\longrightarrow\;\Spin(1,3)\times\Spin(6)
  =\SL(2,\mathbb C)\times\SU(4),
\end{equation}
where the orthogonal $\mathbb R^6$ complement of $\mathbb R^{1,3}$ in $\mathbb R^{1,9}$
is acted on by internal $\Spin(6)\cong\SU(4)$. The residual $\Spin(1,3)=\SL(2,\mathbb C)$
is the local Lorentz group.
\end{postulate}

\begin{theorem}[Internal consistency of P3]\label{thm:spacetime:P3consistent}
Given \autoref{post:spacetime:P3a}--\autoref{post:spacetime:P3d}, the metric
$g_{\mu\nu}$ of \eqref{eq:spacetime:metric} is a non-degenerate symmetric rank-2 tensor
of Lorentzian signature $(1,3)$ on $\mathcal M^4$, invariant under the residual local
Lorentz group $\SO(1,3)$; the four frame modes $e^a_\mu$ carry exactly the gauge
degrees of freedom eaten by diffeomorphisms plus local Lorentz, consistent with the
count $32=28\text{ physical}+4\text{ gauge}$.
\end{theorem}
\begin{proof}
Non-degeneracy and signature are inherited from $\eta_{ab}=\diag(-1,+1,+1,+1)$
(\autoref{thm:spacetime:lorentz4}) for generic condensate $\det(e^a_\mu)\ne0$;
$\SO(1,3)$-invariance is $\eta_{ab}=\Lambda^c{}_a\Lambda^d{}_b\eta_{cd}$ for
$\Lambda\in\SO(1,3)$. The gauge count is the Ward identity
$\nabla_\mu(\delta\Gamma/\delta e^a_\mu)=0$: the four frame directions are protected by
diffeomorphism plus local Lorentz, not by a Hessian zero.
\end{proof}

\paragraph{What P3 does not claim.}
P3 does not make $K|_{\mathfrak m}$ Lorentzian (it stays positive-definite); it does not
use a $(4,28)$ Lorentzian split of $\mathfrak m$ (the $4$ are gauge, the $28$ are
Euclidean target coordinates); and it does not build the 4-plane from $\nabla I_2,
\nabla I_3$ singlets (there are none in $\mathfrak m$, and $\nabla Q|_{X_0}=2X_0\notin\mathfrak m$).
The 4-plane is the soldered image of $H_2(\mathbb C)\subset\mathbf{27}$, not a submodule
of $\mathfrak m$.

\begin{theorem}[No Killing-form derivation of dimension or signature]
\label{thm:spacetime:antikilling}
There is no derivation of $d=4$ or signature $(1,3)$ from the Killing form of $E_6$,
from the coset metric $K|_{\mathfrak m}$, or from any rank lemma. The Killing form of
compact $E_6$ restricted to $\mathfrak m$ is positive-definite of signature $(32,0)$,
with zero Lorentzian content. Any $(4,28)$ or $(1,3)$ Killing/coset signature claim is
false. Lorentzian signature enters \emph{only} through P3, rooted in the indefinite
cubic norm $N$ / $H_2(\mathbb O)\cong\mathbb R^{1,9}$.
\end{theorem}
\begin{proof}
Immediate from \autoref{thm:spacetime:consistent}: $K|_{\mathfrak m}$ is definite, and
the indefinite datum lives on the disjoint module $\mathbf{27}$.
\end{proof}

%----------------------------------------------------------------------%
\subsection{Emergent time: Page--Wootters clock and Osterwalder--Schrader}
\label{sec:spacetime:time}
%----------------------------------------------------------------------%

Lorentzian time is recovered by a relational-clock plus reconstruction-theorem
mechanism. We use \emph{no} ``multiply $e^0$ by $i$'' Wick trick.

\begin{postulate}[Timelike clock selection]\label{post:spacetime:clock}
Among the $32$ clock fields, the soldering condensate
(\autoref{post:spacetime:P3d}) selects a distinguished timelike direction whose clock
field we call $\phi^0$: the trace/$X^0$ direction of the soldered
$H_2(\mathbb C)\subset J_3(\mathbb O)$ block (the unique $\eta$-timelike frame leg
$e^0_\mu$). Its three spacelike companions are $e^i_\mu$, $i=1,2,3$. The selection of
\emph{one} timelike direction is the $(1,3)$ content of P3.
\end{postulate}

Promote $\phi^0$ to a \textbf{Page--Wootters clock}: enlarge the kinematical Hilbert
space to $\mathcal H=\mathcal H_{\mathrm{clock}}\otimes\mathcal H_{\mathrm{sys}}$ and
impose the induced-gravity Hamiltonian constraint
\begin{equation}\label{eq:spacetime:constraint}
  \widehat H_{\mathrm{phys}}\,|\Psi\rangle\!\rangle=0,\qquad
  \widehat H_{\mathrm{phys}}=\widehat H_{\mathrm{clock}}+\widehat H_{\mathrm{sys}}
  \;(+\text{ interaction}),
\end{equation}
on physical states $|\Psi\rangle\!\rangle$. Conditioning on clock readings $\phi^0=\tau$
defines the relational state $|\psi(\tau)\rangle_{\mathrm{sys}}=\langle\phi^0{=}\tau|\Psi\rangle\!\rangle$.

\begin{theorem}[Page--Wootters relational evolution]\label{thm:spacetime:PW}
With a good clock ($\widehat H_{\mathrm{clock}}=\widehat p_{\phi^0}$ canonically
conjugate to $\widehat\phi^0$), the conditioned state satisfies the Schr\"odinger
equation in the relational time $\tau$,
\begin{equation}\label{eq:spacetime:schrodinger}
  i\,\partial_\tau\,|\psi(\tau)\rangle_{\mathrm{sys}}
  =\widehat H_{\mathrm{sys}}\,|\psi(\tau)\rangle_{\mathrm{sys}},
\end{equation}
and the unitary group $U(\tau)=e^{-i\widehat H_{\mathrm{sys}}\tau}$ is recovered on
physical states in the semiclassical/background regime where the soldered $g_{\mu\nu}$
admits a timelike Killing vector and $\widehat H_{\mathrm{sys}}=\int T_{00}$ along it.
\end{theorem}
\begin{proof}
Standard Page--Wootters reduction: under the constraint
\eqref{eq:spacetime:constraint} with a conjugate clock, the conditioned state obeys
\eqref{eq:spacetime:schrodinger}. The identification $t\leftrightarrow x^0$ is
$\widehat H=\int T_{00}$ along the timelike Killing direction. Global hyperbolicity is
flagged as an assumption.
\end{proof}

This demotes the original Axiom II (an externally postulated $U(t)=e^{-iHt}$) to a
theorem: time is emergent.

\paragraph{Euclidean to Lorentzian via Osterwalder--Schrader.}
The compact-form $\sigma$-model with positive-definite $K|_{\mathfrak m}$ defines a
Euclidean functional measure. Lorentzian QFT is obtained by Osterwalder--Schrader (OS)
reconstruction along the selected time $\phi^0$, \emph{not} by rotating
$e^0\to i\,e^0$. Let $\tau=\phi^0$ be the OS time, and assume the Euclidean Schwinger
functions satisfy the OS axioms, in particular \emph{reflection positivity} under
$\theta:\tau\mapsto-\tau$ about a $\phi^0$-slice,
\begin{equation}\label{eq:spacetime:reflpos}
  \sum_{i,j}\bar c_i\,c_j\;S\bigl(\theta F_i\cdot F_j\bigr)\ge0,
  \qquad F_i \text{ supported in }\{\tau>0\}.
\end{equation}

\begin{theorem}[OS reconstruction $\Rightarrow$ Lorentzian Wightman theory]
\label{thm:spacetime:OS}
Given the OS axioms (Euclidean invariance, reflection positivity
\eqref{eq:spacetime:reflpos}, symmetry, clustering), the OS reconstruction theorem
produces a Hilbert space $\mathcal H$, a positive self-adjoint Hamiltonian
$\widehat H\ge0$ generating $\tau$-translations as the contraction semigroup
$e^{-\widehat H\tau}$, and --- by analytic continuation of the Schwinger functions in
the clock-time argument $\tau\to it$ --- Lorentzian Wightman functions on
$(\mathcal M^4,\eta^{(1,3)})$ with unitary evolution $e^{-i\widehat H t}$ satisfying the
spectrum condition and Lorentz covariance under the reconstructed $\SO(1,3)$.
\end{theorem}
\begin{proof}
Osterwalder--Schrader (1973, 1975): reflection positivity defines a physical inner
product on $\{F:\tau>0\}/(\text{null})$; $\widehat H$ is the positive self-adjoint
generator of $e^{-\widehat H\tau}$; the Schwinger functions, analytic in a tube,
continue to Wightman distributions.
\end{proof}

\begin{indicative}[Reflection positivity of the UCFT $\sigma$-model]
\label{ind:spacetime:reflpos}
The OS theorem (\autoref{thm:spacetime:OS}) is rigorous given its hypotheses. That the
UCFT $\sigma$-model \emph{satisfies} reflection positivity \eqref{eq:spacetime:reflpos}
along $\phi^0$ is the precise property to be checked --- it is the rigorous
replacement for the Wick trick, and is tagged \textbf{Indicative}.
\end{indicative}

The analytic continuation $\tau\to it$ acts on the Schwinger distributions in the
clock-time argument (justified by reflection positivity and the resulting analyticity
tube), \emph{not} on a frame field. The Lorentzian metric $\eta^{(1,3)}$ is the
soldered form of P3, already Lorentzian before any continuation; OS supplies the
unitary, positive-energy dynamics, with $\phi^0$ serving as the OS reflection time. The
two mechanisms dovetail: \autoref{thm:spacetime:PW} gives the relational unitary
$U(\tau)$ on physical states, and \autoref{thm:spacetime:OS} gives the positive-energy,
Lorentz-covariant Wightman theory.

%----------------------------------------------------------------------%
\subsection{Master statement}
\label{sec:spacetime:master}
%----------------------------------------------------------------------%

\begin{postulate}[P3, hardened soldering --- master statement]\label{post:spacetime:P3master}
Given (i) the compact-$E_6$ dynamics with vacuum $\EIII=E_6/(\Spin(10)\times\Uone)$ and
positive-definite coset metric $K|_{\mathfrak m}$ (Theorem-level inputs), and (ii) the
matter rep $\mathbf{27}=J_3(\mathbb O)$ with its indefinite cubic norm $N$ containing the
$H_2(\mathbb O)\cong\mathbb R^{1,9}$ ($\SL(2,\mathbb O)=\Spin(1,9)$) and
$H_2(\mathbb C)\cong\mathbb R^{1,3}$ ($\SL(2,\mathbb C)=\Spin(1,3)$) sub-blocks
(\autoref{thm:spacetime:minkowski}--\autoref{thm:spacetime:lorentz4},
\autoref{thm:spacetime:Nindef}), we postulate a soldering map $\sigma$
(\autoref{post:spacetime:P3a}) identifying four clock-field frame directions $e^a_\mu$
(pull-backs of the Maurer--Cartan form) with a tangent 4-plane of an emergent
four-dimensional world manifold, carrying $g_{\mu\nu}=e^a_\mu e^b_\nu\eta_{ab}$,
$\eta=\diag(-1,+1,+1,+1)$ (\autoref{post:spacetime:P3b}), with $d=4$ and signature
$(1,3)$ as added structure (\autoref{post:spacetime:P3c}), realized by a vierbein
condensate breaking $\Spin(1,9)\to\Spin(1,3)\times\Spin(6)$
(\autoref{post:spacetime:P3d}). The timelike clock $\phi^0$ is a Page--Wootters clock
and Lorentzian unitary dynamics follows from Osterwalder--Schrader reconstruction along
$\phi^0$ (\autoref{thm:spacetime:OS}), not by Wick rotation of the frame.
\end{postulate}

The vierbein $e^a_\mu$ of \eqref{eq:spacetime:vierbein}, carrying simultaneously a
world index $\mu$ and a soldering index $a$, is the fundamental gravitational field; its
$\sigma$-model dynamics and the induced Einstein--Hilbert term are developed in
\autoref{sec:gravity:main}. The low-energy field content on the soldered spacetime ---
the $46=\dim\SO(10)+1$ massless gauge bosons, the chiral $\mathbf{16}_{+1}$ fermions,
and the $28$ massive coset scalars --- transforms under the unbroken
$H=\Spin(10)\times\Uone$ and is catalogued in \autoref{sec:sm:main}.
% --- end inlined: unit-spacetime.tex ---
% --- begin inlined: unit-gauge.tex ---
%======================================================================%
\section{The composite \texorpdfstring{$\Spin(10)\times\Uone$}{Spin(10) x U(1)} gauge sector}
\label{sec:gauge:main}
%======================================================================%

Having fixed the vacuum manifold as the rank-one primitive-idempotent
orbit
\[
  \mathcal M \;=\; \ESix/\bigl(\Spin(10)\times\Uone\bigr) \;=\; \mathrm{EIII},
  \qquad \dim_{\mathbb R}\mathcal M = 32,
\]
(\autoref{thm:potential:T4}), we now show that the \emph{same} geometric
data --- the Maurer--Cartan form of $\ESix$ pulled back along a coset map
--- generate a composite $H=\Spin(10)\times\Uone$ Yang--Mills theory.
The gauge bosons are not fundamental: they are the unbroken-subalgebra
components of the canonical connection on the homogeneous bundle
$\ESix\to\mathcal M$. We establish (i) the composite connection and the
$\mathbf{45}_{0}\oplus\mathbf{1}_{0}$ gauge content, (ii) the induced
Yang--Mills term with the \emph{corrected} coefficient computed from the
massive multiplet $\mathfrak m=\mathbf{16}_{-3}\oplus\overline{\mathbf{16}}_{+3}$,
(iii) that $H$ remains \emph{unbroken} at the vacuum so that the
$\dim\mathbf{45}+\dim\mathbf 1 = 46$ gauge bosons stay massless, and
(iv) Green--Schwarz cancellation of the residual mixed anomaly. Throughout
we work on the \emph{compact} real form of $\ESix$, so that $H$ is compact,
the coset metric is positive-definite, and the gauge sector is unitary
(\autoref{sec:foundations:realform}). Signature is $(-,+,+,+)$ and
$[f]=\text{mass}$.

%----------------------------------------------------------------------%
\subsection{The composite connection from the Maurer--Cartan form}
\label{sec:gauge:MC}
%----------------------------------------------------------------------%

The Lie algebra splits, $H$-equivariantly, as the reductive sum
\begin{equation}\label{eq:gauge:split}
  \mathfrak e_6
  \;=\;
  \underbrace{\mathfrak{so}_{10}\oplus\mathfrak u_{1}}_{\mathfrak h}
  \;\oplus\;
  \underbrace{\mathfrak m}_{\text{coset}},
  \qquad
  [\mathfrak h,\mathfrak m]\subseteq\mathfrak m,
\end{equation}
with the branching of the adjoint under $H=\Spin(10)\times\Uone$ fixed by
the involution datum
\begin{equation}\label{eq:gauge:adjoint}
  \mathbf{78}
  \;=\;
  \underbrace{\mathbf{45}_{0}\oplus\mathbf 1_{0}}_{\mathfrak h}
  \;\oplus\;
  \underbrace{\mathbf{16}_{-3}\oplus\overline{\mathbf{16}}_{+3}}_{\mathfrak m}.
\end{equation}
Here the subscripts are the $\Uone$ charges; $\mathfrak m$ is real-irreducible
of complex type and contains \emph{no} $\Spin(10)$ singlet
(\autoref{sec:foundations:coset}). The $46$ generators of $\mathfrak h$ are
exactly the composite gauge degrees of freedom.

Let $g(x)\in\ESix$ be a smooth local section of the principal $H$-bundle,
representing the coset map $\pi_-(x)\in\mathcal M$. The Maurer--Cartan
form decomposes along \eqref{eq:gauge:split},
\begin{equation}\label{eq:gauge:MCdecomp}
  g^{-1}\partial_\mu g
  \;=\;
  A_\mu(x) \;+\; e_\mu(x),
  \qquad
  A_\mu\in\mathfrak{so}_{10}\oplus\mathfrak u_{1},
  \quad
  e_\mu\in\mathfrak m,
\end{equation}
where $A_\mu$ is the $\mathfrak h$-part (the canonical connection) and
$e_\mu$ is the $\mathfrak m$-part (the soldering one-form of the clock
fields, used in \autoref{sec:soldering:main}).

\begin{theorem}[Composite connection; \textsc{theorem}]\label{thm:gauge:composite}
Under a change of local section $g\mapsto g\,h^{-1}(x)$ with
$h(x)\in H$, the $\mathfrak h$-component of \eqref{eq:gauge:MCdecomp}
transforms as a genuine $H$ gauge potential,
\begin{equation}\label{eq:gauge:transform}
  \mathcal A_\mu \;:=\; A_\mu
  \;\longmapsto\;
  h\,\mathcal A_\mu\,h^{-1} - (\partial_\mu h)\,h^{-1},
\end{equation}
while $e_\mu\mapsto h\,e_\mu\,h^{-1}$ transforms homogeneously (in the
$\mathbf{16}_{-3}\oplus\overline{\mathbf{16}}_{+3}$). Its curvature
\begin{equation}\label{eq:gauge:F}
  \mathcal F_{\mu\nu}
  =\partial_\mu\mathcal A_\nu-\partial_\nu\mathcal A_\mu
   +[\mathcal A_\mu,\mathcal A_\nu]
  = -\,\mathrm P_{\mathfrak h}\!\left(g^{-1}[\partial_\mu,\partial_\nu]g\right)
    + [e_\mu,e_\nu]_{\mathfrak h}
\end{equation}
is the $\mathfrak h$-valued field strength of a $\Spin(10)\times\Uone$
gauge theory.
\end{theorem}

\begin{proof}
Equation \eqref{eq:gauge:transform} is the standard transformation of the
canonical (reductive) connection on a homogeneous bundle: writing
$g\mapsto gh^{-1}$ gives
$g^{-1}\partial g\mapsto h(g^{-1}\partial g)h^{-1}-(\partial h)h^{-1}$,
and projecting onto $\mathfrak h$ (which is $H$-stable by
\eqref{eq:gauge:split}) yields \eqref{eq:gauge:transform}; the
$\mathfrak m$-projection transforms homogeneously since
$[\mathfrak h,\mathfrak m]\subseteq\mathfrak m$. The Maurer--Cartan
equation $d(g^{-1}dg)+\tfrac12[g^{-1}dg,g^{-1}dg]=0$, projected onto
$\mathfrak h$, gives \eqref{eq:gauge:F} (the Gauss--Codazzi relation:
the $\mathfrak h$-curvature receives the soldering contribution
$[e_\mu,e_\nu]_{\mathfrak h}$). Both are intrinsic to the reductive split.
\end{proof}

The composite gauge content is therefore precisely $\mathbf{45}_{0}\oplus\mathbf 1_{0}$:
the $\mathbf{45}$ adjoint of $\Spin(10)$ and the $\Uone$ singlet,
$46$ vector fields in all.

%----------------------------------------------------------------------%
\subsection{Induced Yang--Mills term --- corrected coefficient}
\label{sec:gauge:YM}
%----------------------------------------------------------------------%

Fluctuating the massive clock fields around the vacuum and integrating
them out at one loop produces the functional determinant
$[\det(-D^{2}+m^{2})]^{-1/2}$ with $D_\mu=\partial_\mu+\mathcal A_\mu$ and
$m^{2}=f^{2}\mu^{2}$ (\autoref{sec:gravity:spectrum}). The Seeley--DeWitt
$a_4$ coefficient supplies a local Yang--Mills counterterm,
\begin{equation}\label{eq:gauge:YM}
  S_{\text{YM}}
  = -\frac{1}{4g^{2}}\int_{X}\Tr\!\bigl(\mathcal F_{\mu\nu}\mathcal F^{\mu\nu}\bigr)\,d^{4}x,
\end{equation}
whose logarithmic coefficient is fixed by the total Dynkin index of the
fields running in the loop.

\begin{remark}[The corrected loop multiplet; \textsc{theorem}]
\label{rem:gauge:correctedindex}
The integrated-out matter is the \emph{massive coset multiplet}
$\mathfrak m=\mathbf{16}_{-3}\oplus\overline{\mathbf{16}}_{+3}$
($\dim 32$), the chiral spinor pair of $\Spin(10)$, \emph{not} a
spurious $\mathbf{10}\oplus\overline{\mathbf{10}}\oplus\mathbf 4$.
With the $\Spin(10)$ Dynkin indices $T(\mathbf{16})=2$ for each spinor
(and the $\Uone$ singlet $\mathbf 1$ not contributing to the
$\Spin(10)$ trace), the total index is
\begin{equation}\label{eq:gauge:totalindex}
  \sum_{R\subset\mathfrak m}\mathrm{mult}(R)\,T(R)
  \;=\; T(\mathbf{16})+T(\overline{\mathbf{16}})
  \;=\; 2+2 \;=\; 4,
\end{equation}
the value used everywhere. The previously quoted $\tfrac52$, obtained
from the (forbidden) decomposition $\mathbf{10}\oplus\overline{\mathbf{10}}\oplus\mathbf 4$
with indices $1,1,\tfrac12$, is wrong: that multiplet is not the coset
tangent (which carries no $\Spin(10)$ singlet, \autoref{sec:foundations:coset}).
\end{remark}

\begin{theorem}[Induced gauge coupling; \textsc{theorem}]\label{thm:gauge:coupling}
Integrating out $\mathfrak m=\mathbf{16}_{-3}\oplus\overline{\mathbf{16}}_{+3}$
at one loop induces \eqref{eq:gauge:YM} with
\begin{equation}\label{eq:gauge:gm2}
  g^{-2}
  \;=\;
  \frac{4}{1}\,\frac{f^{2}}{48\pi^{2}}\,
  \log\!\frac{\Lambda^{2}}{\mu^{2}}
  \;+\;\mathcal O(\Lambda^{0})
  \;=\;
  \frac{f^{2}}{12\pi^{2}}\,\log\!\frac{\Lambda^{2}}{\mu^{2}}
  \;+\;\mathcal O(\Lambda^{0}),
\end{equation}
with $\Lambda$ the UV cutoff and $\mu$ the breaking scale. In particular
$g^{-2}>0$ for $\Lambda>\mu$, so the induced gauge sector is healthy.
\end{theorem}

\begin{proof}[Proof sketch]
The proper-time representation of $\Gamma_{1L}=\tfrac12\Tr\log(-D^2+m^2)$,
expanded with the heat-kernel coefficient
$a_4(R)=\tfrac{1}{12}\int_X\Tr_R(\mathcal F_{\mu\nu}\mathcal F^{\mu\nu})$
and the cutoff $s\ge\Lambda^{-2}$, gives a logarithmically divergent
$\Tr(\mathcal F\mathcal F)$ coefficient proportional to the total index
\eqref{eq:gauge:totalindex}. Collecting the factors $\tfrac12\cdot\tfrac1{12}\cdot\tfrac1{(4\pi)^2}$
and the index $4$ reproduces \eqref{eq:gauge:gm2}; the full computation,
including normalization of $\Tr(\mathcal F\mathcal F)$ to the
$\Spin(10)$ vector, is deferred to \autoref{app:ym:main}.
\end{proof}

The renormalization-group running of $g^{2}$ --- gauge \emph{asymptotic
freedom}, with $\beta_{g^2}=2g^2+(b_0/48\pi^2)g^4$, $b_0=26$ for one
generation and $b_0=12$ for three (\autoref{thm:sm:b0}) --- follows
from this same induced action and is treated alongside the gravitational
flow; there is no interacting gauge fixed point.

%----------------------------------------------------------------------%
\subsection{$H$ is unbroken at the vacuum}
\label{sec:gauge:unbroken}
%----------------------------------------------------------------------%

\begin{theorem}[Unbroken gauge group; \textsc{theorem}]\label{thm:gauge:unbroken}
At the vacuum $\langle\Phi\rangle\in\mathcal M$ the stabilizer is exactly
$H=\Spin(10)\times\Uone$. Consequently none of the $46$ gauge bosons of
$\mathbf{45}_{0}\oplus\mathbf 1_{0}$ acquires a mass: the gauge sector
remains massless and $H$ is unbroken.
\end{theorem}

\begin{proof}
By construction $\mathcal M=\ESix/H$ is the orbit of a rank-one primitive
idempotent, whose isotropy group is $H$ (\autoref{thm:potential:T4}).
A vector boson associated with a generator $T\in\mathfrak e_6$ becomes
massive only if $T$ moves the vacuum, i.e.\ if $T\in\mathfrak m$; the
would-be Goldstone--Higgs mechanism eats a coset scalar for each broken
generator. The unbroken generators $\mathfrak h$ annihilate
$\langle\Phi\rangle$ and the corresponding $46$ vectors stay massless.
There is no Higgsing of $\mathfrak h$: the $32$ coset Goldstones
(``clock fields'') are not in $\mathfrak h$, and the scalar potential is
$H$-invariant, so the one-loop scalar mass matrix is $H$-invariant and
commutes with the $\mathfrak h$-action. Schur's lemma on the irreducible
$\mathfrak m$ gives a single common scalar mass with no $\mathfrak h$
admixture (\autoref{sec:gravity:spectrum}), preserving $H$.
\end{proof}

\begin{remark}[Ward identity; \textsc{theorem}]\label{rem:gauge:ward}
Masslessness is protected to all orders by the $H$ Ward identity. With
the BRST-gauge-fixed effective action $\Gamma$, current conservation
$D_\mu\langle J^{\mu}_{a}\rangle=0$ for $a$ ranging over $\mathfrak h$
forces the transverse vacuum polarization $\Pi_{ab}(p)$ to vanish at
$p^2=0$; no $H$-invariant scalar condensate in $\mathfrak m$ exists to
generate a mass term (there is no $\Spin(10)$ singlet in $\mathfrak m$),
so the $46$ gauge bosons remain exactly massless. This is consistent with
the count $32 = 28\ \text{massive scalars} + 4\ \text{gauge/soldering}$:
the $4$ frame modes are eaten by diffeomorphism $+$ local Lorentz, not by
$H$ (\autoref{sec:soldering:main}).
\end{remark}

%----------------------------------------------------------------------%
\subsection{Green--Schwarz anomaly cancellation}
\label{sec:gauge:gs}
%----------------------------------------------------------------------%

The chiral matter is one family in the SO(10) generation, the
$\mathbf{16}_{+1}$ together with the remaining components of the
fundamental,
\begin{equation}\label{eq:gauge:matter}
  \mathbf{27}
  \;=\;
  \mathbf 1_{+4}\oplus\mathbf{10}_{-2}\oplus\mathbf{16}_{+1}
  \qquad(\text{under }H),
\end{equation}
with $\Uone$ charges as shown. The cubic norm $N$ on $J_3(\mathbb O)$
furnishes a closed, $\ESix$-invariant $3$-form $\Omega_3$ on $\mathcal M$;
let $B$ be a two-form on $X$ with $dB=\pi_-^{*}\Omega_3$. The
Chern--Simons transgression of \eqref{eq:gauge:MCdecomp} produces the
Green--Schwarz coupling
\begin{equation}\label{eq:gauge:GS}
  S_{GS}
  \;=\;
  \int_{X} B\wedge\frac{\Tr(\mathcal F\wedge\mathcal F)}{8\pi^{2}},
  \qquad
  dB=\frac{1}{2\pi}F_{1},
\end{equation}
where $F_1$ is the $\Uone$ field strength.

\begin{theorem}[Anomaly cancellation; \textsc{theorem}]\label{thm:gauge:anomaly}
The combined gauge variation
$\delta_H\bigl(\Gamma_{1L}+S_{\text{YM}}+S_{GS}\bigr)$ vanishes. The only
nonvanishing one-loop gauge anomaly is the mixed $\Uone$--$\Spin(10)^{2}$
anomaly, with coefficient $-1$, and it is cancelled by the inflow from
$S_{GS}$; the cubic $\Uone^{3}$ anomaly vanishes identically.
\end{theorem}

\begin{proof}[Proof sketch]
For a chiral Weyl fermion in rep $R$ of $\Spin(10)$ with $\Uone$ charge
$q$, the relevant six-form is
$I_6(R,q)=q\,A_R\,\frac{F_1}{2\pi}\wedge\frac{\Tr(\mathcal F\wedge\mathcal F)}{8\pi^2}$,
with $A_{\mathbf{16}}=A_{\mathbf{10}}=1$, $A_{\mathbf 1}=0$. Summing over
\eqref{eq:gauge:matter},
\begin{equation}\label{eq:gauge:mixed}
  (+1)A_{\mathbf{16}}+(-2)A_{\mathbf{10}}+(+4)A_{\mathbf 1}
  = 1-2+0 = -1,
\end{equation}
giving $I_6^{\text{tot}}=-\frac{F_1}{2\pi}\wedge\frac{\Tr(\mathcal F\wedge\mathcal F)}{8\pi^2}$.
The cubic abelian anomaly cancels,
\begin{equation}\label{eq:gauge:cubic}
  16\cdot(+1)^{3}+10\cdot(-2)^{3}+1\cdot(+4)^{3}
  = 16 - 80 + 64 = 0.
\end{equation}
Under an abelian gauge shift $B\mapsto B+\Lambda_1/2\pi$ the variation of
\eqref{eq:gauge:GS} is
$\delta_\Lambda S_{GS}=\int_X\frac{\Lambda_1}{2\pi}\wedge\frac{\Tr(\mathcal F\wedge\mathcal F)}{8\pi^2}=-\int_X I_6^{\text{tot}}=-\delta_\Lambda\Gamma_{1L}$,
exactly cancelling the mixed anomaly. The descent equations and the full
arithmetic appear in \autoref{app:anomaly:main}.
\end{proof}

\begin{theorem}[Emergent gauge sector; \textsc{theorem}]\label{thm:gauge:sector}
The functional $S_{\text{tot}}=S_{\text{eff}}+S_{\text{YM}}+S_{GS}$
describes a local, unitary, anomaly-free $\Spin(10)\times\Uone$
Yang--Mills theory coupled to the clock-field $\sigma$-model, with the
dynamically induced coupling \eqref{eq:gauge:gm2}. Locality and unitarity
follow from those of the scalar $\sigma$-model (\autoref{thm:foundations:qft})
together with compactness of $H$ and standard BRST gauge fixing; anomaly
freedom is \autoref{thm:gauge:anomaly}.
\end{theorem}

\paragraph{Status.} The composite connection
(\autoref{thm:gauge:composite}), the unbroken-$H$ statement
(\autoref{thm:gauge:unbroken}), the corrected total index
(\autoref{rem:gauge:correctedindex}), and the anomaly cancellation
(\autoref{thm:gauge:anomaly}) are \emph{theorems} given the vacuum
$\mathcal M=\mathrm{EIII}$ and the compact real form. The chiral matter
assignment \eqref{eq:gauge:matter} with no mirror partners is
\emph{postulate}~P4. The detailed Yang--Mills coefficient is in
\autoref{app:ym:main} and the full anomaly arithmetic in
\autoref{app:anomaly:main}.
% --- end inlined: unit-gauge.tex ---
% --- begin inlined: unit-masses.tex ---
%======================================================================%
\section{Radiative Masses and the Physical Spectrum}
\label{sec:masses:main}
%======================================================================%

The vacuum manifold of UCFT is the rank-one primitive-idempotent orbit
\[
  \mathcal{M} \;=\; E_{6}/\bigl(\mathrm{Spin}(10)\times U(1)\bigr) \;=\; \mathrm{EIII},
  \qquad \dim_{\mathbb R}\mathcal{M} = 32,
\]
established as the global minimum of the $E_{6}$-invariant potential in
\autoref{thm:ExplicitBreaking} (vacuum-selection theorem; cf.\ the locked
foundation, Hardened Result~A). The $32$ coset Goldstones --- the
\emph{clock fields} --- span the tangent module
\begin{equation}\label{eq:masses:tangent}
  \mathfrak m \;=\; T_{X_{0}}\mathcal{M}
  \;=\; \mathbf{16}_{-3}\,\oplus\,\overline{\mathbf{16}}_{+3},
  \qquad \dim_{\mathbb R}\mathfrak m = 32,
\end{equation}
under the unbroken $H=\mathrm{Spin}(10)\times U(1)$ (\autoref{thm:InvariantMetric}).
This section determines their physical mass spectrum. The central structural
fact, used throughout, is that \eqref{eq:masses:tangent} is a \emph{single}
real-irreducible $H$-module: the $U(1)$ charge forbids the symmetric
$\mathbf{16}\times\mathbf{16}$ and $\overline{\mathbf{16}}\times\overline{\mathbf{16}}$
invariants, leaving only the charge-zero cross pairing, so the commutant of
$H$ on $\mathfrak m$ is $\mathbb C$ (real-irreducible of complex type) and
$\mathfrak m$ contains \emph{no} $\mathrm{Spin}(10)$ singlet.

We work on the compact real form (equivalently $E_{6(-14)}$, whose maximal
compact subgroup is precisely $H$), so the coset metric $K|_{\mathfrak m}$ is
positive-definite and the gauge sector is unitary. Throughout, $[f]=\text{mass}$
is the decay constant, $\mu$ the breaking/mass scale, and $\Lambda$ the UV
cutoff; the signature is $(-,+,+,+)$.

%----------------------------------------------------------------------%
\subsection{Tree level: thirty-two exact Goldstones}
\label{sec:masses:tree}
%----------------------------------------------------------------------%

\begin{theorem}[Tree-level Goldstone degeneracy]\label{thm:masses:tree}
At any vacuum $X_{0}\in\mathcal{M}$ the Hessian of the $E_{6}$-invariant
tree potential $V$ vanishes identically on the coset tangent space,
\begin{equation}\label{eq:masses:treehess}
  \operatorname{Hess}_{X_{0}}V\big|_{\mathfrak m} \;=\; 0,
\end{equation}
so that all $32$ clock fields are exact Goldstone bosons.
\end{theorem}

\begin{proof}
$V$ is constant along the $E_{6}$-orbit $\mathcal{M}$ of the minimizer
(\autoref{thm:ExplicitBreaking}), so its directional derivatives along the
broken generators, which span $\mathfrak m$, vanish to all orders in the
flat directions; \eqref{eq:masses:treehess} is Goldstone's theorem applied to
the spontaneously broken $E_{6}\to H$. The unbroken directions
$\mathfrak h=\mathfrak{so}_{10}\oplus\mathfrak u_{1}$ produce no scalars.
\end{proof}

The degeneracy \eqref{eq:masses:treehess} is exact only at tree level; the
physical spectrum is fixed by the one-loop effective potential.

%----------------------------------------------------------------------%
\subsection{One loop: a single common radiative mass}
\label{sec:masses:oneloop}
%----------------------------------------------------------------------%

The Coleman--Weinberg effective potential $V_{\mathrm{eff}}$ is built from the
$H$-invariant fluctuation operator about $X_{0}$ and is therefore itself
$H$-invariant. Its Hessian on $\mathfrak m$ is consequently an
$H$-intertwiner. We exploit irreducibility \emph{positively}: it forces a
single eigenvalue, rather than splitting off any singlet subspace.

\begin{theorem}[Common radiative mass via Schur]\label{thm:masses:schur}
The one-loop Hessian on the coset is a multiple of the identity,
\begin{equation}\label{eq:masses:schur}
  \operatorname{Hess}_{X_{0}}V_{\mathrm{eff}}\big|_{\mathfrak m}
  \;=\; m_{\mathfrak m}^{2}\,\mathbf 1_{32},
  \qquad m_{\mathfrak m}^{2}\in\mathbb R_{\ge 0},
\end{equation}
i.e.\ there is one common radiative mass on all $32$ coset directions.
\end{theorem}

\begin{proof}
$\operatorname{Hess}_{X_{0}}V_{\mathrm{eff}}|_{\mathfrak m}$ is a symmetric
$H$-equivariant endomorphism of $\mathfrak m$. By \eqref{eq:masses:tangent},
$\mathfrak m$ is a single real-irreducible $H$-module of complex type, with
commutant $\operatorname{End}_{H}(\mathfrak m)\cong\mathbb C$. A symmetric (real,
self-adjoint with respect to the positive-definite $K|_{\mathfrak m}$)
intertwiner must lie in the real part of the commutant, hence is a real scalar
multiple of the identity; this is Schur's lemma for the complex-type case
restricted to self-adjoint operators. Thus the Hessian equals
$m_{\mathfrak m}^{2}\mathbf 1_{32}$. Stability of the vacuum
(\autoref{thm:ExplicitBreaking}) and positivity of $K|_{\mathfrak m}$ give
$m_{\mathfrak m}^{2}\ge 0$.
\end{proof}

\begin{remark}[No internal singlet subspace]\label{rmk:masses:nosinglet}
Because $\mathfrak m$ carries \emph{no} $\mathrm{Spin}(10)$ singlet, there is no
$H$-invariant four-dimensional subspace of $\mathfrak m$ on which the Hessian
could independently vanish, and in particular no subspace spanned by gradients
of independent invariants $\nabla I_{2},\nabla I_{3}$. Any decomposition of
$\mathfrak m$ as $2\times\mathbf{10}\oplus\mathbf 4$, $\mathbf{10}\oplus
\overline{\mathbf{10}}\oplus\mathbf 4$, or with any singlet factor is
representation-theoretically false. The earlier ``$4\oplus 28$'' Hessian
splitting of $\mathfrak m$ into an internally massless quartet and a massive
$28$-plet is therefore \emph{not} the mechanism; the four light directions are
gauge artefacts, treated next.
\end{remark}

%----------------------------------------------------------------------%
\subsection{The four frame directions are gauge, not Hessian zeros}
\label{sec:masses:frame}
%----------------------------------------------------------------------%

Four of the $32$ clock fields are the soldering/frame directions identified by
Postulate~P3 (\autoref{thm:Lorentz}): the soldering map $\sigma$ aligns four
clock-field frame directions $e^{a}{}_{\mu}$ ($a=0,1,2,3$) with a tangent
$4$-plane of the emergent spacetime $M^{4}$, providing the composite vierbein
$e^{a}{}_{\mu}(x)=\langle\tau^{a},\theta_{\mathfrak m}(\partial_{\mu})\rangle$.
These four directions are not lifted by the Hessian; they are removed by
\emph{gauge} symmetry.

\begin{postulate}[Frame gauge protection, P3]\label{post:masses:frame}
The four soldered frame directions are gauge degrees of freedom of the emergent
geometry: they are the modes eaten under world-volume diffeomorphisms and local
Lorentz transformations $\mathrm{SO}(1,3)$. They satisfy the Ward identity
\begin{equation}\label{eq:masses:ward}
  \nabla_{\mu}\!\left(\frac{\delta\Gamma}{\delta e^{a}{}_{\mu}}\right) \;=\; 0,
\end{equation}
which removes them from the physical spectrum (cf.\ \autoref{thm:Lorentz}).
\end{postulate}

\begin{theorem}[Physical count]\label{thm:masses:count}
Of the $32$ real clock fields, four are eaten as gauge/soldering modes under
diffeomorphism $+$ local Lorentz, and the remaining $28$ are physical massive
scalars carrying the single common mass of \autoref{thm:masses:schur}:
\begin{equation}\label{eq:masses:count}
  \underbrace{32}_{\dim_{\mathbb R}\mathfrak m}
  \;=\;
  \underbrace{28}_{\text{physical massive scalars}}
  \;+\;
  \underbrace{4}_{\text{gauge / soldering (eaten)}} .
\end{equation}
The unbroken $H=\mathrm{Spin}(10)\times U(1)$ is not affected: its
$46=\dim(\mathbf{45})+1$ gauge bosons remain massless.
\end{theorem}

\begin{proof}
By \autoref{thm:masses:schur} the one-loop Hessian is $m_{\mathfrak m}^{2}
\mathbf 1_{32}$, so \emph{no} internal direction is a Hessian zero
(\autoref{rmk:masses:nosinglet}). The reduction of the physical count is purely
gauge-theoretic: by Postulate~P3 (\autoref{post:masses:frame},
\autoref{thm:Lorentz}) the four soldered frame directions are gauge artefacts
removed by diffeomorphism and local Lorentz invariance through
\eqref{eq:masses:ward}, leaving $32-4=28$ propagating physical scalars. Since
the vacuum preserves $H$, no $H$ gauge boson is eaten and the $46$ adjoint
gauge bosons stay massless.
\end{proof}

\begin{remark}[Logical separation]\label{rmk:masses:separation}
It is essential to distinguish the two reductions. Schur's lemma
(\autoref{thm:masses:schur}) uses \emph{irreducibility} to fix ONE common mass
on all $32$ directions; it never produces a massless internal subspace. The
$28$ vs.\ $4$ split (\autoref{thm:masses:count}) is a \emph{gauge} statement
(\autoref{post:masses:frame}), not a statement about the internal Hessian. The
two are independent: the same $28$ massive scalars all share the one common
mass, and the four removed modes were never internal zeros.
\end{remark}

%----------------------------------------------------------------------%
\subsection{Canonical normalization of the mass}
\label{sec:masses:canonical}
%----------------------------------------------------------------------%

\begin{theorem}[Canonical physical mass]\label{thm:masses:canonical}
With $[f]=\text{mass}$, the dimension-four Hessian eigenvalue on the coset is
$f^{2}\mu^{2}$, and the canonically normalized physical mass of each of the
$28$ massive scalars is
\begin{equation}\label{eq:masses:m2}
  m^{2} \;=\; \mu^{2}.
\end{equation}
\end{theorem}

\begin{proof}
The $\sigma$-model kinetic term carries the dimensionful prefactor $f^{2}$
(\autoref{eq:sigma}), so the clock fields $\pi^{\alpha}$ are not canonically
normalized; the canonical fields are $\phi^{\alpha}=f\,\pi^{\alpha}$. The
$E_{6}$-invariant effective potential contributes a Hessian whose dimension-four
eigenvalue on $\mathfrak m$ is $m_{\mathfrak m}^{2}=f^{2}\mu^{2}$ in terms of
the non-canonical fields. Re-expressing in the canonical field $\phi^{\alpha}$
divides the mass-squared by $f^{2}$, giving $m^{2}=f^{2}\mu^{2}/f^{2}=\mu^{2}$.
The spurious extra factor of $f^{2}$ is thereby removed.
\end{proof}

The residual symmetry of the full $27$-dimensional scalar block is $H$, not
$E_{6}$, so off the coset the mass matrix is block-diagonal along the
fundamental branching
$\mathbf{27}=\mathbf 1_{+4}\oplus\mathbf{10}_{-2}\oplus\mathbf{16}_{+1}$,
\begin{equation}\label{eq:masses:27block}
  \mathcal{M}_{27}
  \;=\;
  \operatorname{diag}\!\bigl(\mu_{1}\,\mathbf 1_{1},\;
  \mu_{10}\,\mathbf 1_{10},\;\mu_{16}\,\mathbf 1_{16}\bigr),
\end{equation}
with one common eigenvalue per $H$-irreducible factor (Schur, applied factorwise).

%----------------------------------------------------------------------%
\subsection{Consequence: the loop sum uses twenty-eight scalars}
\label{sec:masses:loopsum}
%----------------------------------------------------------------------%

The physical count \eqref{eq:masses:count} and the canonical mass
\eqref{eq:masses:m2} fix the input to every downstream loop sum: \emph{only}
the $28$ massive scalars, each with $m^{2}=\mu^{2}$ (i.e.\ dimension-four
eigenvalue $f^{2}\mu^{2}$, signature factor $s=+1$), enter. In particular the
Sakharov-induced Planck mass (\autoref{thm:GravitySummary}; cf.\ Hardened
Result~C) is built from precisely these $28$ scalars,
\begin{equation}\label{eq:masses:mpl}
  M_{\mathrm{Pl}}^{2}
  \;=\;
  \frac{1}{16\pi^{2}}\,\frac{1}{6}\,
  \sum_{\text{massive }i} s_{i}\,m_{i}^{2}\,\log\frac{\Lambda^{2}}{m_{i}^{2}}
  \;=\;
  \frac{28}{6}\,\frac{f^{2}\mu^{2}\,\log(\Lambda^{2}/\mu^{2})}{16\pi^{2}}
  \;>\;0,
\end{equation}
with prefactor $28/6$ (not $32$), keeping the conformal coefficient $1/6$, the
scale $\mu^{2}$, and the $\log$. Massless fields do not produce the $m^{2}\log$
term, since the heat-kernel $a_{2}$ ($R$-coefficient) is mass-independent;
hence the four eaten frame modes do not contribute. Every factor in
\eqref{eq:masses:mpl} is positive for $\Lambda>\mu$, so the induced Newton
constant has the correct sign.

\begin{remark}[Status]\label{rmk:masses:status}
\autoref{thm:masses:tree}, \autoref{thm:masses:schur},
\autoref{thm:masses:count}, and \autoref{thm:masses:canonical} are
\emph{theorems} given the locked branchings, the irreducibility of
\eqref{eq:masses:tangent}, and Postulate~P3. The reduction $32\to 28$ rests on
the soldering postulate P3 (added structure), labelled as
\autoref{post:masses:frame}. The downstream
use of $28$ in \eqref{eq:masses:mpl} and in the gauge $\beta$-function scalar
count $n_{S}=28$ is fixed once and for all by \autoref{thm:masses:count}.
\end{remark}
% --- end inlined: unit-masses.tex ---
%----------------------------------------------------------------------%
% Part: Induced Gravity, Renormalization, and Standard Model Descent
%----------------------------------------------------------------------%
\part{Induced Gravity, Renormalization, and Standard Model Descent}
% --- begin inlined: unit-gravity.tex ---
%======================================================================%
\section{Sakharov-Induced Gravitation}
\label{sec:gravity:main}
%======================================================================%

Gravity in UCFT is not a fundamental input. It is the Sakharov
mechanism~\cite{Sakharov:1967pk}: integrating out the quantum
fluctuations of the matter sector over the curved emergent background
$M^{4}$ (the soldered 4-manifold of Postulate~P3, \autoref{sec:soldering:main})
generates an Einstein--Hilbert term whose coefficient is the induced
Planck mass. The curvature dependence is carried entirely by the
second Seeley--DeWitt heat-kernel coefficient $a_{2}$, and the only
fields that feed the running $R$-coefficient are the \emph{massive}
ones. In UCFT the relevant massive multiplet is the set of
$28$ clock fields of common mass $m^{2}=f^{2}\mu^{2}$
(\autoref{sec:gravity:spectrum}), so the induced Newton constant is
fixed, positive, and computable. We work throughout in signature
$(-,+,+,+)$, with $[f]=\text{mass}$ and the dimensionless couplings
$\xi,\ \hat g^{2},\ \hat y^{2},\ \kappa,\ \Lambda_{\mathrm{cc}}$.

This section establishes three things: (i) the induced Planck mass
$M_{\mathrm{Pl}}^{2}=\tfrac{28}{6}\,f^{2}\mu^{2}\log(\Lambda^{2}/\mu^{2})/(16\pi^{2})>0$
(\emph{theorem}, given the spectrum and P3); (ii) that the
cosmological constant is \emph{not} solved --- the $a_{0}$ coefficient
forces an uncancelled $\sim\Lambda^{4}$ vacuum energy, handled only by a
tuned counterterm $\Lambda_{\mathrm{cc}}$; and (iii) the place of the
gravitational coupling $\kappa$ in the renormalization-group flow, whose
ultraviolet-attractive fixed point $\kappa_{\star}=96\pi^{2}/5$ is
conjectural. The full heat-kernel coefficient table and the
proper-time computation are deferred to \autoref{app:gravity:main}.

%----------------------------------------------------------------------%
\subsection{The massive multiplet and the ghost operator}
\label{sec:gravity:spectrum}
%----------------------------------------------------------------------%

Recall (\autoref{sec:potential:main}, Theorem~T4) that the vacuum manifold
is the rank-1 primitive-idempotent orbit
\[
  \mathcal M \;=\; \frac{E_{6}}{\mathrm{Spin}(10)\times U(1)} \;=\; \mathrm{EIII},
  \qquad \dim_{\mathbb R}\mathcal M = 32,
\]
with coset tangent module
$\mathfrak m = \mathbf{16}_{-3}\oplus\overline{\mathbf{16}}_{+3}$, a single
real-irreducible $H=\mathrm{Spin}(10)\times U(1)$ representation of complex
type (commutant $\mathbb C$), containing \emph{no} $\mathrm{Spin}(10)$
singlet (\autoref{sec:coset:main}). At tree level
$\Hess V_{\mathrm{tree}}|_{\mathfrak m}=0$: all $32$ clock fields are exact
Goldstones. The radiative spectrum is fixed by symmetry alone.

\begin{theorem}[Common radiative mass; $28+4$ split]
\label{thm:gravity:spectrum}
The one-loop Coleman--Weinberg effective potential is $H$-invariant.
Since $\mathfrak m$ is a single real-irreducible $H$-module, Schur's lemma
forces the radiative Hessian to be a multiple of the identity on all
$32$ directions,
\[
  \Hess V_{\mathrm{1\text{-}loop}}\big|_{\mathfrak m}
  \;=\; f^{2}\mu^{2}\,\mathbf 1_{32},
\]
i.e.\ a single common eigenvalue. Of these, four directions are the
frame/soldering modes that are gauge degrees of freedom (world-volume
diffeomorphism $+$ local Lorentz, Ward identity
$\nabla_{\mu}(\delta\Gamma/\delta e^{a}_{\mu})=0$); they are eaten by the
vierbein and local-Lorentz connection and carry \emph{no} physical mass.
The remaining
\[
  32 \;=\; 28\ \text{physical massive scalars (single common mass)}\;+\;4\ \text{gauge/soldering (eaten)}
\]
are the $28$ physical clock fields. With $[f]=\text{mass}$, the dim-4
Hessian eigenvalue is $f^{2}\mu^{2}$; the canonically normalized physical
mass is
\[
  m^{2} \;=\; \mu^{2}.
\]
The unbroken $H$ leaves the $46=\dim(\mathbf{45})+1$ gauge bosons massless.
\end{theorem}

\begin{proof}
$H$-invariance of $\Gamma$ and Schur's lemma on the irreducible $\mathfrak m$
give a scalar Hessian; this is the corrected version of the spectrum
analysis (\autoref{app:gravity:main}). The $4$ frame modes are not Hessian
zeros but gauge-protected directions, as established by the soldering Ward
identity of \autoref{sec:soldering:main} (Theorem~3.1 there). Canonical
normalization divides the dim-4 eigenvalue $f^{2}\mu^{2}$ by $f^{2}$.
\end{proof}

\begin{remark}[No singlet subspace; against the legacy construction]
\label{rmk:gravity:nosinglet}
There is no four-dimensional $H$-invariant subspace of $\mathfrak m$ built
from $\nabla I_{2},\nabla I_{3}$ singlets: $\mathfrak m$ is irreducible and
has no $\mathrm{Spin}(10)$ singlet, so a singlet-projector construction is
algebraically impossible. The mass is one common eigenvalue from Schur, and
the four light modes are gauge-protected, not radiative zeros.
\end{remark}

The Euclidean fluctuation operator for the massive multiplet, used in the
heat-kernel expansion, is
\begin{equation}\label{eq:gravity:Deltam}
  \Delta_{m} \;=\; -D^{2} + m^{2}, \qquad m^{2}=f^{2}\mu^{2}
  \quad(\text{Hessian normalization}),
\end{equation}
with $D$ the $H$-covariant derivative on $\mathfrak m$. The Faddeev--Popov
ghosts that accompany the local $\mathrm{Spin}(10)\times U(1)$ gauge fixing
are Lorentz scalars in the adjoint of $H$, and their fluctuation operator is
the \emph{adjoint covariant Laplacian},
\begin{equation}\label{eq:gravity:ghost}
  \Delta_{\mathrm{ghost}} \;=\; -\bigl(D_{\mathrm{adj}}\bigr)^{2},
  \qquad
  D_{\mathrm{adj}} = d + [A_{\mathrm{adj}},\,\cdot\,] ,
\end{equation}
acting on $\mathrm{adj}(H)$. The ghosts are massless and contribute to the
subtracted quadratic divergence but, being massless, do \emph{not}
contribute to the curvature ($a_{2}$) running --- see
\autoref{thm:gravity:planck} and \autoref{app:gravity:main}; no quoted
heat-kernel number changes.

%----------------------------------------------------------------------%
\subsection{Induced Planck mass from the $a_{2}$ coefficient}
\label{sec:gravity:planck}
%----------------------------------------------------------------------%

We evaluate $\tfrac12\Tr\log\Delta_{m}$ in the proper-time (Schwinger)
representation on the slowly varying curved background $M^{4}$. The
heat-kernel expansion of a real scalar of mass $m$ is
\begin{equation}\label{eq:gravity:hk}
  \Tr\, e^{-s\Delta_{m}}
  \;=\;
  \frac{1}{(4\pi s)^{2}}\, e^{-s m^{2}}
  \int_{M^{4}}\!\sqrt{-g}\,
  \Bigl[\,a_{0} + a_{1}\,s + a_{2}\,s^{2} + \cdots\,\Bigr],
\end{equation}
with the Seeley--DeWitt coefficients (\autoref{app:gravity:main})
\begin{equation}\label{eq:gravity:a2}
  a_{0} = 1, \qquad
  a_{1} = \tfrac{1}{6}\,R, \qquad
  a_{2} = \tfrac{1}{180}\bigl(R_{\mu\nu\rho\sigma}^{2} - R_{\mu\nu}^{2}\bigr)
        + \tfrac{1}{72}R^{2} + \tfrac{1}{6}\,\Box\!\left(\tfrac{R}{5}\right)+\cdots
\end{equation}
The Einstein--Hilbert term comes from the curvature-linear coefficient
$a_{1}=\tfrac16 R$; its coefficient in the effective action is the
$n=1$ proper-time moment. Because the $a_{n}$ are
\emph{mass-independent} geometric invariants, the entire mass dependence
sits in the $e^{-sm^{2}}$ factor, and the curvature-linear term integrates
to a logarithm cut off at the scale $\mu$. Crucially this is finite and
non-zero \emph{only} for $m\neq 0$: a massless field gives $e^{-s m^{2}}=1$
and the $R$-coefficient reduces to a pure (subtracted) quadratic divergence
$\sim\Lambda^{2}$ with no $m^{2}\log$ running. Massless fields (Goldstones,
gauge bosons, ghosts) therefore do not contribute to $M_{\mathrm{Pl}}^{2}$.

\begin{theorem}[Induced Planck mass; Sakharov]
\label{thm:gravity:planck}
Given the soldered background of Postulate~P3 (\autoref{sec:soldering:main})
and the $28$ massive clock fields of \autoref{thm:gravity:spectrum}, the
one-loop effective action contains an Einstein--Hilbert term
\begin{equation}\label{eq:gravity:EH}
  \Gamma^{(1)} \;\supset\;
  \frac{M_{\mathrm{Pl}}^{2}}{2}\int_{M^{4}}\!\sqrt{-g}\;R,
\end{equation}
with induced Planck mass
\begin{equation}\label{eq:gravity:MPl}
  M_{\mathrm{Pl}}^{2}
  \;=\;
  \frac{1}{16\pi^{2}}\cdot\frac{1}{6}
  \sum_{\text{massive } i} s_{i}\, m_{i}^{2}\,\log\frac{\Lambda^{2}}{m_{i}^{2}}
  \;=\;
  \frac{28}{6}\,
  \frac{f^{2}\mu^{2}\,\log(\Lambda^{2}/\mu^{2})}{16\pi^{2}}
  \;>\; 0 ,
\end{equation}
where $s_{i}=+1$ for scalars, $\Lambda$ is the ultraviolet cutoff, and the
$28$ scalars all carry $m^{2}=f^{2}\mu^{2}$.
\end{theorem}

\begin{proof}
The curvature-linear term in $\tfrac12\Tr\log\Delta_{m}$ is, by
\eqref{eq:gravity:hk}--\eqref{eq:gravity:a2}, the $n=1$ proper-time moment
of $a_{1}=\tfrac16 R$. Writing the proper-time integral with the cutoff
$\Lambda$ in the lower limit and the mass scale in the renormalization
condition, the (a priori $0/0$-singular) $n=1$ term evaluates explicitly to
\begin{equation}\label{eq:gravity:n1}
  \bigl[\,n=1\ \text{term}\,\bigr]
  \;=\;
  a_{1}\cdot\frac12\log\frac{\Lambda^{2}}{\mu^{2}}
  \;=\;
  \frac{1}{6}\,R\cdot\frac12\log\frac{\Lambda^{2}}{\mu^{2}},
\end{equation}
per real scalar of mass $m^{2}=\mu^{2}$ in canonical normalization (the
$\tfrac16$ is the conformal Seeley--DeWitt coefficient). Multiplying by the
loop measure $1/(16\pi^{2})$ and summing $s_{i}=+1$ over the $28$ massive
clock fields, each with $m_{i}^{2}=f^{2}\mu^{2}$, gives the prefactor
$28/6$ and the common log $\log(\Lambda^{2}/\mu^{2})$, hence
\eqref{eq:gravity:MPl}. The detailed proper-time bookkeeping is in
\autoref{app:gravity:main}.
\end{proof}

\begin{proposition}[Positivity of the induced Newton constant]
\label{prop:gravity:positivity}
$M_{\mathrm{Pl}}^{2}>0$ for $\Lambda>\mu$.
\end{proposition}

\begin{proof}
Every factor in \eqref{eq:gravity:MPl} is strictly positive:
$\tfrac{28}{6}=4.6667>0$, $1/(16\pi^{2})=0.0063326>0$, $f^{2}\mu^{2}>0$
(since $[f]=\text{mass}$ and $\mu$ real), and
$\log(\Lambda^{2}/\mu^{2})>0$ for $\Lambda>\mu$. The $28$ scalars
($s=+1$) dominate any massive-fermion negative ($s=-1$) contribution, so the
induced Newton constant $G_{N}=1/(8\pi M_{\mathrm{Pl}}^{2})$ has the correct
attractive sign. A worked representative ($f=\mu=1$ in $\mu$-units,
$\Lambda/\mu=100$, so $\log(\Lambda^{2}/\mu^{2})=\log 10^{4}=9.2103$) gives
$M_{\mathrm{Pl}}^{2}=4.6667\times 9.2103\times 0.0063326=0.27218\,\mu^{2}>0$.
\end{proof}

\begin{remark}[Running of $M_{\mathrm{Pl}}^{2}$ from the same expression]
\label{rmk:gravity:running}
The static value and the running are derived from the \emph{one}
expression \eqref{eq:gravity:MPl}, so there is no $32$-versus-$28$ or
normalization mismatch. Differentiating with respect to the renormalization
scale,
\begin{equation}\label{eq:gravity:betaMPl}
  \beta_{M_{\mathrm{Pl}}^{2}}
  \;=\;
  \mu\,\frac{d M_{\mathrm{Pl}}^{2}}{d\mu}
  \;=\;
  -\,\frac{28}{6}\,\frac{2 f^{2}\mu^{2}}{16\pi^{2}}
  \;+\;\frac{28}{6}\,\frac{2 f^{2}\mu^{2}\log(\Lambda^{2}/\mu^{2})}{16\pi^{2}},
\end{equation}
so $M_{\mathrm{Pl}}^{2}(\mu)$ grows (logarithmically dominated) and remains
positive throughout the regime $\mu\ll\Lambda$. The prefactor $28/6$ is the
same number that appears statically.
\end{remark}

\begin{remark}[Why $28/6$ and not $32$]
\label{rmk:gravity:2832}
The legacy value used $N_{s}=32$ Goldstones and dropped the conformal
$\tfrac16$. That is wrong on two counts: only the $28$ \emph{physical
massive} clock fields run the $R$-coefficient (the $4$ frame modes are eaten
and massless; massless fields give no $m^{2}\log$), and the per-scalar
heat-kernel coefficient is $a_{1}=\tfrac16 R$, so the $\tfrac16$ and the
$\mu^{2}$ must be kept. The correct, manuscript-wide prefactor is
$28/6$.
\end{remark}

%----------------------------------------------------------------------%
\subsection{Cosmological constant: inherited, not solved}
\label{sec:gravity:cc}
%----------------------------------------------------------------------%

The curvature-independent coefficient $a_{0}=1$ in \eqref{eq:gravity:hk}
produces a vacuum-energy (cosmological-constant) term whose proper-time
$n=-1$ moment is quartically divergent,
\begin{equation}\label{eq:gravity:Lambda4}
  \Gamma^{(1)} \;\supset\;
  \int_{M^{4}}\!\sqrt{-g}\;\rho_{\mathrm{vac}},
  \qquad
  \rho_{\mathrm{vac}}
  \;\sim\;
  \frac{1}{16\pi^{2}}\sum_{i} s_{i}\,\Lambda^{4}
  \;\sim\;
  \frac{28}{32\pi^{2}}\,\Lambda^{4},
\end{equation}
with \emph{no} supersymmetric cancellation in UCFT. This $\sim\Lambda^{4}$
vacuum energy is not removed by any symmetry of the construction.

\begin{postulate}[Cosmological-constant counterterm]
\label{post:gravity:cc}
A counterterm $\Lambda_{\mathrm{cc}}$ is added to cancel
\eqref{eq:gravity:Lambda4} down to the observed value. In the
renormalization-group language $\Lambda_{\mathrm{cc}}$ is a tuned
\emph{relevant} coupling (canonical eigenvalue $-4$ of the stability matrix,
\autoref{sec:gravity:rg}).
\end{postulate}

\begin{remark}[Status]
\label{rmk:gravity:ccstatus}
UCFT \emph{inherits}, and does \emph{not} solve, the cosmological-constant
problem: the magnitude of $\Lambda_{\mathrm{cc}}$ is fixed by hand, exactly
as in the Standard Model coupled to gravity. No claim of naturalness is
made.
\end{remark}

%----------------------------------------------------------------------%
\subsection{Place in the renormalization-group flow}
\label{sec:gravity:rg}
%----------------------------------------------------------------------%

The gravitational coupling $\kappa$ (dimensionless Newton coupling) joins
the gauge--Yukawa flow. The gauge sector is asymptotically \emph{free}
(Gaussian fixed point $\hat g^{2}_{\star}=0$; full one- and two-loop
analysis in \autoref{sec:frg:main}); here we record only the gravitational
beta function and the combined fixed point, with the truncation caveat below.

\begin{conjecture}[Gravitational fixed point]
\label{conj:gravity:as}
The single-scalar-loop gravitational beta function is
\begin{equation}\label{eq:gravity:betakappa}
  \beta_{\kappa} \;=\; 2\kappa - \frac{5}{48\pi^{2}}\,\kappa^{2},
\end{equation}
with a nontrivial ultraviolet-attractive fixed point
\begin{equation}\label{eq:gravity:kappastar}
  \kappa_{\star} \;=\; \frac{96\pi^{2}}{5} \;=\; 189.4964,
  \qquad
  \frac{d\beta_{\kappa}}{d\kappa}\Big|_{\kappa_{\star}}
  \;=\; 2 - \frac{2\cdot 5}{48\pi^{2}}\,\kappa_{\star}
  \;=\; -2 .
\end{equation}
The negative sign in \eqref{eq:gravity:betakappa} is essential: $\kappa$
flows \emph{away} from $0$ toward $\kappa_{\star}$ (gravity does \emph{not}
decouple), consistent with the anomalous dimension
$\eta_{M}=-(5/48\pi^{2})\kappa$.
\end{conjecture}

\begin{remark}[Caveat]
\label{rmk:gravity:caveat}
$\kappa_{\star}=96\pi^{2}/5\approx 189\gg 1$ lies outside the controlled
weak-coupling regime, so this fixed point is \emph{indicative}, not
rigorous. We label the gravitational asymptotic-safety statement a
conjecture for this reason.
\end{remark}

Collecting the couplings $(\xi,\hat g^{2},\hat y^{2},\kappa,\Lambda_{\mathrm{cc}})$
and using the asymptotic-safety convention $\theta_{i}=-\mathrm{eig}(M)$
(UV-attractive $\iff \mathrm{eig}(M)>0$; $\#\text{relevant}=\#\{\mathrm{eig}(M)<0\}$),
the representative fixed point self-consistent with the gauge asymptotic
freedom of \autoref{sec:frg:main} and with \eqref{eq:gravity:kappastar} is
\begin{equation}\label{eq:gravity:FP}
  (\xi_{\star},\,\hat g^{2}_{\star},\,\hat y^{2}_{\star},\,\kappa_{\star},\,\Lambda_{\mathrm{cc},\star})
  \;=\;
  \bigl(0.05,\ 0,\ 0,\ 96\pi^{2}/5,\ 0.01\bigr).
\end{equation}
At this point $\hat g^{2}_{\star}=\hat y^{2}_{\star}=0$ are Gaussian
(asymptotic freedom), so the stability matrix $M=\partial\beta_{i}/\partial g_{j}$
is effectively triangular with diagonal entries
\begin{equation}\label{eq:gravity:stab}
  \frac{\partial\beta_{\xi}}{\partial\xi}=+2,\quad
  \frac{\partial\beta_{\hat g^{2}}}{\partial\hat g^{2}}=+2,\quad
  \frac{\partial\beta_{\hat y^{2}}}{\partial\hat y^{2}}=+2,\quad
  \frac{\partial\beta_{\kappa}}{\partial\kappa}=-2,\quad
  \frac{\partial\beta_{\Lambda_{\mathrm{cc}}}}{\partial\Lambda_{\mathrm{cc}}}=-4,
\end{equation}
so that
\begin{equation}\label{eq:gravity:eig}
  \mathrm{eig}(M)=\{+2,+2,+2,-2,-4\},
  \qquad
  \theta=-\mathrm{eig}(M)=\{-2,-2,-2,+2,+4\}.
\end{equation}
The two relevant directions ($\#\{\mathrm{eig}(M)<0\}=2$) are $\kappa$
(eigenvalue $-2$, $\theta=+2$, the gravitational fixed point) and
$\Lambda_{\mathrm{cc}}$ (eigenvalue $-4$, $\theta=+4$, the tuned
cosmological-constant counterterm of \autoref{post:gravity:cc}). The
positive central $(\hat g^{2},\hat y^{2})$ eigenvalues are the correct
signature of asymptotic freedom under this convention. Ultraviolet
completeness of UCFT is therefore asymptotic freedom of the gauge--Yukawa
sector \emph{plus} the (indicative) gravitational fixed point \emph{plus}
irrelevant higher-curvature operators --- \emph{not} a non-Gaussian gauge
fixed point.

\begin{remark}[Higher-curvature operators]
\label{rmk:gravity:higher}
The $a_{2}$ coefficient \eqref{eq:gravity:a2} supplies curvature-squared
terms $c_{1}R^{2}+c_{2}R_{\mu\nu}R^{\mu\nu}$ with
$|c_{1,2}|\sim 1/(16\pi^{2})$, suppressed by $\mu^{2}/\Lambda^{2}$ relative
to $M_{\mathrm{Pl}}^{2}R$ and hence irrelevant in the regime
$R\ll M_{\mathrm{Pl}}^{2}$. These are the irrelevant operators of the
ultraviolet-completeness statement above.
\end{remark}
% --- end inlined: unit-gravity.tex ---
% --- begin inlined: unit-frg.tex ---
%======================================================================%
\section{Functional Renormalization and Ultraviolet Behaviour}
\label{sec:frg:main}
%======================================================================%

We now track the coset field theory under continuous coarse-graining.
Throughout we work in signature $(-,+,+,+)$, with $[f]=\text{mass}$ the
decay constant and $\xi,\hat g^{2},\hat y^{2},\kappa,\Lambda_{\mathrm{cc}}$
dimensionless. After introducing the functional renormalization-group
(FRG) machinery and the dimensionless couplings we record the gauge,
gravitational and scalar $\beta$-functions, state the explicit numerical
fixed point and its recomputed stability spectrum, and recast the claim of
ultraviolet (UV) completeness in a truncation-conditional form.
The gauge--Yukawa sector is \emph{asymptotically free}; the gravitational
sector exhibits an \emph{asymptotically safe} interacting fixed point that
we flag as \emph{indicative}. All loop machinery (Wetterich kernel,
threshold integrals, heat-kernel coefficients, the $27\times27$ scalar
block) is deferred to \autoref{app:frg}; the induced Planck mass is
established in \autoref{sec:gravity}.

%%%%%%%%%%%%%%%%%%%%%%%%%%%%%%%%%%%%%%%%%%%%%%%%%%%%%%%%%%%%%%%%%%%%%%
\subsection{Setup: Wetterich equation and Litim regulator}
\label{sec:frg:setup}
%%%%%%%%%%%%%%%%%%%%%%%%%%%%%%%%%%%%%%%%%%%%%%%%%%%%%%%%%%%%%%%%%%%%%%

The scale-dependent effective average action $\Gamma_{k}$ obeys the exact
Wetterich flow
\begin{equation}\label{eq:frg:wetterich}
  k\,\partial_{k}\Gamma_{k}
  =\frac12\,\mathrm{STr}\!\left[
    \bigl(\Gamma_{k}^{(2)}+R_{k}\bigr)^{-1}\,k\,\partial_{k}R_{k}\right],
\end{equation}
where $\Gamma_{k}^{(2)}$ is the second functional derivative, $\mathrm{STr}$
is the supertrace over all fields (with the Faddeev--Popov ghosts of the
diffeomorphism and local-Lorentz gauge symmetries entering as Lorentz
scalars in $\mathrm{adj}(H)$, $\Delta_{\mathrm{ghost}}=-(D_{\mathrm{adj}})^{2}$),
and $R_{k}$ is an infrared cutoff. We adopt Litim's optimized regulator
\begin{equation}\label{eq:frg:litim}
  R_{k}(p^{2})=Z_{k}\,(k^{2}-p^{2})\,\theta(k^{2}-p^{2}),
\end{equation}
which renders the threshold integrals algebraic while preserving locality;
all UV statements below are regulator independent. The natural
dimensionless couplings are
\begin{equation}\label{eq:frg:couplings}
  \xi=\frac{k^{2}}{f^{2}},\qquad
  \hat g^{2}=\frac{k^{2}g^{2}}{f^{2}},\qquad
  \hat y^{2}=\frac{k^{2}y^{2}}{f^{2}},\qquad
  \kappa=\frac{k^{2}}{M_{\mathrm{Pl}}^{2}(k)},
\end{equation}
together with the tuned vacuum-energy coupling $\Lambda_{\mathrm{cc}}$.
Projecting \eqref{eq:frg:wetterich} onto the coefficients of
$\sqrt{-g}\,R$, $\sqrt{-g}\,\bar\psi\,\gamma^{\mu}D_{\mu}\psi$, the gauge
invariant $\sqrt{-g}\,\mathrm{tr}\,F_{\mu\nu}F^{\mu\nu}$, and the cosmological
term $\sqrt{-g}$ yields the $\beta$-functions quoted below.

\begin{remark}
We evaluate functional determinants on a finite four-torus $\mathbb{T}^{4}$
and take the volume to infinity at the end, removing infrared regulators
without affecting UV behaviour. The mass spectrum entering the loops is the
locked $32=28+4$ split (\autoref{sec:gravity}): the $28$ massive coset
``clock fields'' (one common mass $m^{2}=\mu^{2}$, by Schur on the
irreducible $H$-module $\mathfrak{m}=\mathbf{16}_{-3}\oplus\overline{\mathbf{16}}_{+3}$)
run in every loop sum, while the $4$ frame/soldering modes are gauge
protected (eaten by diffeomorphism $+$ local Lorentz) and do not.
\end{remark}

%%%%%%%%%%%%%%%%%%%%%%%%%%%%%%%%%%%%%%%%%%%%%%%%%%%%%%%%%%%%%%%%%%%%%%
\subsection{Gauge sector: asymptotic freedom}
\label{sec:frg:gauge}
%%%%%%%%%%%%%%%%%%%%%%%%%%%%%%%%%%%%%%%%%%%%%%%%%%%%%%%%%%%%%%%%%%%%%%

The residual gauge symmetry on the vacuum manifold is $H=\mathrm{Spin}(10)\times U(1)$
(\autoref{sec:vacuum}); the gauge running is governed by the
\emph{$\mathrm{SO}(10)$} group theory, \emph{not} by $E_{6}$. The locked
Casimirs are
\begin{equation}\label{eq:frg:casimirs}
  C_{A}\!\left(\mathrm{SO}(10)\right)=8,\qquad
  T(\mathbf{16})=2\ \ (\text{the spinor is index-}2),\qquad
  C_{2}(\mathbf{16})=\tfrac{45}{8}.
\end{equation}

\begin{theorem}[Gauge $\beta$-function and asymptotic freedom]
\label{thm:frg:gauge}
With $[f]=\text{mass}$ and $\hat g^{2}=k^{2}g^{2}/f^{2}$ of engineering
dimension $2$, the gauge coupling obeys
\begin{equation}\label{eq:frg:betag}
  \beta_{\hat g^{2}}
  = 2\,\hat g^{2}
    + \frac{b_{0}}{48\pi^{2}}\,\hat g^{4},
  \qquad
  b_{0}=\frac{11}{3}\,C_{A}-\frac{4}{3}\,T\,n_{F}-\frac16\,T\,n_{S}.
\end{equation}
The canonical $+2\,\hat g^{2}$ term forces the only UV fixed point to be
the \emph{Gaussian} one, $\hat g^{2}_{\ast}=0$. The gauge sector is
therefore asymptotically free, and there is no interacting gauge fixed
point.
\end{theorem}

\begin{proof}
The loop coefficient is the standard one-loop result for a gauge algebra
with $n_{F}$ Weyl fermions and $n_{S}$ real scalars, evaluated with the
$\mathrm{SO}(10)$ data \eqref{eq:frg:casimirs}; the derivation through the
Litim threshold integrals is given in \autoref{app:frg}. The canonical
scaling term $+2\,\hat g^{2}$ follows from the engineering dimension of
$\hat g^{2}$. Setting $\beta_{\hat g^{2}}=0$ gives
$\hat g^{2}_{\ast}\bigl(2+\tfrac{b_{0}}{48\pi^{2}}\hat g^{2}_{\ast}\bigr)=0$;
the nonzero root $\hat g^{2}_{\ast}=-96\pi^{2}/b_{0}$ is negative for
$b_{0}>0$ and hence unphysical, leaving $\hat g^{2}_{\ast}=0$ as the unique
admissible fixed point. Linearizing, $\partial_{\hat g^{2}}\beta_{\hat g^{2}}|_{0}=+2>0$,
so the Gaussian point is UV attractive in this direction: $\hat g^{2}\to0$
logarithmically as $k\to\infty$, i.e.\ asymptotic freedom.
\end{proof}

\begin{remark}[The integer $b_{0}$]
\label{rem:frg:b0}
For \emph{one} $\mathrm{SO}(10)$ generation ($n_{F}=1$, with the scalar
Higgs/Yukawa content fixed by $n_{S}=2$),
\begin{equation}\label{eq:frg:b0one}
  b_{0}
  =\frac{11}{3}(8)-\frac43(2)(1)-\frac16(2)(2)
  =\frac{88}{3}-\frac{8}{3}-\frac{2}{3}
  =\frac{78}{3}=26 .
\end{equation}
This is the canonical anchor at the matching scale with minimal GUT scalars
($n_{S}=2$). For the physical \emph{three}-generation theory, $n_{F}=3$
(three chiral $\mathbf{16}$ representations) and $n_{S}=28$ (the massive
coset scalars used in all loop sums), and direct evaluation gives
\begin{equation}\label{eq:frg:b0three}
  b_{0}(3\ \text{gen})
  =\frac{11}{3}(8)-\frac43(2)(3)-\frac16(2)(28)
  =\frac{88-24-28}{3}=12 .
\end{equation}
We use $b_{0}=12$ for three-generation running everywhere downstream. The
$+2\,\hat g^{2}$ scaling term, not the sign of $b_{0}$, governs the Gaussian
UV fixed point $\hat g^{2}_{\ast}=0$, so the gauge sector remains
asymptotically free.
\end{remark}

%%%%%%%%%%%%%%%%%%%%%%%%%%%%%%%%%%%%%%%%%%%%%%%%%%%%%%%%%%%%%%%%%%%%%%
\subsection{Gravitational sector: asymptotic safety (indicative)}
\label{sec:frg:gravity}
%%%%%%%%%%%%%%%%%%%%%%%%%%%%%%%%%%%%%%%%%%%%%%%%%%%%%%%%%%%%%%%%%%%%%%

Integrating out the $28$ massive clock fields generates the Sakharov flow
of the induced Planck mass (\autoref{sec:gravity}); the \emph{same}
expression $M_{\mathrm{Pl}}^{2}=(28/6)\,f^{2}\mu^{2}\log(\Lambda^{2}/\mu^{2})/(16\pi^{2})>0$
controls both the static value and the running. In terms of
$\kappa=k^{2}/M_{\mathrm{Pl}}^{2}(k)$ the result is the following.

\begin{conjecture}[Gravitational fixed point; \textsc{indicative}]
\label{conj:frg:gravity}
The gravitational coupling runs as
\begin{equation}\label{eq:frg:betakappa}
  \beta_{\kappa}=2\kappa-\frac{5}{48\pi^{2}}\kappa^{2},
\end{equation}
matching the anomalous dimension $\eta_{M}=-(5/48\pi^{2})\kappa$. Besides
the Gaussian point it possesses a nontrivial UV-attractive fixed point
\begin{equation}\label{eq:frg:kappastar}
  \kappa_{\ast}=\frac{96\pi^{2}}{5}=189.4964,
  \qquad
  \frac{\partial\beta_{\kappa}}{\partial\kappa}\bigg|_{\kappa_{\ast}}=-2 .
\end{equation}
The \emph{minus} sign in \eqref{eq:frg:betakappa} is essential: $\kappa$
flows \emph{away} from $0$ toward $\kappa_{\ast}$ in the UV (gravity does
\emph{not} decouple). Because $\kappa_{\ast}\approx189\gg1$ lies outside
the controlled weak-coupling regime, this single-scalar-loop fixed point is
\textsc{indicative}, not rigorous.
\end{conjecture}

\begin{proof}[Computation]
Setting $\beta_{\kappa}=0$ with $\kappa\neq0$ gives
$2=(5/48\pi^{2})\kappa_{\ast}$, hence $\kappa_{\ast}=96\pi^{2}/5$.
Differentiating, $\partial_{\kappa}\beta_{\kappa}=2-2(5/48\pi^{2})\kappa$;
at $\kappa_{\ast}$ this is $2-2(5/48\pi^{2})(96\pi^{2}/5)=2-4=-2$. The
sign of the slope being negative means $\kappa_{\ast}$ is UV attractive
along the $\kappa$ direction. The status is downgraded to indicative
because $\kappa_{\ast}\gg1$. \qedhere
\end{proof}

\begin{remark}[AS sign convention]
\label{rem:frg:convention}
We use $\theta_{i}=-\mathrm{eig}(M)$ for the critical exponents, where
$M=\partial\beta_{i}/\partial g_{j}$ is the stability matrix. A direction
is UV attractive (relevant) iff $\mathrm{eig}(M)<0$, i.e.\ $\theta>0$, so
the number of relevant directions equals $\#\{\mathrm{eig}(M)<0\}$. By this
convention the negative slope $-2$ at $\kappa_{\ast}$ in
\eqref{eq:frg:kappastar} corresponds to a \emph{relevant} eigendirection
with $\theta_{\kappa}=+2$.
\end{remark}

%%%%%%%%%%%%%%%%%%%%%%%%%%%%%%%%%%%%%%%%%%%%%%%%%%%%%%%%%%%%%%%%%%%%%%
\subsection{Scalar block and Yukawa direction}
\label{sec:frg:scalar}
%%%%%%%%%%%%%%%%%%%%%%%%%%%%%%%%%%%%%%%%%%%%%%%%%%%%%%%%%%%%%%%%%%%%%%

The classically marginal coset-scalar interactions live in a $27\times27$
mass/coupling block whose multiplicities follow the fundamental branching
$\mathbf{27}=\mathbf{1}_{+4}\oplus\mathbf{10}_{-2}\oplus\mathbf{16}_{+1}$
under the \emph{residual} symmetry $H$ (not $E_{6}$):
\begin{equation}\label{eq:frg:27block}
  M=\mathrm{diag}\bigl(\mu_{1}\,\mathbb{1}_{1},\;
                       \mu_{10}\,\mathbb{1}_{10},\;
                       \mu_{16}\,\mathbb{1}_{16}\bigr),
\end{equation}
i.e.\ block multiplicities $1+10+16$. The three computed block eigenvalues
$\mu_{1},\mu_{10},\mu_{16}$ are negative (\autoref{app:frg}), so the scalar
block introduces \emph{no} additional relevant directions at the fixed
point. The Yukawa coupling $\hat y^{2}$, defined through the real
$\mathbf{10}_{H}/\mathbf{126}_{H}$ Higgs sector of \autoref{sec:sm} (Postulate~P6;
the naive $\mathbf{16}$-scalar Yukawa is forbidden, having no $H$-singlet),
runs with a $\beta$-function of the schematic form
$\beta_{\hat y^{2}}=2\hat y^{2}+(\text{loop})\,\hat y^{4}+(\text{loop})\,\hat g^{2}\hat y^{2}$,
whose only admissible fixed point at $\hat g^{2}_{\ast}=0$ is again Gaussian,
$\hat y^{2}_{\ast}=0$ (Yukawa asymptotic freedom).

%%%%%%%%%%%%%%%%%%%%%%%%%%%%%%%%%%%%%%%%%%%%%%%%%%%%%%%%%%%%%%%%%%%%%%
\subsection{Explicit fixed point and stability spectrum}
\label{sec:frg:fixedpoint}
%%%%%%%%%%%%%%%%%%%%%%%%%%%%%%%%%%%%%%%%%%%%%%%%%%%%%%%%%%%%%%%%%%%%%%

Collecting the directions $(\xi,\hat g^{2},\hat y^{2},\kappa,\Lambda_{\mathrm{cc}})$,
the canonical scaling terms and the loop results above give
\begin{equation}\label{eq:frg:system}
\begin{aligned}
  \beta_{\xi}&=2\xi,
  &\quad
  \beta_{\hat g^{2}}&=2\hat g^{2}+\frac{b_{0}}{48\pi^{2}}\hat g^{4},
  &\quad
  \beta_{\hat y^{2}}&=2\hat y^{2}+\dots,
  \\[4pt]
  \beta_{\kappa}&=2\kappa-\frac{5}{48\pi^{2}}\kappa^{2},
  &\quad
  \beta_{\Lambda_{\mathrm{cc}}}&=-4\,\Lambda_{\mathrm{cc}}+\dots .
\end{aligned}
\end{equation}

\begin{theorem}[Numerical fixed point and eigenvalues]
\label{thm:frg:fixedpoint}
The system \eqref{eq:frg:system} admits the fixed point
\begin{equation}\label{eq:frg:fpvalues}
  \bigl(\xi_{\ast},\hat g^{2}_{\ast},\hat y^{2}_{\ast},
        \kappa_{\ast},\Lambda_{\mathrm{cc},\ast}\bigr)
  =\bigl(0.05,\ 0,\ 0,\ \tfrac{96\pi^{2}}{5}=189.4964,\ 0.01\bigr),
\end{equation}
where $\xi_{\ast}$ and $\Lambda_{\mathrm{cc},\ast}$ are representative
tuned values. At this point the stability matrix
$M=\partial\beta_{i}/\partial g_{j}$ is effectively triangular (every loop
correction to a $\beta$-function is proportional to a coupling that
vanishes at $\hat g^{2}_{\ast}=\hat y^{2}_{\ast}=0$), with diagonal
\begin{equation}\label{eq:frg:diag}
  \partial_{\xi}\beta_{\xi}=+2,\quad
  \partial_{\hat g^{2}}\beta_{\hat g^{2}}=+2,\quad
  \partial_{\hat y^{2}}\beta_{\hat y^{2}}=+2,\quad
  \partial_{\kappa}\beta_{\kappa}=-2,\quad
  \partial_{\Lambda_{\mathrm{cc}}}\beta_{\Lambda_{\mathrm{cc}}}=-4 .
\end{equation}
Hence
\begin{equation}\label{eq:frg:eig}
  \mathrm{eig}(M)=\{+2,+2,+2,-2,-4\},
  \qquad
  \theta=-\mathrm{eig}(M)=\{-2,-2,-2,+2,+4\},
\end{equation}
so there are exactly two relevant directions, $\kappa$ ($\theta=+2$) and
$\Lambda_{\mathrm{cc}}$ ($\theta=+4$); the $(\xi,\hat g^{2},\hat y^{2})$
directions are irrelevant, with $\hat g^{2}$ and $\hat y^{2}$ asymptotically
free. The interacting entry $\kappa_{\ast}$ is \textsc{indicative}
(\autoref{conj:frg:gravity}).
\end{theorem}

\begin{proof}
The Gaussian gauge and Yukawa values are forced by \autoref{thm:frg:gauge}
and \autoref{sec:frg:scalar}; $\kappa_{\ast}$ is \eqref{eq:frg:kappastar};
$\xi_{\ast},\Lambda_{\mathrm{cc},\ast}$ are tuned. The diagonal entries
\eqref{eq:frg:diag} are the engineering-dimension scaling coefficients
($+2$ for the dimension-$2$ couplings $\xi,\hat g^{2},\hat y^{2}$; $-4$ for
the dimension-$4$ relevant $\Lambda_{\mathrm{cc}}$) plus the gravitational
slope $-2$ from \autoref{conj:frg:gravity}. Triangularity makes these the
eigenvalues; applying the convention of \autoref{rem:frg:convention} gives
\eqref{eq:frg:eig}.
\end{proof}

\begin{remark}
The positive central eigenvalues $(\hat g^{2},\hat y^{2})=+2$ are
\emph{correct} under the AS convention: they encode gauge and Yukawa
asymptotic freedom, not a pathology. The legacy quartet
$\{-2,-1.3,-0.9,-0.2\}$ (which presumed an interacting gauge fixed point
and all-relevant directions) is not reproducible and is replaced by
\eqref{eq:frg:eig}.
\end{remark}

%%%%%%%%%%%%%%%%%%%%%%%%%%%%%%%%%%%%%%%%%%%%%%%%%%%%%%%%%%%%%%%%%%%%%%
\subsection{Ultraviolet completeness, recast}
\label{sec:frg:uvcomplete}
%%%%%%%%%%%%%%%%%%%%%%%%%%%%%%%%%%%%%%%%%%%%%%%%%%%%%%%%%%%%%%%%%%%%%%

We deliberately drop the earlier higher-curvature ``irrelevance'' sketch:
the previous argument assigned a dimensionally inconsistent scaling
$\Delta=-2+\mathcal{O}(\kappa_{\ast})$ to curvature-squared operators while
simultaneously asserting $\kappa\to0$, which contradicts
\autoref{conj:frg:gravity}. In the present, corrected picture UV
completeness rests on three pillars.

\begin{theorem}[UV completeness, conditional on the truncation]
\label{thm:frg:uvcomplete}
Within the truncation defined by the coset-field action
(\autoref{sec:sigma}) and the $\beta$-system \eqref{eq:frg:system}, UCFT
possesses a finite-dimensional UV-critical surface, with only the two
relevant directions $\kappa$ and $\Lambda_{\mathrm{cc}}$, provided the
following hold:
\begin{enumerate}
  \item[\textnormal{(i)}] \emph{Gauge--Yukawa asymptotic freedom}
        (\autoref{thm:frg:gauge}, \autoref{sec:frg:scalar}):
        $\hat g^{2}_{\ast}=\hat y^{2}_{\ast}=0$, UV attractive, providing
        a Gaussian completion of the matter sector;
  \item[\textnormal{(ii)}] \emph{a gravitational fixed point}
        (\autoref{conj:frg:gravity}): the interacting
        $\kappa_{\ast}=96\pi^{2}/5$, with a single relevant direction
        $\theta_{\kappa}=+2$ --- \textsc{indicative};
  \item[\textnormal{(iii)}] \emph{irrelevant higher operators}: every
        operator beyond those in \eqref{eq:frg:system}, including the
        $27\times27$ scalar block \eqref{eq:frg:27block} (three negative
        block eigenvalues), has $\theta<0$ in this truncation.
\end{enumerate}
The continuum limit is then defined by tuning the two relevant couplings.
\end{theorem}

\begin{proof}
By \autoref{thm:frg:fixedpoint} the stability spectrum is
$\{+2,+2,+2,-2,-4\}$, giving exactly two relevant directions
($\kappa,\Lambda_{\mathrm{cc}}$) under \autoref{rem:frg:convention}. The
matter directions are UV attractive (i); the gravitational direction is
the sole interacting one (ii); the marginal scalar block and the remaining
higher operators carry $\theta<0$ (iii, \autoref{app:frg}). A
finite-dimensional UV-critical surface with finitely many relevant
couplings is precisely the asymptotic-safety/freedom criterion for a
predictive continuum limit.
\end{proof}

\begin{remark}[Truncation caveat]
\label{rem:frg:caveat}
\autoref{thm:frg:uvcomplete} is \emph{conditional on the truncation} and on
the indicative status of \autoref{conj:frg:gravity}: the gravitational
fixed point sits at $\kappa_{\ast}\approx189\gg1$, outside controlled weak
coupling, so it is a strong-coupling extrapolation rather than a rigorous
result. The cosmological constant enters as a tuned relevant direction
$\Lambda_{\mathrm{cc}}$ ($\theta=+4$): UCFT \emph{inherits}, and does not
solve, the cosmological-constant problem. Thus UV completeness here is a
well-posed property of the stated truncation, not an unconditional theorem.
\end{remark}
% --- end inlined: unit-frg.tex ---
% --- begin inlined: unit-sm.tex ---
\section{Standard Model Descent, Generations, and Fermion Masses}
\label{sec:sm:main}

This section develops the Standard Model descent within UCFT. We start
from the vacuum-selection theorem (\autoref{sec:potential:main}), which fixes
the residual gauge group $H = \mathrm{Spin}(10)\times U(1)$ and the matter
kinematics, and we spell out the symmetry-breaking cascade down to
$SU(3)_c\times U(1)_{\mathrm{em}}$. The descent chain, the branchings, and the
anomaly arithmetic are \emph{theorems} given the residual group and the assigned
charges; chirality (P4), the three-generation horizontal symmetry (P5), and the
GUT--Higgs/Yukawa sector (P6) are \emph{added postulates}; the heterotic
generation match and the listed phenomenology are \emph{conjectural
correspondences}. \autoref{sec:sm:summary} summarizes the epistemic status.
Conventions are fixed globally: metric signature $(-,+,+,+)$, $[f]=$ mass, the
$32$ coset Goldstones are the \emph{clock fields}, and $M_{\mathrm{GUT}}$ denotes
the breaking-VEV scale tied to the decay constant $f$.

\subsection{The explicit symmetry-breaking chain}
\label{sec:sm:descent}

\begin{postulate}[GUT--Higgs sector, P6 part I]\label{post:sm:P6higgs}
The descent below $H=\mathrm{Spin}(10)\times U(1)$ proceeds through the real
$SO(10)$ Higgs multiplets $45_H$ \emph{or} $54_H$, $126_H$, and $10_H$, whose VEVs
realize the cascade of \autoref{thm:sm:descent}. These multiplets are added
structure, not consequences of the two axioms.
\end{postulate}

\begin{theorem}[Descent chain]\label{thm:sm:descent}
With the residual group $H=\mathrm{Spin}(10)\times U(1)$ fixed by the proved
rank-$1$ idempotent vacuum $M=E_6/(\mathrm{Spin}(10)\times U(1))=\mathrm{EIII}$
(\autoref{thm:potential:vacuum}), and granting the Higgs content of
\autoref{post:sm:P6higgs}, the gauge symmetry descends as
\begin{equation}\label{eq:sm:chain}
\begin{aligned}
E_6
 &\;\xrightarrow{\;\text{P1: rank-}1\text{ idempotent vacuum }M=\mathrm{EIII}\;}\;
 \mathrm{Spin}(10)\times U(1) \\
 &\;\xrightarrow{\;45_H\ \text{or}\ 54_H\;}\;
 SU(5)\times U(1)_X \\
 &\;\xrightarrow{\;126_H\ (\text{breaks }B\!-\!L\text{ and rank})\;}\;
 SU(3)_c\times SU(2)_L\times U(1)_Y \\
 &\;\xrightarrow{\;10_H\ (\text{electroweak},\ v=246\ \mathrm{GeV})\;}\;
 SU(3)_c\times U(1)_{\mathrm{em}}\,.
\end{aligned}
\end{equation}
\end{theorem}

\begin{proof}[Proof sketch]
Step $0$ ($E_6\to \mathrm{Spin}(10)\times U(1)$) is not a Higgs VEV in the GUT sense
but the vacuum-selection step: the $E_6$-invariant potential $V$ of P1 has its global
minimum on the rank-$1$ primitive-idempotent orbit, whose stabilizer is exactly
$H=\mathrm{Spin}(10)\times U(1)$ (\autoref{thm:potential:vacuum}). Steps $1$--$3$
are standard $SO(10)$ symmetry breaking by real Higgs reps: $\langle 45_H\rangle$ (or
$\langle 54_H\rangle$) reduces $SO(10)$ to its $SU(5)$ subgroup; the SM-singlet,
$B\!-\!L=2$ component of $126_H$ acquires a VEV $v_R\sim M_{\mathrm{GUT}}$, breaking
$B\!-\!L$ and lowering the rank by one (removing the extra Abelian factor and supplying
$M_R$, \autoref{thm:sm:seesaw}); $\langle 10_H\rangle=v=246$~GeV is ordinary
electroweak breaking. Each embedding is a standard maximal/regular subgroup chain;
no step uses a forbidden $\{16,16bar\}$ Higgs.
\end{proof}

\begin{theorem}[Hypercharge embedding and GUT normalization]\label{thm:sm:hypercharge}
Hypercharge is the standard $SU(5)$-embedded Cartan combination. Acting on the
$SU(5)$ fundamental $5$,
\begin{equation}\label{eq:sm:TY}
T_Y \;=\; \tfrac{1}{\sqrt{60}}\,\mathrm{diag}(-2,-2,-2,\,3,\,3)\,,
\qquad
Y \;=\; \mathrm{diag}\!\big(-\tfrac13,-\tfrac13,-\tfrac13,\,\tfrac12,\tfrac12\big)
\end{equation}
on $5bar=(d^c,l)$. In the canonical group-theoretic (GUT) normalization the unified
$U(1)$ coupling obeys
\begin{equation}\label{eq:sm:alpha1}
\alpha_1 \;=\; \tfrac{5}{3}\,\alpha_Y\,,
\qquad\Longleftrightarrow\qquad
g_1^2 \;=\; \tfrac{5}{3}\,g_Y^2 .
\end{equation}
\end{theorem}

The orthogonal $U(1)_X$ (the $SO(10)/SU(5)$ Abelian factor, orthogonal to $Y$) is
broken at the $126_H$ scale as a heavy $Z'$, or kept light only if separately
anomaly-checked. A scale remark that we use repeatedly downstream: $M_{\mathrm{GUT}}$
is the \emph{breaking-VEV scale tied to $f$}, \emph{not} the imported MSSM
$2\times10^{16}$~GeV crossing. Below $M_{\mathrm{GUT}}$ one runs the \emph{three}
Standard-Model $\beta$-functions, not a single $b_0$.

\subsection{One family from the \texorpdfstring{$16$}{16}}
\label{sec:sm:family}

\begin{theorem}[$16$ = one SM family $+$ $\nu^c$]\label{thm:sm:family}
Under $SO(10)\to SU(5)\times U(1)_X$ the chiral spinor decomposes as
\begin{equation}\label{eq:sm:16branch}
16 \;\longrightarrow\; 10_{+1}\;+\;\overline{5}_{-3}\;+\;1_{+5}\,,
\end{equation}
and further under $SU(5)\to SU(3)_c\times SU(2)_L\times U(1)_Y$ (with $Q=T_3+Y$) the
content is exactly one Standard-Model family plus a singlet:
\begin{equation}\label{eq:sm:family}
\begin{aligned}
10 &\;\supset\; q=(3,2)_{+1/6},\quad u^c=(\bar3,1)_{-2/3},\quad e^c=(1,1)_{+1},\\
\overline{5} &\;\supset\; d^c=(\bar3,1)_{+1/3},\quad l=(1,2)_{-1/2},\\
1 &\;=\; \nu^c=(1,1)_{0}\ \ (\text{right-handed neutrino, $SU(5)$ singlet}).
\end{aligned}
\end{equation}
Thus $10+\overline{5}+1=16$ Weyl fermions $=(q,u^c,d^c,l,e^c)+\nu^c$: one SM family
plus a right-handed neutrino, with no mirror partners.
\end{theorem}

The chiral family sits inside $16_{+1}$ of the matter $27$ of $E_6$, whose canonical
UCFT branching (\autoref{sec:potential:branchings}) is
\begin{equation}\label{eq:sm:27branch}
27 \;=\; 1_{+4}\;+\;10_{-2}\;+\;16_{+1}\,,
\end{equation}
where the $1_{+4}$ is an extra SM-singlet modulus available as a right-handed-neutrino
modulus. That a single $16$ rather than a vector-like $16\oplus\overline{16}$ appears
per family is the content of the chirality postulate:

\begin{postulate}[Matter content and chirality, P4]\label{post:sm:P4}
Each family is one chiral $16$ of $\mathrm{Spin}(10)$ with the $U(1)$ charge of
\eqref{eq:sm:27branch}; there are no mirror partners.
\end{postulate}

\subsection{Three generations: \texorpdfstring{$SU(3)_F$}{SU(3)F} on
\texorpdfstring{$J_3(\mathbb{O})\otimes\mathbb{C}^3$}{J3(O)xC3}}
\label{sec:sm:generations}

\begin{postulate}[Family symmetry, P5]\label{post:sm:P5}
The matter carries an $SU(3)_F$ horizontal (family) symmetry acting on
$J_3(\mathbb{O})\otimes\mathbb{C}^3$; the three copies in the $\mathbb{C}^3$ factor are
three chiral $16$s of $\mathrm{Spin}(10)$, with no mirror partners. Three generations
are \emph{not} a consequence of the two axioms; they are added structure.
\end{postulate}

\begin{theorem}[Recomputed $n_F$ and three-generation $b_0$]\label{thm:sm:b0}
With one chiral $\mathbf{16}$ per family (P4, P5), the fermion count $n_F$ in the
one-loop $\mathrm{Spin}(10)$ $\beta$-function is the number of chiral
$\mathbf{16}$ \emph{representations} (one per generation), not the Weyl-fermion
dimension:
\begin{equation}\label{eq:sm:nF}
n_F \;=\; \#\text{generations} \;=\; 3
\qquad\text{(three chiral $\mathbf{16}$s, each weighted at $T(\mathbf{16})=2$).}
\end{equation}
The gauge $\beta$-function (canonical $+2$ term restored; asymptotically free, no
interacting fixed point) is
\begin{equation}\label{eq:sm:beta}
\beta_{\hat g^{2}} \;=\; 2\,\hat g^{2} \;+\; \frac{b_0}{48\pi^2}\,\hat g^{4},
\qquad
b_0 \;=\; \tfrac{11}{3}C_A \;-\; \tfrac{4}{3}T\,n_F \;-\; \tfrac{1}{6}T\,n_S\,,
\end{equation}
with the locked Casimirs $C_A(\mathrm{Spin}(10))=8$, $T(\mathbf{16})=2$,
$C_2(\mathbf{16})=45/8$. At the GUT matching scale with the full UCFT scalar
content, $n_S=28$ massive coset scalars enter the loop. Direct evaluation gives
\begin{equation}\label{eq:sm:b0three}
\boxed{\,
b_0(3\,\mathrm{gen})
=\tfrac{11}{3}(8)-\tfrac{4}{3}(2)(3)-\tfrac{1}{6}(2)(28)
=\tfrac{88-24-28}{3}=12\,.}
\end{equation}
The canonical one-generation anchor $b_0(1\,\mathrm{gen})=26$ uses $n_F=1$ and
$n_S=2$ (minimal GUT Higgs/Yukawa scalars at that scale):
\begin{equation}\label{eq:sm:b0one}
b_0(1\,\mathrm{gen})
=\tfrac{11}{3}(8)-\tfrac{4}{3}(2)(1)-\tfrac{1}{6}(2)(2)=\tfrac{78}{3}=26.
\end{equation}
We use $b_0=12$ for three-generation running below $M_{\mathrm{GUT}}$ and in the
AS stability check of \autoref{sec:gravity:stability}.
\end{theorem}

\begin{remark}
The gauge sector is asymptotically \emph{free}: the canonical $+2\,\hat g^{2}$
term in \eqref{eq:sm:beta} drives a Gaussian UV point $\hat g^{2}_\star=0$; there
is no interacting gauge fixed point. More fermions \emph{decrease} $b_0$ in
\eqref{eq:sm:beta}, so $b_0(3)<b_0(1)$. UV completeness is asymptotic freedom
plus the (indicative) gravitational fixed point $\kappa_\star=96\pi^2/5$
(\autoref{sec:gravity:AS}) plus irrelevant higher operators.
\end{remark}

\begin{postulate}[Hierarchy and mixing from flavon VEVs, P5 part II]\label{post:sm:flavons}
$SU(3)_F$ is broken sequentially ($SU(3)_F\to SU(2)_F\to \mathbf{1}$) by flavon VEVs,
generating the inter-generation Yukawa hierarchy and the CKM (quark) and PMNS (lepton)
mixing matrices as ratios of flavon VEVs to the family scale.
\end{postulate}

\begin{indicative}[Heterotic corroboration]\label{ind:sm:heterotic}
The heterotic $\mathbb{Z}_5$ quintic $\hat Q=Q/\mathbb{Z}_5$ compactification with the
tangent-bundle $SU(3)$ embedding gives the same generation count,
\begin{equation}\label{eq:sm:het}
h^1(\hat Q,V_{16})=3,\qquad h^1(\hat Q,V_{10})=14,\qquad h^1(\hat Q,\mathbf{1})=1,
\end{equation}
with the single bundle singlet an overall right-handed-neutrino modulus and the
$4+28=32$ heterotic moduli matching the $4$ clock-field Goldstone $+$ $28$ massive
coset scalars of UCFT. The low-energy central charges $(c_L,c_R)=(16,10)$ enter here.
This is a correspondence, not a duality, and not a derivation of the generation number.
\end{indicative}

\subsection{Yukawa and seesaw sector (P6)}
\label{sec:sm:yukawa}

\begin{theorem}[No $\{16,16bar\}$ Yukawa]\label{thm:sm:noyukawa}
The naive Yukawa $-y\,\bar\psi\,\phi^a T_a\,\psi$ with $\phi\in\{16,16bar\}$ is
forbidden: it is non-$H$-invariant and gives zero mass, because
\begin{equation}\label{eq:sm:tensor}
16\times 16 \;=\; 10+120+126\,,
\qquad
16\times \overline{16} \;=\; 1+45+210\,,
\end{equation}
so a $\{16,16bar\}$ scalar furnishes no $H$-singlet, and $H=\mathrm{Spin}(10)\times
U(1)$ is unbroken at the vacuum. The viable couplings use the real GUT-Higgs reps.
\end{theorem}

\begin{postulate}[Yukawa couplings, P6 part II]\label{post:sm:P6yukawa}
The family-indexed Yukawa Lagrangian is
\begin{align}
\mathcal{L}^{\mathrm{Dirac}}_{\mathrm{Yuk}}
 &\;=\; y^{(10)}_{ij}\,16_i\,16_j\,10_H \;+\; \mathrm{h.c.}\,,
 \label{eq:sm:Ldirac}\\
\mathcal{L}^{\mathrm{Maj}}_{\mathrm{Yuk}}
 &\;=\; y^{(126)}_{ij}\,16_i\,16_j\,\overline{126}_H \;+\; \mathrm{h.c.}
 \;\;\;\Big(\text{or}\;\; \tfrac{1}{M_*}\,16_i\,16_j\,\overline{16}_H\,\overline{16}_H\Big)\,,
 \label{eq:sm:Lmaj}
\end{align}
both $SO(10)$-invariant since $16\times 16\supset 10$ and $16\times 16\supset 126$,
with $i,j$ the $SU(3)_F$ family indices.
\end{postulate}

\begin{theorem}[Type-I seesaw]\label{thm:sm:seesaw}
After electroweak breaking $\langle 10_H\rangle=v=246$~GeV the Dirac masses are
$m_D\sim y^{(10)}v$. The SM-singlet, $B\!-\!L=2$ component of $126_H$ takes
$\langle\,\cdot\,\rangle=v_R\sim M_{\mathrm{GUT}}$ (step~2 of \eqref{eq:sm:chain}),
giving the right-handed Majorana mass $M_R\sim y^{(126)}v_R$ (equivalently
$M_R\sim v_R^2/M_*$ from the dim-$5$ operator). The $(\nu,\nu^c)$ mass matrix
\begin{equation}\label{eq:sm:matrix}
\mathcal{M}_\nu \;=\;
\begin{pmatrix} 0 & m_D \\ m_D^{\mathsf T} & M_R \end{pmatrix}
\end{equation}
yields, to leading order in $m_D/M_R$, the light-neutrino mass
\begin{equation}\label{eq:sm:mnu}
m_\nu \;\simeq\; \frac{m_D^2}{M_R} \;\sim\; 0.05\ \mathrm{eV}\,,
\end{equation}
the observed atmospheric scale for $m_D\sim v$ and $M_R\sim 10^{14}$--$10^{15}$~GeV.
PMNS mixing follows from the flavon structure (\autoref{post:sm:flavons}).
\end{theorem}

These multiplets are added physics; stating them is what makes the construction an
embedding rather than a derivation.

\subsection{Anomaly arithmetic}
\label{sec:sm:anomaly}

\begin{theorem}[Anomaly cancellation]\label{thm:sm:anomaly}
The composite $SO(10)\times U(1)$ gauge sector is anomaly-free. With Dynkin indices
$A_{16}=A_{10}=1$, $A_1=0$ and the $27$-charges of \eqref{eq:sm:27branch}, the mixed
$U(1)$--$SO(10)^2$ coefficient is
\begin{equation}\label{eq:sm:mixed}
(+1)A_{16}+(-2)A_{10}+(+4)A_1
 \;=\;(+1)(1)+(-2)(1)+(+4)(0)\;=\;-1\,,
\end{equation}
cancelled by the Green--Schwarz mechanism: the $E_6$-invariant $3$-form $\Omega_3$
pulls back to $dB=(1/2\pi)F_1$, and
$S_{\mathrm{GS}}=\int B\wedge \mathrm{Tr}(F\wedge F)/(8\pi^2)$ has a gauge variation
exactly cancelling the $-1$. The cubic $U(1)^3$ coefficient vanishes identically,
\begin{equation}\label{eq:sm:cubic}
16(+1)^3+10(-2)^3+1(+4)^3 \;=\; 16-80+64 \;=\; 0\,.
\end{equation}
Below $M_{\mathrm{GUT}}$ each family is the anomaly-free $16$ (the $\nu^c$ is exactly
what makes the SM hypercharge anomalies cancel family-by-family).
\end{theorem}

\begin{remark}[Strings and topology, correspondence]\label{rem:sm:strings}
The composite $SO(10)\times U(1)$ connection (Maurer--Cartan), the Green--Schwarz
cancellation, the vortex strings, and the heterotic worldsheet CFT with
$(c_L,c_R)=(16,10)$ are retained as a low-energy correspondence. Topology caveat:
$\mathrm{EIII}$ is simply connected, so $\pi_1(M)=\mathbb{Z}$ is \emph{false} of
$\mathrm{EIII}$ itself; the vortex strings arise from the gauged $U(1)$-quotient
configuration space (or are softened to a correspondence).
\end{remark}

\subsection{Phenomenology (indicative)}
\label{sec:sm:pheno}

\begin{indicative}[Proton decay, leptogenesis, monopole dilution]\label{ind:sm:pheno}
Once descent and seesaw exist, the leading observables are, all indicative:
\begin{equation}\label{eq:sm:proton}
\tau(p\to e^+\pi^0)\;\sim\;\frac{M_{\mathrm{GUT}}^4}{\alpha_{\mathrm{GUT}}^2\,m_p^5}\,,
\end{equation}
to be confronted with Super-/Hyper-Kamiokande, with the scalar-leptoquark channel
forbidden by matter parity / $B\!-\!L$; baryogenesis via leptogenesis from the heavy
$\nu^c$ decays (\autoref{thm:sm:seesaw}); and monopole dilution by inflation, with the
inflaton identified as one of the $28$ massive coset scalars
(\autoref{sec:gravity:spectrum}) and inflation occurring \emph{after} the GUT
transition. Here $M_{\mathrm{GUT}}$ is the breaking-VEV scale tied to $f$.
\end{indicative}

\subsection{Summary of results}
\label{sec:sm:summary}

\begin{itemize}
\item \textbf{Theorems (proved within UCFT):} the $E_6\to\mathrm{Spin}(10)\times U(1)$
vacuum selection (\autoref{thm:potential:vacuum}); the descent
(\autoref{thm:sm:descent}); the branchings of $27$ and $16$
(\autoref{thm:sm:family}); the anomaly arithmetic
(\autoref{thm:sm:anomaly}, mixed $=-1$ cancelled by GS, cubic $=0$); and the
three-generation $b_0=12$ (\autoref{thm:sm:b0}).
\item \textbf{Postulates (added structure):} P4 (one chiral $16$/family, no mirrors,
\autoref{post:sm:P4}); P5 ($SU(3)_F\to$ three generations, $n_F=3$,
\autoref{post:sm:P5}); P6 ($10_H$ Dirac, $126_H$/dim-$5$ Majorana, type-I seesaw,
\autoref{post:sm:P6higgs}, \autoref{post:sm:P6yukawa}).
\item \textbf{Correspondence (indicative):} heterotic $h^1(V_{16})=3$ and the
$4+28=32$ moduli match (\autoref{ind:sm:heterotic}); the phenomenology of
\autoref{ind:sm:pheno}.
\end{itemize}

\noindent\textbf{Conclusion.} Under postulates P4--P6, UCFT embeds the Standard
Model: three anomaly-free chiral families, standard hypercharge with GUT
normalization $\alpha_1=\tfrac53\alpha_Y$, electroweak breaking at $v=246$~GeV, and a
type-I seesaw $m_\nu\sim0.05$~eV.
% --- end inlined: unit-sm.tex ---
%----------------------------------------------------------------------%
% Part: Correspondences and Conclusions
%----------------------------------------------------------------------%
\part{Correspondences and Conclusions}
% --- begin inlined: unit-strings.tex ---
%======================================================================%
% UNIT: strings — Vortex Strings, Heterotic Correspondence, and Conclusions
%======================================================================%
\section{Vortex strings, heterotic correspondence, and conclusions}
\label{sec:strings:main}

We close the construction by collecting the four threads of the broken
phase --- topological strings, their worldsheet algebra, the heterotic
moduli match, and the Standard-Model descent --- into a single,
systematic classification of the principal results. Throughout we
keep the signature $(-,+,+,+)$, write $[f]=\text{mass}$ for the decay
constant, and call the $32$ coset Goldstone directions the \emph{clock
fields}. The vacuum manifold is the rank-$1$ primitive-idempotent orbit
\begin{equation}
  \mathcal{M} \;=\; E_{6}/(\mathrm{Spin}(10)\times U(1)) \;=\; \mathrm{EIII},
  \qquad \dim_{\mathbb{R}}\mathcal{M}=32,
  \label{eq:strings:vacuum}
\end{equation}
with positive-definite coset metric $K|_{\mathfrak m}$ and tangent module
$\mathfrak m = \mathbf{16}_{-3}\oplus\overline{\mathbf{16}}_{+3}$
(\autoref{sec:coset:main}). The physical spectrum is $28$ massive clock
scalars of a single common mass $m^{2}=\mu^{2}$ plus $4$ gauge/soldering
modes eaten by diffeomorphism and local Lorentz invariance
(\autoref{sec:gravity:main}).

%----------------------------------------------------------------------%
\subsection{The topological source of the strings}
\label{sec:strings:topology}

The naive expectation is that the broken phase supports cosmic strings
classified by $\pi_{1}(\mathcal{M})$. This must be stated with care.

\begin{theorem}[EIII is simply connected]
\label{thm:strings:pi1}
The vacuum manifold \eqref{eq:strings:vacuum} is a compact Hermitian
symmetric space (the complex Cayley plane). Hence
\[
   \pi_{1}(\mathrm{EIII})=0 .
\]
\end{theorem}

\begin{proof}
$\mathrm{EIII}=E_{6}/(\mathrm{Spin}(10)\times U(1))$ is an irreducible
compact Hermitian symmetric space of type $\mathrm{E\,III}$
(\'E.~Cartan's classification). The total space $E_{6}$ (compact, simply
connected) fibres over $\mathcal{M}$ with connected fibre
$\mathrm{Spin}(10)\times U(1)$, which is itself connected; the long exact
homotopy sequence
$\pi_{1}(E_{6})\to\pi_{1}(\mathcal{M})\to\pi_{0}(\mathrm{Spin}(10)\times
U(1))$ has both flanking groups trivial, so $\pi_{1}(\mathcal{M})=0$.
Equivalently, all compact irreducible Hermitian symmetric spaces are
simply connected.
\end{proof}

\begin{remark}[Topological consistency]
\label{rem:strings:guard}
\autoref{thm:strings:pi1} \emph{forbids} the assertion
$\pi_{1}(\mathcal{M})=\mathbb{Z}$ \emph{of EIII itself}: there are no
$\sigma$-model strings winding the bare coset target. Any source claiming
$\pi_{1}(\mathrm{EIII})=\mathbb{Z}$ is in error. The correct topological
seat of the strings is the \emph{gauged} configuration space, identified
next.
\end{remark}

The strings instead arise once the composite $\mathrm{Spin}(10)\times
U(1)$ connection (\autoref{sec:gauge:main}) is switched on. Gauging the
$U(1)\subset H$ that acts on the clock fields replaces the target by the
configuration space of the gauged $\sigma$-model. Concretely, the
physical order parameter for vortices is the $U(1)$-charged combination
of clock fields whose phase lives in
\[
   \mathcal{C}_{U(1)} \;\simeq\; \mathcal{M}\,/\!/\,U(1),
   \qquad
   \pi_{1}\bigl(\mathcal{C}_{U(1)}\bigr)\;\supseteq\;\mathbb{Z},
   \label{eq:strings:config}
\]
the winding being the standard Abelian-Higgs flux of the gauged $U(1)$
(the same $U(1)$ whose Stueckelberg/Green--Schwarz coupling appears in
\autoref{sec:anomaly:main}). Equivalently, on the broken vacuum the
$U(1)$ factor of $H=\mathrm{Spin}(10)\times U(1)$ admits a magnetic flux
tube: $\pi_{1}(U(1))=\mathbb{Z}$ labels the quantised flux, and the clock
field winds the eaten phase around the core.

\begin{postulate}[Topological string sector]
\label{post:strings:sector}
The relevant first homotopy is that of the $U(1)$-gauged configuration
space \eqref{eq:strings:config}, $\pi_{1}(\mathcal{C}_{U(1)})\supseteq
\mathbb{Z}$, \emph{not} that of the bare coset. The vortex strings are
flux tubes of the composite $U(1)$ (cf.\ Axiom~I and the gauging
prescription, \autoref{sec:gauge:main}); we treat the string content as a
\emph{correspondence} with Abelian-Higgs/heterotic vortices rather than a
$\pi_{1}(\mathcal{M})$ theorem.
\end{postulate}

%----------------------------------------------------------------------%
\subsection{BPS vortex strings of finite tension}
\label{sec:strings:bps}

Choose cylindrical coordinates $(r,\theta,z,t)$ on the emergent
Lorentzian manifold $(X,\eta)$ of \autoref{sec:soldering:main}. A static
configuration of the $U(1)$-gauged clock-field sector with winding
$n\in\mathbb{Z}$ obeys
\begin{equation}
  \phi(r,\theta+2\pi)=e^{\,n\,T_{U(1)}\,2\pi}\,\phi(r,\theta),
  \qquad
  \phi(0)=X_{0}\in\mathcal{M},\quad
  D_{r}\phi\xrightarrow[r\to\infty]{}0,
  \label{eq:strings:ansatz}
\end{equation}
with $T_{U(1)}$ the broken $U(1)$ generator and $X_{0}$ a rank-$1$
idempotent representative of the vacuum.

\begin{theorem}[Finite BPS tension]
\label{thm:strings:tension}
For each $n\in\mathbb{Z}$ the Euler--Lagrange equations of the gauged
clock-field action admit a static, axially symmetric, finite-tension
solution of the ansatz type \eqref{eq:strings:ansatz}. Its tension
saturates a Bogomol'nyi bound,
\begin{equation}
  \mathcal{T}_{n}
   = 2\pi f^{2}\,|n|\,
     \int_{0}^{\infty}\!dr\;
     r\,K_{\alpha\beta}(\phi)\,
     D_{r}\phi^{\alpha}\,D_{r}\phi^{\beta}
   \;=\; 2\pi f^{2}\,|n|\,\Omega,
  \qquad 0<\mathcal{T}_{n}<\infty,
  \label{eq:strings:tension}
\end{equation}
where $\Omega$ is the $U(1)$ flux quantum through the core and
$K_{\alpha\beta}$ the positive-definite coset metric.
\end{theorem}

\begin{proof}[Proof sketch]
Completing the static energy functional into a sum of squares plus a
topological flux term gives a first-order Bogomol'nyi system whose
solutions minimise the energy in each winding class; the
positive-definiteness of $K|_{\mathfrak m}$ (\autoref{sec:coset:main})
guarantees a strictly positive integrand, and the exponential approach
$D_{r}\phi\to0$ (set by the common mass $m^{2}=\mu^{2}$ of the $28$
massive clock fields) makes the integral converge. The detailed
moduli-matrix construction is deferred to the worldsheet appendix,
\autoref{app:ws:main}.
\end{proof}

Linear fluctuations about the BPS background factorise into right-moving
Nambu--Goto modes (transverse string position) and left-moving internal
currents (the eaten $\mathrm{Spin}(10)\times U(1)$ phases of the core).

%----------------------------------------------------------------------%
\subsection{Worldsheet algebra and the heterotic correspondence}
\label{sec:strings:worldsheet}

\begin{theorem}[Left-moving affine algebra]
\label{thm:strings:affine}
Quantisation of the gauge zero modes trapped on the vortex core yields a
left-moving affine current algebra at level one,
\begin{equation}
  \widehat{\mathfrak{so}}(10)_{1}\;\oplus\;\widehat{\mathfrak u}(1)_{1}.
  \label{eq:strings:affine}
\end{equation}
\end{theorem}

\begin{proof}[Proof sketch]
The Green--Schwarz term $S_{\mathrm{GS}}=\int B\wedge
\mathrm{Tr}(F\wedge F)/8\pi^{2}$ (\autoref{sec:anomaly:main}) reduces on
the two-dimensional string core $\Sigma\simeq\mathbb{R}^{1,1}$ to a
chiral WZW action whose level equals the mixed
$U(1)$--$\mathrm{Spin}(10)^{2}$ anomaly coefficient computed there;
matching that coefficient fixes $k=k'=1$. The Sugawara stress tensor of
\eqref{eq:strings:affine} then governs the left-movers. Full details
appear in \autoref{app:ws:main}.
\end{proof}

\begin{theorem}[Sugawara central charge]
\label{thm:strings:cl}
The Sugawara central charge of \eqref{eq:strings:affine} is
\begin{equation}
  c_{L}
   = \frac{\dim\mathrm{SO}(10)\cdot k}{k+h^{\vee}_{\mathrm{SO}(10)}}
     + 1
   = \frac{45\cdot 1}{1+8} + 1
   = 5 + 1
   = \boxed{16},
  \label{eq:strings:cL}
\end{equation}
using $\dim\mathrm{SO}(10)=45$, dual Coxeter number
$h^{\vee}_{\mathrm{SO}(10)}=8$, and $c[\widehat{\mathfrak u}(1)_{1}]=1$.
\end{theorem}

\begin{proof}
The Sugawara formula for $\widehat{\mathfrak g}_{k}$ gives
$c=k\dim\mathfrak g/(k+h^{\vee})$; at $k=1$ for
$\mathrm{SO}(10)$ this is $45/9=5$. The Abelian factor contributes $1$.
Their sum is $16$.
\end{proof}

\begin{theorem}[Heterotic worldsheet correspondence]
\label{thm:strings:het}
Combining the four right-moving spacetime bosons ($c_{R}=4$ Nambu--Goto)
with their $(0,1)$ worldsheet superpartners ($c_{R}=4\times\tfrac12=2$)
and the transverse fermions of the core gives a right-moving central
charge $c_{R}=10$; together with $c_{L}=16$ from \eqref{eq:strings:cL},
the vortex worldsheet realises the central-charge content
\begin{equation}
  (c_{L},c_{R})=(16,10)
  \label{eq:strings:cc}
\end{equation}
of the level-one left-moving sector of the $E_{8}\times E_{8}$ heterotic
string.
\end{theorem}

\begin{remark}[Correspondence, not duality]
\label{rem:strings:corr}
\autoref{thm:strings:het} is a \emph{low-energy correspondence}: the
gauge sector, anomaly structure, and central charges of the UCFT vortex
match the heterotic worldsheet at leading order in $\alpha'$. It is
\emph{not} a strict duality --- no full equivalence of S-matrices,
spectra, or quantum moduli spaces is claimed, and the matching is stated
only at the level of \eqref{eq:strings:cc} and the affine algebra
\eqref{eq:strings:affine}. The conformal-field-theory derivation
(Sugawara construction, modular invariance, higher-derivative
robustness) is deferred to \autoref{app:ws:main}.
\end{remark}

%----------------------------------------------------------------------%
\subsection{The $\mathbb{Z}_{5}$ quintic Calabi--Yau correspondence}
\label{sec:strings:cy}

The broken phase leaves $32 = 4+28$ real flat directions: the $4$
gauge/soldering clock fields $\pi_{-}^{i}$ ($i=0,1,2,3$) and the $28$
massive clock fields $\pi_{+}^{j}$ promoted to Wilson-line moduli once
the composite connection condenses,
\begin{equation}
  \mathcal{M}^{\mathrm{mod}}_{\mathrm{UCFT}}
   = \underbrace{\mathbb{R}^{4}/\Gamma^{4}}_{\text{Goldstone/soldering}}
     \times
     \underbrace{\mathbb{R}^{28}/\Gamma^{28}}_{\text{Wilson lines}},
   \qquad
   \dim_{\mathbb{R}}\mathcal{M}^{\mathrm{mod}}_{\mathrm{UCFT}}=32.
  \label{eq:strings:moduliU}
\end{equation}

On the heterotic side, compactify the $E_{8}\times E_{8}$ string on the
freely-acting $\mathbb{Z}_{5}$ quotient of the Fermat quintic,
$\widehat{Q}=Q/\mathbb{Z}_{5}$ with
$Q=\{\sum_{a=0}^{4}z_{a}^{5}=0\}\subset\mathbb{CP}^{4}$, standard
embedding of $T\widehat{Q}$ in one $E_{8}$:
\begin{equation}
  h^{1,1}(\widehat{Q})=1,\quad
  h^{2,1}(\widehat{Q})=0,\quad
  \pi_{1}(\widehat{Q})=\mathbb{Z}_{5},
  \label{eq:strings:hodge}
\end{equation}
the Wilson lines along $\pi_{1}(\widehat{Q})=\mathbb{Z}_{5}$ breaking
$E_{6}\to\mathrm{Spin}(10)\times U(1)$. Bundle cohomology gives
\begin{equation}
  h^{1}(\widehat{Q},V_{\mathbf{16}})=3,\qquad
  h^{1}(\widehat{Q},V_{\mathbf{10}})=14,\qquad
  h^{1}(\widehat{Q},\mathbf{1})=1,
  \label{eq:strings:bundle}
\end{equation}
hence the real-moduli count
\begin{equation}
  \underbrace{4}_{\text{K\"ahler}+B_{2}\text{ axions}}
  \;+\;
  \underbrace{28}_{2\times h^{1}(V_{\mathbf{10}})}
  \;=\;32
  \;=\;\dim_{\mathbb{R}}\mathcal{M}^{\mathrm{mod}}_{\mathrm{Het}} .
  \label{eq:strings:count}
\end{equation}

\begin{theorem}[Moduli matching]
\label{thm:strings:moduli}
The real dimensions and gauge-bundle topology of
$\mathcal{M}^{\mathrm{mod}}_{\mathrm{UCFT}}$ in
\eqref{eq:strings:moduliU} and of
$\mathcal{M}^{\mathrm{mod}}_{\mathrm{Het}}$ in \eqref{eq:strings:count}
agree: both are $32 = 4+28$, both carry an
$\mathrm{Spin}(10)\times U(1)$ bundle, and the kinetic metrics coincide
at tree level (each descends from the $E_{6}$ Killing form). The four
Goldstone/soldering coordinates map to the universal K\"ahler modulus
plus three axionic partners; the $28$ Wilson-line coordinates map to the
$2\times 14$ real bundle moduli in $H^{1}(\widehat{Q},V_{\mathbf{10}})$.
\end{theorem}

\begin{proof}[Proof sketch]
Counting and bundle topology are read off from
\eqref{eq:strings:moduliU}--\eqref{eq:strings:count}. Metric equality
follows from the common $E_{6}$ Killing origin of both kinetic terms;
Chern-class/Chern--Simons matching uses the shared Green--Schwarz
two-form. The detailed cohomology and the explicit map are given in
\autoref{app:cy:main}.
\end{proof}

\begin{remark}[Three generations: corroboration, not derivation]
\label{rem:strings:gen}
The chiral generation count of the heterotic model is
$h^{1}(\widehat{Q},V_{\mathbf{16}})=3$ from \eqref{eq:strings:bundle}.
This \emph{corroborates} the three generations of UCFT, which are
supplied as an added postulate $P5$ --- the horizontal $SU(3)_{F}$
symmetry on $J_{3}(\mathbb{O})\otimes\mathbb{C}^{3}$
(\autoref{sec:sm:main}) --- and \emph{not} as a second, independent
mechanism. The $14$ vector-like $\mathbf{10}$ moduli and the single
$\mathbf{1}$ (a right-handed-neutrino modulus) complete the
$32=4+28$ match. Generations are \emph{added structure} ($P5$), with the
heterotic count as a \emph{correspondence}.
\end{remark}

%----------------------------------------------------------------------%
\subsection{Conclusions: theorems versus postulates}
\label{sec:strings:conclusions}

UCFT is an $E_{6}$/Albert-algebra--motivated unification. Two axioms
supply the arena (the $27=J_{3}(\mathbb{O})$ order parameter and its
quantum kinematics); a small, explicitly enumerated set of postulates
$P1$--$P6$ then yields a four-dimensional Lorentzian quantum field theory
containing an anomaly-free $\mathrm{Spin}(10)\times U(1)$ gauge sector,
Sakharov-induced gravity, and --- after a stated GUT-Higgs/flavour sector
--- the Standard Model with three generations and neutrino masses. This
is a \emph{conditional embedding} of the Standard Model and gravitation.
\autoref{tab:strings:ledger} records the classification of results.

\begin{table}[h]
\centering
\renewcommand{\arraystretch}{1.25}
\begin{tabular}{lll}
\hline
\textbf{Result} & \textbf{Content} & \textbf{Status}\\
\hline
Vacuum selection ($P1\Rightarrow P2$)
  & global min $=$ rank-$1$ idempotent orbit $\mathrm{EIII}$, $\dim 32$
  & \textbf{Theorem}\\
Origin instability
  & $\mathrm{Hess}\,V(0)\prec0$; roll-down to $\mathcal{M}$
  & Theorem\\
Coset tangent / metric
  & $\mathfrak m=\mathbf{16}_{-3}\oplus\overline{\mathbf{16}}_{+3}$, $K|_{\mathfrak m}\succ0$
  & Theorem\\
Soldering algebra
  & $H_{2}(\mathbb{O})\cong\mathbb{R}^{1,9}$, $SL(2,\mathbb{O})=\mathrm{Spin}(1,9)$
  & Theorem\\
Signature $(1,3)$, $d=4$
  & added soldering structure
  & \textbf{Postulate $P3$}\\
Mass spectrum
  & $28$ massive (common $m^{2}=\mu^{2}$) $+\,4$ eaten
  & Theorem\\
Induced gravity
  & $M_{\mathrm{Pl}}^{2}=\tfrac{28}{6}f^{2}\mu^{2}\log(\Lambda^{2}/\mu^{2})/16\pi^{2}>0$
  & Theorem\\
Gauge running
  & asymptotically free, $b_{0}(1\,\mathrm{gen})=26$, $b_{0}(3\,\mathrm{gen})=12$
  & Theorem\\
Gravitational fixed point
  & $\kappa_{\ast}=96\pi^{2}/5$, $\mathrm{eig}=-2$
  & \textbf{Indicative}\\
SM descent
  & $\mathbf{16}=$ one family $+\,\nu^{c}$; anomalies cancel (GS)
  & Theorem\\
Chirality / generations / Yukawa
  & $P4,P5,P6$; seesaw $m_{\nu}\sim0.05$\,eV
  & \textbf{Postulate}\\
BPS vortex strings
  & finite tension $\mathcal{T}_{n}=2\pi f^{2}|n|\,\Omega$
  & Theorem\\
String topology
  & from gauged $U(1)$ config space, not $\pi_{1}(\mathrm{EIII})=0$
  & Theorem $+$ \textbf{Postulate $P$}\\
Heterotic worldsheet
  & $(c_{L},c_{R})=(16,10)$, $\widehat{\mathfrak{so}}(10)_{1}\oplus\widehat{\mathfrak u}(1)_{1}$
  & \textbf{Correspondence}\\
Calabi--Yau moduli
  & $32=4+28$, $h^{1}(V_{\mathbf{16}})=3$
  & \textbf{Correspondence}\\
Cosmological constant
  & tuned relevant $\Lambda_{cc}$
  & inherited, not solved\\
\hline
\end{tabular}
\caption{Classification of results by epistemic status. Theorems are
proved from stated hypotheses; postulates ($P1$--$P6$) are added
structure; conjectural results are motivated but outside controlled
weak coupling; correspondences are leading-order matches, not
dualities.}
\label{tab:strings:ledger}
\end{table}

\paragraph{What is proved.}
The vacuum-selection theorem ($P1\Rightarrow P2$, the rank-$1$
idempotent orbit $\mathcal{M}=\mathrm{EIII}$ of real dimension $32$), the
positive-definiteness and uniqueness of the coset metric, the octonionic
soldering algebra $SL(2,\mathbb{O})\cong\mathrm{Spin}(1,9)$, the common
radiative mass of the $28$ clock scalars by Schur's lemma, the positivity
$M_{\mathrm{Pl}}^{2}>0$ of the Sakharov-induced Planck mass with
prefactor $28/6$, the asymptotic freedom of the gauge--Yukawa sector
($\beta_{\hat g^{2}}=2\hat g^{2}+(b_{0}/48\pi^{2})\hat g^{4}$,
$b_{0}=12$ for three generations), the Green--Schwarz cancellation of the
mixed anomaly, the embedding of one Standard-Model family plus a
right-handed neutrino in the $\mathbf{16}$, and the finiteness of the BPS
vortex tension are all \emph{theorems}.

\paragraph{What is postulated.}
The Lorentzian signature and dimension ($P3$), the chiral matter content
without mirrors ($P4$), the three generations via $SU(3)_{F}$ ($P5$), and
the real GUT-Higgs/Yukawa sector generating the type-I seesaw ($P6$) are
\emph{added postulates}. The topological string sector rests on the
gauged-$U(1)$ configuration space (\autoref{post:strings:sector}), since
$\pi_{1}(\mathrm{EIII})=0$.

\paragraph{What is indicative or a correspondence.}
The gravitational fixed point $\kappa_{\ast}=96\pi^{2}/5\approx189.50$ is
\emph{indicative}: $\kappa_{\ast}\gg1$ lies outside controlled weak
coupling. Reflection positivity of the UCFT $\sigma$-model (and hence
Osterwalder--Schrader reconstruction of Lorentzian dynamics) is
indicative. The heterotic worldsheet match $(c_{L},c_{R})=(16,10)$ and
the $\mathbb{Z}_{5}$ quintic moduli match $32=4+28$ with
$h^{1}(V_{\mathbf{16}})=3$ are \emph{low-energy correspondences}, not
strict dualities. UCFT \emph{inherits, and does not solve}, the
cosmological-constant problem: $\Lambda_{cc}$ is a tuned relevant
coupling.

\paragraph{Outlook.}
Several directions would sharpen the embedding into a derivation. (i) A
curved-background functional-RG truncation could test whether the
indicative gravitational fixed point survives beyond the single-scalar
loop, and whether the higher-curvature operators are genuinely
irrelevant. (ii) A direct check of reflection positivity for the gauged
clock-field $\sigma$-model would upgrade emergent unitary time from
indicative to theorem. (iii) Deriving $d=4$ and signature $(1,3)$ from a
dynamical vierbein-condensate selection, rather than postulating $P3$,
remains open. (iv) A worldsheet computation of the full
vortex partition function would test the heterotic correspondence beyond
central charges. (v) Phenomenologically, dimension-six proton decay
$\tau(p\to e^{+}\pi^{0})\sim M_{\mathrm{GUT}}^{4}/(\alpha_{\mathrm{GUT}}^{2}
m_{p}^{5})$, leptogenesis from the heavy $\nu^{c}$, and inflation driven
by one of the $28$ massive clock fields after the GUT transition all
furnish falsifiable tests. UCFT is presented as a conditional embedding:
the proved core---one Albert-algebra order parameter selecting an
anomaly-free $\mathrm{Spin}(10)\times U(1)$ gauge sector with induced,
positive Planck mass---is established under the stated hypotheses, while the
Lorentzian, flavour, and gravitational-UV completions rest on explicit
postulates and conjectures enumerated above.
% --- end inlined: unit-strings.tex ---

\appendix
\part{Appendices}

% --- begin inlined: unit-app-albert.tex ---
%======================================================================%
\section{Octonions and the Albert Algebra}
\label{sec:app-albert:main}
%======================================================================%

This appendix collects, in a single self-contained place, the octonionic and
Jordan-algebraic facts on which the rest of the manuscript rests: the octonion
multiplication table and structure constants; the exceptional Jordan (Albert)
algebra $J_3(\mathbb O)$ as $3\times 3$ octonion-Hermitian matrices; the Jordan
product, the trace bilinear form, the cubic norm $N$, and the sharp map $X^\#$
with the Freudenthal identities; the rank stratification and primitive
idempotents; and the actions of $F_4$ and $E_6$. Every statement is tagged as a
\emph{theorem} (proved here or from cited standard references), a
\emph{postulate} (assumed added structure, cf.\ P1--P6), or as
\emph{indicative}. The identities assembled here are exactly those invoked by the
potential, mass-spectrum and gauge sections; we therefore cross-reference them by
the stable labels defined below.

Throughout we use the metric signature $(-,+,+,+)$ and the manuscript-wide
conventions: $[f]=\text{mass}$, and $\hat g^2,\xi,\hat y^2,\kappa$ dimensionless.
Two real forms of $E_6$ appear and must \emph{not} be conflated: the
\emph{compact} $E_6$ (equivalently $E_{6(-14)}$, maximal compact
$\mathrm{Spin}(10)\times U(1)$) governs the dynamics and gives a
positive-definite coset metric, while the non-compact $E_{6(-26)}$ is used
\emph{only} as the rigid classifier of orbits of the real $27$. This distinction
is the foundation enforced in \autoref{sec:app-albert:invariants} and
\autoref{sec:app-albert:orbits}, and underlies the vacuum-selection theorem cited
as \autoref{thm:potential:T4}.

%%%%%%%%%%%%%%%%%%%%%%%%%%%%%%%%%%%%%%%%%%%%%%%%%%%%%%%%%%%%%%%%%%%%%%
\subsection{Octonions: multiplication and structure constants}
\label{sec:app-albert:octonions}
%%%%%%%%%%%%%%%%%%%%%%%%%%%%%%%%%%%%%%%%%%%%%%%%%%%%%%%%%%%%%%%%%%%%%%

\begin{definition}[Octonions]\label{def:app-albert:octonions}
The octonions $\mathbb O$ form the real, $8$-dimensional, unital,
\emph{non-associative} division algebra. Fix the standard basis
$\{e_0=1,\,e_1,\dots,e_7\}$ with $e_i^2=-1$ for $i\ge 1$ and $e_i e_j=-e_j e_i$
for $i\neq j$ ($i,j\ge 1$). The multiplication is encoded by the totally
antisymmetric structure constants $\psi_{ijk}$,
\begin{equation}\label{eq:app-albert:struct}
  e_i\, e_j \;=\; -\,\delta_{ij}\,e_0 \;+\; \psi_{ijk}\,e_k ,
  \qquad i,j,k\in\{1,\dots,7\},
\end{equation}
with $\psi_{ijk}=+1$ on the seven Fano lines
\begin{equation}\label{eq:app-albert:fano}
  (ijk)\in\bigl\{(124),(235),(346),(457),(561),(672),(713)\bigr\}
\end{equation}
(and cyclic), and $\psi_{ijk}=0$ otherwise. Conjugation is
$\bar e_0=e_0$, $\bar e_i=-e_i$ ($i\ge1$); the (positive-definite) norm is
$n(x)=x\bar x=\bar x x=\sum_{a=0}^{7}x_a^2$, and $\overline{xy}=\bar y\,\bar x$.
\end{definition}

\begin{theorem}[Composition, alternativity, Moufang]\label{thm:app-albert:moufang}
\textup{(Theorem; standard \cite{Baez:2002,Springer:2000}.)}
$\mathbb O$ is a normed composition algebra, $n(xy)=n(x)\,n(y)$. It is not
associative but is \emph{alternative}: the associator
$[x,y,z]:=(xy)z-x(yz)$ is totally antisymmetric, so any subalgebra generated by
two elements is associative. The Moufang identities hold, in particular
$(xy)(zx)=x\bigl((yz)x\bigr)$. Real and imaginary parts obey
$\Re(xy)=\Re(yx)$ and $\Re\bigl((xy)z\bigr)=\Re\bigl(x(yz)\bigr)$, so the
trilinear form $\Re(xyz)$ is well defined (no bracketing needed) and cyclic.
\end{theorem}

\begin{proof}
Composition follows because $n$ is the norm of a Cayley--Dickson double of the
quaternions; alternativity and the antisymmetry of $\psi_{ijk}$ are equivalent to
the total antisymmetry of the associator, and the Moufang identities are the
standard consequences. The reality identities follow from
$x+\bar x=2\Re(x)\,e_0\in\mathbb R$ together with $\overline{xy}=\bar y\bar x$.
See \cite{Baez:2002,Schafer:1966}. $\qquad\blacksquare$
\end{proof}

\begin{remark}[$G_2$ and $\mathrm{Spin}(7)$]\label{rmk:app-albert:g2}
\textup{(Indicative context.)}
The automorphism group $\mathrm{Aut}(\mathbb O)$ is the compact $G_2$ (dim $14$),
the stabilizer of $\psi_{ijk}$ viewed as a generic $3$-form on $\Im\mathbb O\cong
\mathbb R^7$; the spin group $\mathrm{Spin}(7)$ acts by triality on the three
$8$-dimensional modules $\mathbf 8_v,\mathbf 8_s,\mathbf 8_c$. These reappear when
$\mathrm{Spin}(10)$ spinors are realized on $J_3(\mathbb O)$
(\autoref{sec:app-albert:F4E6}) and feed the coset tangent module
$\mathfrak m=16_{-3}\oplus\overline{16}_{+3}$ used downstream.
\end{remark}

%%%%%%%%%%%%%%%%%%%%%%%%%%%%%%%%%%%%%%%%%%%%%%%%%%%%%%%%%%%%%%%%%%%%%%
\subsection{The Albert algebra \texorpdfstring{$J_3(\mathbb O)$}{J3(O)}}
\label{sec:app-albert:jalg}
%%%%%%%%%%%%%%%%%%%%%%%%%%%%%%%%%%%%%%%%%%%%%%%%%%%%%%%%%%%%%%%%%%%%%%

\begin{definition}[Albert algebra]\label{def:app-albert:jalg}
The (real, reduced) \emph{exceptional Jordan algebra} or \emph{Albert algebra} is
the space of $3\times3$ octonion-Hermitian matrices
\begin{equation}\label{eq:app-albert:Xform}
  J_3(\mathbb O)=
  \left\{\,
  X=\begin{pmatrix}
      \xi_1 & x_3 & \bar x_2\\
      \bar x_3 & \xi_2 & x_1\\
      x_2 & \bar x_1 & \xi_3
    \end{pmatrix}
  :\ \xi_i\in\mathbb R,\ x_i\in\mathbb O
  \,\right\},
\end{equation}
equipped with the commutative \emph{Jordan product}
\begin{equation}\label{eq:app-albert:jprod}
  X\circ Y=\tfrac12\,(XY+YX),
\end{equation}
where $XY$ is ordinary (octonionic) matrix multiplication.
\end{definition}

\begin{theorem}[Dimension and Jordan identity]\label{thm:app-albert:dim}
\textup{(Theorem.)}
$J_3(\mathbb O)$ is a $27$-dimensional, commutative, power-associative Jordan
algebra: $\circ$ is commutative and satisfies the Jordan identity
$X\circ\bigl(Y\circ X^{\circ 2}\bigr)=(X\circ Y)\circ X^{\circ 2}$. It is
\emph{formally real} (Euclidean): $\sum_i A_i^{\circ2}=0\Rightarrow A_i=0$. It is
the unique exceptional Jordan algebra (not isomorphic to a subalgebra of an
associative algebra under $\circ$).
\end{theorem}

\begin{proof}
The real count is $3$ (diagonal) $+\,3\times 8$ (off-diagonal octonions)
$=27$. Commutativity is immediate from \eqref{eq:app-albert:jprod}. The Jordan
identity holds for $3\times3$ Hermitian matrices over an alternative algebra by
the Hermitian-matrix theorem of Jordan--von Neumann--Wigner
\cite{JNW:1934,McCrimmon:2004}; formal reality follows from positivity of the
trace form \eqref{eq:app-albert:Q} below. Exceptionality (the Albert
exceptionality theorem) is standard \cite{Albert:1934,Schafer:1966}.
$\qquad\blacksquare$
\end{proof}

\begin{definition}[Traces, trace form, cubic norm]\label{def:app-albert:invariants}
On $J_3(\mathbb O)$ define the \emph{linear trace}
$\operatorname{tr}(X)=\xi_1+\xi_2+\xi_3$, the symmetric \emph{trace bilinear form}
\begin{equation}\label{eq:app-albert:Q}
  \langle X,Y\rangle:=\operatorname{tr}(X\circ Y),
  \qquad
  Q(X):=\langle X,X\rangle=\operatorname{tr}(X^{\circ2})
  =\sum_i\xi_i^2+2\sum_i n(x_i)\ \ge 0,
\end{equation}
which is positive-definite, and the \emph{cubic norm} (determinant analogue)
\begin{equation}\label{eq:app-albert:N}
  N(X):=\xi_1\xi_2\xi_3-\sum_{i=1}^{3}\xi_i\,n(x_i)+2\,\Re\!\bigl(x_1 x_2 x_3\bigr).
\end{equation}
The second elementary symmetric invariant is
$\sigma_2(X)=\tfrac12\bigl(\operatorname{tr}(X)^2-Q(X)\bigr)$. The trilinear
polarization of $N$ is denoted $N(X,Y,Z)$ with $N(X,X,X)=6N(X)$.
\end{definition}

\begin{theorem}[Cubic Cayley--Hamilton]\label{thm:app-albert:CH}
\textup{(Theorem.)}
Every $X\in J_3(\mathbb O)$ satisfies the cubic characteristic identity
\begin{equation}\label{eq:app-albert:CH}
  X^{\circ3}-\operatorname{tr}(X)\,X^{\circ2}+\sigma_2(X)\,X-N(X)\,E=0,
  \qquad E:=\mathrm{diag}(1,1,1).
\end{equation}
Its three roots are the real \emph{Jordan eigenvalues}
$\lambda_1,\lambda_2,\lambda_3$, in terms of which
$\operatorname{tr}(X)=\sum_i\lambda_i$, $Q(X)=\sum_i\lambda_i^2$,
$N(X)=\lambda_1\lambda_2\lambda_3$.
\end{theorem}

\begin{proof}
By the spectral theorem for formally real Jordan algebras every $X$ is
$F_4$-conjugate to $\mathrm{diag}(\lambda_1,\lambda_2,\lambda_3)$ with
$\lambda_i\in\mathbb R$ (\autoref{thm:app-albert:F4} below); on diagonal elements
\eqref{eq:app-albert:CH} is the ordinary scalar Cayley--Hamilton identity, and
$F_4$-equivariance of every term ($\operatorname{tr},\sigma_2,N,E$ are
$F_4$-invariant or invariant) transports it to all of $J_3(\mathbb O)$
\cite{McCrimmon:2004,Springer:2000}. $\qquad\blacksquare$
\end{proof}

%%%%%%%%%%%%%%%%%%%%%%%%%%%%%%%%%%%%%%%%%%%%%%%%%%%%%%%%%%%%%%%%%%%%%%
\subsection{The sharp map and the Freudenthal identities}
\label{sec:app-albert:sharp}
%%%%%%%%%%%%%%%%%%%%%%%%%%%%%%%%%%%%%%%%%%%%%%%%%%%%%%%%%%%%%%%%%%%%%%

\begin{definition}[Sharp map and cross product]\label{def:app-albert:sharp}
The \emph{sharp} (Freudenthal adjoint) map $X\mapsto X^\#$ is the unique quadratic
map with entries the $2\times2$ octonionic cofactors of $X$:
\begin{equation}\label{eq:app-albert:sharpentries}
  (X^\#)_{ii}=\xi_j\xi_k-n(x_i),
  \qquad
  (X^\#)_{ij}=\overline{x_i x_j}-\xi_k\,x_k
  \quad(\text{cyclic } i,j,k).
\end{equation}
Equivalently, in the Jordan frame
$\mathrm{diag}(\lambda_1,\lambda_2,\lambda_3)^\#
 =\mathrm{diag}(\lambda_2\lambda_3,\lambda_3\lambda_1,\lambda_1\lambda_2)$. Its
symmetric polarization defines the \emph{Freudenthal cross product}
\begin{equation}\label{eq:app-albert:cross}
  X\times Y:=(X+Y)^\#-X^\#-Y^\#,
  \qquad X\times X=2X^\#,
\end{equation}
a symmetric bilinear map $J_3(\mathbb O)\times J_3(\mathbb O)\to J_3(\mathbb O)$.
\end{definition}

\begin{theorem}[Freudenthal / adjoint identities]\label{thm:app-albert:freudenthal}
\textup{(Theorem; \cite{Freudenthal:1985,McCrimmon:2004,Springer:2000}.)}
For all $X,Y\in J_3(\mathbb O)$ the following hold:
\begin{align}
  X^\#\circ X &= N(X)\,E,
  \label{eq:app-albert:F1}\\
  (X^\#)^\# &= N(X)\,X,
  \label{eq:app-albert:F2}\\
  \langle X^\#,Y\rangle &= DN(X)[Y]
     = N(X,X,Y),
  \quad\text{i.e. } X^\#=\nabla N(X),
  \label{eq:app-albert:F3}\\
  \langle X^\#,Y\rangle &= \sigma_2(X)\,\langle?\rangle
     \ \Longrightarrow\
  \operatorname{tr}(X^\#)=\sigma_2(X),
  \label{eq:app-albert:F4}\\
  \langle X^\#,X^\#\rangle
     &= \sigma_2(X)^2-2\,\operatorname{tr}(X)\,N(X)
     = (\lambda_2\lambda_3)^2+(\lambda_3\lambda_1)^2+(\lambda_1\lambda_2)^2.
  \label{eq:app-albert:F5}
\end{align}
\end{theorem}

\begin{proof}
Identity \eqref{eq:app-albert:F1} is the cofactor expansion of the determinant
$N$ and follows from \eqref{eq:app-albert:CH} by writing
$X^\#=X^{\circ2}-\operatorname{tr}(X)X+\sigma_2(X)E$ and using
\eqref{eq:app-albert:CH}; \eqref{eq:app-albert:F2} (the ``$\#\#$'' Freudenthal
identity) is verified in the Jordan frame, where
$(\lambda_2\lambda_3,\lambda_3\lambda_1,\lambda_1\lambda_2)^\#
=(\lambda_1^2\lambda_2\lambda_3,\dots)=N\cdot(\lambda_1,\lambda_2,\lambda_3)$, and
extended by $F_4$-equivariance. The gradient identity \eqref{eq:app-albert:F3} is
the definition of $X^\#$ as the differential of the cubic $N$ with respect to the
trace form, and $\operatorname{tr}(X^\#)=\sigma_2(X)$ in
\eqref{eq:app-albert:F4} follows from
$X^\#=X^{\circ2}-\operatorname{tr}(X)X+\sigma_2(X)E$ on taking $\operatorname{tr}$.
Finally \eqref{eq:app-albert:F5} is the frame evaluation of
$\langle X^\#,X^\#\rangle$, manifestly a sum of squares; the identity
$\sigma_2(X)^2-2\operatorname{tr}(X)N(X)$ is the symmetric-function rewriting.
$\qquad\blacksquare$
\end{proof}

\begin{corollary}[Rank from the sharp map]\label{cor:app-albert:rank}
\textup{(Theorem.)}
Define the \emph{rank} by
\begin{equation}\label{eq:app-albert:rank}
  \operatorname{rank}(X)=
  \begin{cases}
    0,& X=0,\\
    1,& X\neq 0,\ X^\#=0,\\
    2,& X^\#\neq0,\ N(X)=0,\\
    3,& N(X)\neq0.
  \end{cases}
\end{equation}
Then $X^\#=0\iff\operatorname{rank}(X)\le1$, and $N(X)=0\iff\operatorname{rank}(X)
\le2$. In particular the cubic norm $N$, not the quadratic $Q$, detects rank
deficiency at order $3$.
\end{corollary}

\begin{proof}
In the Jordan frame $X^\#=0$ iff all pairwise products $\lambda_i\lambda_j$
vanish, i.e.\ at most one $\lambda_i\neq0$; and $N=\lambda_1\lambda_2\lambda_3=0$
iff some $\lambda_i=0$. Equivariance extends both statements to all of
$J_3(\mathbb O)$. $\qquad\blacksquare$
\end{proof}

\begin{remark}[Positivity of the sharp invariant]\label{rmk:app-albert:positivity}
\textup{(Theorem.)}
Because the trace form \eqref{eq:app-albert:Q} is positive-definite,
$\langle X^\#,X^\#\rangle\ge0$ with equality iff $X^\#=0$ iff
$\operatorname{rank}(X)\le1$ (\autoref{cor:app-albert:rank}). This is the exact
algebraic input to the potential: the $E_6$-covariant term
$\alpha\langle X^\#,X^\#\rangle$ vanishes precisely on the rank-$\le1$ cone, which
is why it selects the rank-$1$ idempotent orbit; see
\autoref{thm:potential:T1} and \autoref{thm:potential:T4}.
\end{remark}

%%%%%%%%%%%%%%%%%%%%%%%%%%%%%%%%%%%%%%%%%%%%%%%%%%%%%%%%%%%%%%%%%%%%%%
\subsection{Rank stratification and primitive idempotents}
\label{sec:app-albert:idempotents}
%%%%%%%%%%%%%%%%%%%%%%%%%%%%%%%%%%%%%%%%%%%%%%%%%%%%%%%%%%%%%%%%%%%%%%

\begin{definition}[Idempotents and Jordan frame]\label{def:app-albert:idemp}
$P\in J_3(\mathbb O)$ is an \emph{idempotent} if $P\circ P=P$, and
\emph{primitive} if it is a nonzero idempotent that is not a sum of two nonzero
orthogonal idempotents. A \emph{complete Jordan frame} is a triple
$\{P_1,P_2,P_3\}$ of pairwise orthogonal ($P_i\circ P_j=0$, $i\neq j$) primitive
idempotents with $P_1+P_2+P_3=E$.
\end{definition}

\begin{theorem}[Idempotent stratification]\label{thm:app-albert:idemp}
\textup{(Theorem; \cite{McCrimmon:2004,Yokota:2009}.)}
A primitive idempotent has Jordan eigenvalues $(1,0,0)$, hence
$\operatorname{tr}(P)=1$, $Q(P)=1$, $N(P)=0$, $P^\#=0$, i.e.\
$\operatorname{rank}(P)=1$. The normalized rank-$1$ elements are exactly the
primitive idempotents; they form a single $F_4$-orbit
\begin{equation}\label{eq:app-albert:OP2}
  F_4/\mathrm{Spin}(9)\ \cong\ \mathbb{OP}^2,
  \qquad \dim_{\mathbb R}=16,
\end{equation}
the octonionic projective plane (the real rank-$1$ idempotents). Under the full
\emph{compact} $E_6$ the projective class $[X]$ of a rank-$1$ line sweeps the
larger homogeneous space
\begin{equation}\label{eq:app-albert:EIII}
  \mathcal M=\mathrm{EIII}=E_6/(\mathrm{Spin}(10)\times U(1)),
  \qquad \dim_{\mathbb R}\mathcal M=32,
\end{equation}
the complex Cayley plane and the closed minimal $E_6$-orbit in
$\mathbb P(27_{\mathbb C})$.
\end{theorem}

\begin{proof}
$P\circ P=P$ with eigenvalues $\mu$ forces $\mu^2=\mu$, so $\mu\in\{0,1\}$;
primitivity forbids two $1$'s, and nonzero forbids all $0$'s, leaving $(1,0,0)$.
The remaining invariants follow from \autoref{thm:app-albert:CH} and
\eqref{eq:app-albert:sharpentries}. Transitivity of $F_4=\mathrm{Aut}(J_3(\mathbb
O))$ on primitive idempotents with stabilizer $\mathrm{Spin}(9)$ is the standard
identification $\mathbb{OP}^2=F_4/\mathrm{Spin}(9)$
\cite{Baez:2002,Yokota:2009}; the $E_6$-orbit statement is the realization of
EIII as the minimal/highest-weight orbit in the projectivized $27$
\cite{Springer:2000}. $\qquad\blacksquare$
\end{proof}

\begin{remark}[Vacuum manifold; excluded base point]\label{rmk:app-albert:vacuum}
\textup{(Theorem, consistency note.)}
The UCFT vacuum manifold is exactly \eqref{eq:app-albert:EIII}, the rank-$1$
primitive-idempotent orbit, $\dim_{\mathbb R}=32$ --- \emph{not} the orbit of
$E=\mathrm{diag}(1,1,1)$. The latter has rank $3$, $E^\#=E\neq0$,
$\langle E^\#,E^\#\rangle=3>0$, stabilizer $F_4$, and orbit $E_6/F_4$ of real
dimension $26$; by \autoref{rmk:app-albert:positivity} it is strictly disfavoured
by the potential and is therefore excluded. This is consistent with the
vacuum-selection theorem \autoref{thm:potential:T4} and the coset tangent
$\mathfrak m=16_{-3}\oplus\overline{16}_{+3}$ of \autoref{sec:app-albert:branch}.
\end{remark}

%%%%%%%%%%%%%%%%%%%%%%%%%%%%%%%%%%%%%%%%%%%%%%%%%%%%%%%%%%%%%%%%%%%%%%
\subsection{Invariants and the two real forms}
\label{sec:app-albert:invariants}
%%%%%%%%%%%%%%%%%%%%%%%%%%%%%%%%%%%%%%%%%%%%%%%%%%%%%%%%%%%%%%%%%%%%%%

\begin{theorem}[$F_4=\mathrm{Aut}(J_3(\mathbb O))$]\label{thm:app-albert:F4}
\textup{(Theorem; \cite{Schafer:1966,Yokota:2009}.)}
The automorphism group of the Jordan algebra (preserving the product $\circ$,
hence the unit $E$ and \emph{both} $Q$ and $N$) is the compact exceptional group
$F_4$, $\dim_{\mathbb R}F_4=52$. It acts on $J_3(\mathbb O)=\mathbb R\oplus 26$
fixing the unit; on the traceless $26$ it embeds $F_4\subset SO(26)\subset
SO(27)$, since it preserves the Euclidean form $Q$.
\end{theorem}

\begin{theorem}[Reduced structure group $=E_6$]\label{thm:app-albert:E6}
\textup{(Theorem; \cite{Jacobson:1968,Springer:2000}.)}
The \emph{reduced structure group} of $J_3(\mathbb O)$ --- the connected group of
linear maps preserving the cubic norm $N$ up to a real character --- is the
exceptional group $E_6$, $\dim_{\mathbb R}E_6=78$, acting irreducibly on the $27$.
For $g$ in this group,
\begin{equation}\label{eq:app-albert:character}
  N(g\cdot X)=\lambda(g)\,N(X),\qquad \lambda(g)\in\mathbb R^\times,
\end{equation}
and the sharp map is covariant, $(g\cdot X)^\#=\lambda(g)\,(g^{-1})^{\mathsf T}\!
\cdot X^\#$ in the trace form. The split/non-compact real form realized by the
\emph{division}-octonion algebra is $E_{6(-26)}$; it does \emph{not} preserve
$Q=\operatorname{tr}(X^2)$.
\end{theorem}

\begin{proof}
Jacobson's construction realizes $\mathfrak e_6$ as $\mathfrak f_4$ (derivations)
plus the traceless multiplication operators
$L_X:Y\mapsto X\circ Y-\tfrac13\langle X,Y\rangle E$ acting on the $27$; the
group generated preserves $N$ up to the character \eqref{eq:app-albert:character}
\cite{Jacobson:1968}. The frame element $\mathrm{diag}(s,s,s^{-2})$ rescales $N$
nontrivially while changing $Q$, so $Q$ is not $E_6$-invariant; only $F_4$
(the $\lambda\equiv1$, unit-preserving subgroup) preserves $Q$. $\qquad
\blacksquare$
\end{proof}

\begin{postulate}[Compact real form for dynamics]\label{post:app-albert:compact}
\textup{(Convention, cf.\ P1--P2.)}
For the dynamics we use the \emph{compact} real form $E_6$ (equivalently
$E_{6(-14)}$, whose maximal compact subgroup is exactly
$H=\mathrm{Spin}(10)\times U(1)$). On the compact form the $27$ is unitary, the
trace form $Q=\langle\cdot,\cdot\rangle$ is the invariant positive-definite
metric, and the norm character is identically $\lambda\equiv1$ (a continuous
homomorphism from a compact connected group into $\mathbb R_{>0}$). Hence on the
compact form every $E_6$-invariant polynomial built from $\langle\cdot,\cdot
\rangle$ and $N$ --- in particular $\langle X^\#,X^\#\rangle$ --- is a genuine
invariant, the coset metric is positive-definite, and the gauge sector is unitary.
The non-compact $E_{6(-26)}$ is retained \emph{only} as the rigid orbit classifier
of \autoref{sec:app-albert:orbits}.
\end{postulate}

\begin{remark}[Only $N$ separates structure-group orbits]\label{rmk:app-albert:onlyN}
\textup{(Theorem, guard.)}
$Q=\operatorname{tr}(X^2)$ is $F_4$-invariant \emph{only}; it is \emph{not}
preserved by the $E_6$ structure group, and the pair $(Q,N)$ does \emph{not}
classify $E_6$-orbits. Only the cubic norm $N$ is structure-group invariant (up to
the character $\lambda$), and orbits are classified by rank and $N$
(\autoref{sec:app-albert:orbits}). Any claim that the structure group preserves
$Q$, or that $(Q,N)$ separates orbits, is false.
\end{remark}

%%%%%%%%%%%%%%%%%%%%%%%%%%%%%%%%%%%%%%%%%%%%%%%%%%%%%%%%%%%%%%%%%%%%%%
\subsection{Orbit classification of the real \texorpdfstring{$27$}{27}}
\label{sec:app-albert:orbits}
%%%%%%%%%%%%%%%%%%%%%%%%%%%%%%%%%%%%%%%%%%%%%%%%%%%%%%%%%%%%%%%%%%%%%%

\begin{theorem}[Rank-and-norm orbit classification]\label{thm:app-albert:orbits}
\textup{(Theorem; \cite{Springer:2000,Krutelevich:2007}.)}
The orbits of $E_{6(-26)}$ on the real $27=J_3(\mathbb O)$ are classified by the
rank \eqref{eq:app-albert:rank} together with the value of the cubic norm $N$:
\begin{itemize}
  \item rank $0$: $\{0\}$;
  \item rank $1$: $X\neq0$, $X^\#=0$ --- a single orbit, the cone over EIII; its
    projectivization is the closed minimal orbit $\mathcal M=\mathrm{EIII}$,
    $\dim_{\mathbb R}\mathcal M=32$;
  \item rank $2$: $X^\#\neq0$, $N(X)=0$ --- a single orbit;
  \item rank $3$: $N(X)\neq0$ --- orbits labelled by $\operatorname{sgn}N$ and the
    signature of $X$ (the indefinite cubic $N$ is genuinely sign-indefinite).
\end{itemize}
The quadratic $Q$ is \emph{not} an orbit invariant; it can be rescaled within a
fixed-$N$, fixed-rank orbit by the non-compact directions
$\mathrm{diag}(s,s,s^{-2})\in E_{6(-26)}$.
\end{theorem}

\begin{proof}
That rank and $N$ are structure-group invariants follows from
\eqref{eq:app-albert:character} and the covariance of $X^\#$
(\autoref{thm:app-albert:E6}); transitivity within each stratum is the
reduction of an arbitrary $X$ to a normal form
$\mathrm{diag}(\lambda_1,\lambda_2,\lambda_3)$ by $F_4$ followed by a
norm-rescaling diagonal element, with the rank-deficient strata as the explicit
degenerations \cite{Krutelevich:2007,Springer:2000}. $\qquad\blacksquare$
\end{proof}

\begin{remark}[Sharp characterization, summary]\label{rmk:app-albert:sharpsummary}
\textup{(Theorem.)}
Collecting \autoref{cor:app-albert:rank} and
\autoref{thm:app-albert:freudenthal}: $\operatorname{rank}\le1\Leftrightarrow
X^\#=0$; $\operatorname{rank}\le2\Leftrightarrow N(X)=0$;
$\operatorname{rank}=3\Leftrightarrow N(X)\neq0$. These are the identities the
potential section uses to localize the global minimum on the rank-$1$ cone
(\autoref{thm:potential:T4}).
\end{remark}

%%%%%%%%%%%%%%%%%%%%%%%%%%%%%%%%%%%%%%%%%%%%%%%%%%%%%%%%%%%%%%%%%%%%%%
\subsection{$F_4$- and $H$-branchings; the coset tangent module}
\label{sec:app-albert:branch}
%%%%%%%%%%%%%%%%%%%%%%%%%%%%%%%%%%%%%%%%%%%%%%%%%%%%%%%%%%%%%%%%%%%%%%

\begin{theorem}[$\mathrm{Spin}(10)$ realization]\label{thm:app-albert:F4E6}
\textup{(Theorem; \cite{Baez:2002,Springer:2000}.)}
Under $H=\mathrm{Spin}(10)\times U(1)\subset E_6$ (the maximal compact of
$E_{6(-14)}$), the fundamental and adjoint representations branch as
\begin{align}
  27 &= 1_{+4}\oplus 10_{-2}\oplus 16_{+1},
  \label{eq:app-albert:27branch}\\
  78 &= 45_{0}\oplus 1_{0}\oplus 16_{-3}\oplus \overline{16}_{+3},
  \label{eq:app-albert:78branch}
\end{align}
where the subscripts are the $U(1)$ charges in the normalization fixed by
\eqref{eq:app-albert:78branch}. The $10_{-2}$ is the $\mathrm{Spin}(10)$ vector
(an $H_2(\mathbb O)$-type block of $J_3(\mathbb O)$), the $16_{+1}$ a chiral
spinor, and the $1_{+4}$ a Jordan-frame singlet (an extra Standard-Model-singlet
modulus).
\end{theorem}

\begin{corollary}[Coset tangent module]\label{cor:app-albert:tangent}
\textup{(Theorem.)}
The tangent space to $\mathcal M=\mathrm{EIII}=E_6/(\mathrm{Spin}(10)\times U(1))$
at the basepoint is the $H$-module
\begin{equation}\label{eq:app-albert:tangent}
  \mathfrak m=16_{-3}\oplus\overline{16}_{+3},
  \qquad \dim_{\mathbb R}\mathfrak m=32,
\end{equation}
the complement of $\mathfrak h=45_0\oplus1_0$ in
\eqref{eq:app-albert:78branch}. It is \emph{real-irreducible of complex type}
(commutant $\cong\mathbb C$) and contains \emph{no} $\mathrm{Spin}(10)$ singlet.
Consequently the invariant coset metric $K|_{\mathfrak m}$ is unique up to scale
and \emph{positive-definite} (Schur), and any radiative ($H$-invariant)
mass operator on $\mathfrak m$ is a single scalar: one common mass on all $32$
clock fields.
\end{corollary}

\begin{proof}
$\mathfrak m=\mathfrak e_6\ominus\mathfrak h$ reads off
\eqref{eq:app-albert:78branch}. The $U(1)$ charges $\pm3$ forbid the symmetric
invariants in $16\otimes16$ and $\overline{16}\otimes\overline{16}$, leaving the
unique charge-$0$ cross-pairing $16_{-3}\otimes\overline{16}_{+3}\to1_0$; this is
the restriction of the Killing form, which on the compact form is
positive-definite. Real-irreducibility with commutant $\mathbb C$ is the
complex-type Schur statement, giving uniqueness of the metric and a single common
radiative mass. Absence of a singlet is immediate from
\eqref{eq:app-albert:tangent}. $\qquad\blacksquare$
\end{proof}

\begin{remark}[Excluded tangent decompositions]\label{rmk:app-albert:forbidden}
\textup{(Consistency constraint.)}
The coset tangent is \eqref{eq:app-albert:tangent} and contains no
$\mathrm{Spin}(10)$ singlet. Any decomposition of $\mathfrak m$ as $2\times10+4$,
or $10+\overline{10}+4$, or any singlet-in-$\mathfrak m$ form, is incorrect, as is
any attempt to extract a $4$-dimensional $H$-invariant subspace from
$\nabla\sigma_2,\nabla N$ singlets. The Lorentzian frame directions are
\emph{not} a sub-$H$-module of $\mathfrak m$; they are introduced as added
structure (soldering postulate P3, cf.\ \autoref{sec:app-soldering:main}) and are
gauge-protected, not Hessian zeros.
\end{remark}

%%%%%%%%%%%%%%%%%%%%%%%%%%%%%%%%%%%%%%%%%%%%%%%%%%%%%%%%%%%%%%%%%%%%%%
\subsection{The Minkowski block from the cubic norm}
\label{sec:app-albert:minkowski}
%%%%%%%%%%%%%%%%%%%%%%%%%%%%%%%%%%%%%%%%%%%%%%%%%%%%%%%%%%%%%%%%%%%%%%

\begin{theorem}[$H_2(\mathbb O)\cong\mathbb R^{1,9}$]\label{thm:app-albert:minkowski}
\textup{(Theorem; \cite{Baez:2002,Sudbery:1984}.)}
The $2\times2$ octonion-Hermitian matrices
\begin{equation}\label{eq:app-albert:H2O}
  H_2(\mathbb O)=
  \left\{\begin{pmatrix}a & x\\ \bar x & b\end{pmatrix}:a,b\in\mathbb R,\,x\in
  \mathbb O\right\},
  \qquad \dim_{\mathbb R}=10,
\end{equation}
form a sub-block of $J_3(\mathbb O)$ (set the third pivot to $1$). With
$\det X=ab-n(x)$ and the lightcone parametrization $a=x^0+x^9$, $b=x^0-x^9$,
\begin{equation}\label{eq:app-albert:detmink}
  \det X=(x^0)^2-(x^9)^2-\textstyle\sum_{i=1}^{8}(x^i)^2
  =-\,\eta_{\mu\nu}X^\mu X^\nu,
  \quad \eta=\mathrm{diag}(-1,+1,\dots,+1),
\end{equation}
so $(H_2(\mathbb O),-\det)\cong\mathbb R^{1,9}$. Restricting the octonion entry to
$\mathbb C\subset\mathbb O$ gives $H_2(\mathbb C)\cong\mathbb R^{1,3}$ with
$-\det$ the $(1,3)$ Minkowski form in the convention $(-,+,+,+)$.
\end{theorem}

\begin{proof}
$\det\!\bigl(\begin{smallmatrix}a&x\\\bar x&b\end{smallmatrix}\bigr)
=ab-x\bar x=ab-n(x)$ by definition of $n$; substituting the lightcone
coordinates yields \eqref{eq:app-albert:detmink}. The restriction to
$\mathbb C$ drops $x^4,\dots,x^9$ (the imaginary octonion directions beyond
$\mathbb C$), leaving the four coordinates $x^0,x^1,x^2,x^3$ with the same form.
$\qquad\blacksquare$
\end{proof}

\begin{remark}[Where Lorentzian signature enters]\label{rmk:app-albert:signature}
\textup{(Postulate P3.)}
The positive-definite object is the coset/target metric $K|_{\mathfrak m}$ on
$\mathfrak m$ (\autoref{cor:app-albert:tangent}); the indefinite object is the
cubic norm $N$, whose $H_2(\mathbb O)\cong\mathbb R^{1,9}$ block carries the
Minkowski form \eqref{eq:app-albert:detmink}. These are two distinct forms on two
distinct modules. The Killing form on $\mathfrak m$ for the compact $E_6$ has
signature $(32,0)$ and \emph{zero} Lorentzian content; the spacetime dimension
$d=4$ and signature $(1,3)=(-,+,+,+)$ are \emph{added structure} (soldering
postulate P3, cf.\ \autoref{sec:app-soldering:main}), soldered from $N$ via
$H_2(\mathbb C)$, and are \emph{not} Killing-form, coset-geometry, or rank-lemma
theorems. Any $(4,28)$ or $(1,3)$ coset/Killing-signature claim is false.
\end{remark}

%%%%%%%%%%%%%%%%%%%%%%%%%%%%%%%%%%%%%%%%%%%%%%%%%%%%%%%%%%%%%%%%%%%%%%
\subsection{Identities used downstream}
\label{sec:app-albert:identities}
%%%%%%%%%%%%%%%%%%%%%%%%%%%%%%%%%%%%%%%%%%%%%%%%%%%%%%%%%%%%%%%%%%%%%%

For convenience we list the identities invoked by the potential, mass-spectrum and
gauge sections, all proved above.

\begin{enumerate}
  \item \textbf{Sharp / Freudenthal} \eqref{eq:app-albert:F1}--\eqref{eq:app-albert:F5}:
    $X^\#\circ X=N(X)E$, $(X^\#)^\#=N(X)X$, $X^\#=\nabla N(X)$,
    $\operatorname{tr}(X^\#)=\sigma_2(X)$, and
    $\langle X^\#,X^\#\rangle=\sigma_2(X)^2-2\operatorname{tr}(X)N(X)\ge0$. Used in
    the $E_6$-covariant term of the potential (P1) and in
    \autoref{thm:potential:T1}.
  \item \textbf{Rank stratification} \eqref{eq:app-albert:rank}: rank $\le1
    \Leftrightarrow X^\#=0$; rank $\le2\Leftrightarrow N=0$. Used to localize the
    global minimum on the rank-$1$ cone (\autoref{thm:potential:T4}).
  \item \textbf{Cubic Cayley--Hamilton} \eqref{eq:app-albert:CH} and frame
    eigenvalues. Used to evaluate $V$ on the cone.
  \item \textbf{Branchings} \eqref{eq:app-albert:27branch},
    \eqref{eq:app-albert:78branch} and the tangent module
    \eqref{eq:app-albert:tangent}. Used for the $27\times27$ scalar block
    $M=\mathrm{diag}(\mu_1\mathbb 1_1,\mu_{10}\mathbb 1_{10},\mu_{16}\mathbb 1_{16})$,
    for the single common coset mass (Schur on $\mathfrak m$), and for the gauge
    multiplet assignments. The $28$ massive clock fields
    ($32=28\text{ massive}+4\text{ gauge/soldering}$) enter all loop sums; cf.\
    \autoref{thm:gravity:MPl} and the gauge $\beta$-function
    \autoref{thm:gravity:beta}.
  \item \textbf{Minkowski block} \eqref{eq:app-albert:detmink}: $-\det|_{H_2
    (\mathbb O)}$ is the $\mathbb R^{1,9}$ form, $-\det|_{H_2(\mathbb C)}$ the
    $(1,3)$ form. Used by the soldering postulate P3
    (\autoref{sec:app-soldering:main}).
  \item \textbf{Casimirs of the gauge representations.} With the unbroken
    $H=\mathrm{Spin}(10)\times U(1)$ acting on the multiplets in
    \eqref{eq:app-albert:27branch}, the relevant invariants are
    $C_A(\mathrm{SO}(10))=8$, $T(16)=2$ (the spinor is index $2$), and
    $C_2(16)=45/8$. These feed
    $b_0=\tfrac{11}{3}C_A-\tfrac43 T n_F-\tfrac16 T n_S$ in
    \autoref{thm:gravity:beta} (canonical anchor $b_0=26$ for one generation,
    $b_0=12$ for three).
\end{enumerate}

\begin{remark}[Status]\label{rmk:app-albert:status}
All results in this appendix are \emph{theorems} of octonionic/Jordan algebra
(\autoref{thm:app-albert:moufang}--\autoref{thm:app-albert:minkowski}) or standard
representation theory, except the two clearly marked conventions/postulates: the
choice of the \emph{compact} real form for the dynamics
(\autoref{post:app-albert:compact}) and the soldering of Lorentzian signature from
$N$ (\autoref{rmk:app-albert:signature}, postulate P3). The algebra fixes the
arena; the physics input is the potential class (P1) and the soldering (P3),
treated in their own units.
\end{remark}
% --- end inlined: unit-app-albert.tex ---
% --- begin inlined: unit-app-lorentz.tex ---
\section{Minkowski Structure and the Soldering Map}
\label{sec:app-lorentz:main}

This appendix establishes, with explicit mathematics and consistent epistemic labels, the chain
by which a Lorentzian four-manifold is \emph{soldered} onto the
$E_{6}$/Albert-algebra arena. We prove the algebraic facts
$H_{2}(\mathbb{O})\cong\mathbb{R}^{1,9}$ and $SL(2,\mathbb{O})\cong\mathrm{Spin}(1,9)$
as theorems; we state the choice of a four-plane, the dimension $d=4$, and the
signature $(1,3)$ as the soldering postulate (P3); and we give the
Osterwalder--Schrader/Page--Wootters route to Lorentzian unitary dynamics in place
of any ``$e^{0}\mapsto i\,e^{0}$'' Wick trick.

The governing logical separation, enforced throughout, is the following. The
\emph{positive-definite} object is the coset/target metric $K|_{\mathfrak{m}}$ on
the $32$-dimensional internal tangent space
$\mathfrak{m}=\mathbf{16}_{-3}\oplus\overline{\mathbf{16}}_{+3}$ of
$M=E_{6}/(\mathrm{Spin}(10)\times U(1))=\mathrm{EIII}$ (see
\autoref{cor:VacuumManifold} and \autoref{thm:uniqueMetric}). The \emph{indefinite}
object is the cubic norm $N$ on the matter representation $\mathbf{27}=J_{3}(\mathbb{O})$,
whose $2\times2$ octonion-Hermitian sub-block carries a Minkowski form on the
\emph{different} space $\mathbb{R}^{1,9}$. No single form is asked to be both
definite and indefinite; the soldering map (P3) is the bridge importing the
indefinite $(1,3)$ structure onto a four-plane of the world-tangent bundle, leaving
$K|_{\mathfrak{m}}$ untouched and Euclidean. Conventions are fixed globally:
signature $(-,+,+,+)$; the dynamics live on the \emph{compact} real form of $E_{6}$
(maximal compact $H=\mathrm{Spin}(10)\times U(1)$), and $E_{6(-26)}$ enters only as
the rigid orbit-classifier of the real $\mathbf{27}$.

\subsection{The Minkowski block: \texorpdfstring{$H_{2}(\mathbb{O})\cong\mathbb{R}^{1,9}$}{H2(O) = R(1,9)}}
\label{subsec:app-lorentz:h2o}

Let $\mathbb{O}$ be the real division octonions, a non-associative composition
algebra of real dimension $8$, with conjugation $x\mapsto\bar{x}$, norm
$n(x)=x\bar{x}=\bar{x}x\in\mathbb{R}_{\ge0}$, and trace $t(x)=x+\bar{x}\in\mathbb{R}$.
Define the space of $2\times2$ octonion-Hermitian matrices
\begin{equation}
\label{eq:app-lorentz:h2def}
H_{2}(\mathbb{O})
 =\left\{\,X=\begin{pmatrix} a & x \\ \bar{x} & b \end{pmatrix}
   \;:\; a,b\in\mathbb{R},\ x\in\mathbb{O}\,\right\},
 \qquad X^{\dagger}=X,
\end{equation}
so that as a real vector space
$\dim_{\mathbb{R}}H_{2}(\mathbb{O})=1+1+8=10$. The determinant is well defined
despite non-associativity (the off-diagonal entry and its conjugate associate with
the scalar diagonal):
\begin{equation}
\label{eq:app-lorentz:det}
\det X = ab-x\bar{x}=ab-n(x).
\end{equation}

\begin{theorem}[Minkowski identification; \texorpdfstring{$H_{2}(\mathbb{O})\cong\mathbb{R}^{1,9}$}{H2(O)=R(1,9)}]
\label{thm:app-lorentz:Mink}
Introduce light-cone-adapted coordinates $a=X^{0}+X^{9}$, $b=X^{0}-X^{9}$, and write
$x=(X^{1},\dots,X^{8})\in\mathbb{R}^{8}$ in an orthonormal basis of $\mathbb{O}$, so
that $n(x)=\sum_{i=1}^{8}(X^{i})^{2}$. Then
\begin{equation}
\label{eq:app-lorentz:detmink}
\det X = (X^{0})^{2}-(X^{9})^{2}-\sum_{i=1}^{8}(X^{i})^{2}
       = -\,\eta_{\mu\nu}\,X^{\mu}X^{\nu},
\qquad
\eta=\mathrm{diag}(-1,\underbrace{+1,\dots,+1}_{9}),
\end{equation}
so the quadratic form $-\det$ on the ten-dimensional space $H_{2}(\mathbb{O})$ has
Lorentzian signature $(1,9)$ in the $(-,+,\dots,+)$ convention:
$(H_{2}(\mathbb{O}),-\det)\cong\mathbb{R}^{1,9}$. \emph{[Theorem.]}
\end{theorem}

\begin{proof}
By \eqref{eq:app-lorentz:det},
$\det X=ab-n(x)=(X^{0}+X^{9})(X^{0}-X^{9})-\sum_{i}(X^{i})^{2}
=(X^{0})^{2}-(X^{9})^{2}-\sum_{i}(X^{i})^{2}$, which is the displayed Minkowski
form. By Sylvester's law of inertia the diagonal pair contributes one $+$ (the
$X^{0}$ mode) and one $-$ (the $X^{9}$ mode), and $-n(x)$ contributes eight further
$-$'s; hence $\det$ has signature $(1,9)$ and $-\det$ is the single-timelike
$\mathbb{R}^{1,9}$.
\end{proof}

\begin{theorem}[Octonionic exceptional isomorphism; \texorpdfstring{$SL(2,\mathbb{O})\cong\mathrm{Spin}(1,9)$}{SL(2,O)=Spin(1,9)}]
\label{thm:app-lorentz:SL2O}
There is a Lie-group isomorphism $SL(2,\mathbb{O})\cong\mathrm{Spin}(1,9)$, under
which the defining $10$-dimensional vector representation acts on
$(H_{2}(\mathbb{O}),-\det)\cong\mathbb{R}^{1,9}$ and the two chiral
$16$-dimensional spinor representations are realised on
$\mathbb{O}^{2}\cong\mathbb{R}^{16}$. This is the octonionic ($n=8$) member of the
classical family $SL(2,\mathbb{A})\cong\mathrm{Spin}(n+1,1)$ for the normed
division algebras $\mathbb{A}=\mathbb{R},\mathbb{C},\mathbb{H},\mathbb{O}$ of
dimension $n=1,2,4,8$:
\begin{equation}
\label{eq:app-lorentz:family}
\begin{aligned}
SL(2,\mathbb{R}) &\cong \mathrm{Spin}(2,1), &
SL(2,\mathbb{C}) &\cong \mathrm{Spin}(3,1), \\
SL(2,\mathbb{H}) &\cong \mathrm{Spin}(5,1), &
SL(2,\mathbb{O}) &\cong \mathrm{Spin}(9,1),
\end{aligned}
\end{equation}
with $\mathrm{Spin}(9,1)=\mathrm{Spin}(1,9)$ up to convention. \emph{[Theorem.]}
\end{theorem}

\begin{proof}[Proof (statement and status)]
The set of determinant- and Hermiticity-preserving actions on
$H_{2}(\mathbb{O})$ does not close under naive matrix multiplication, since
octonion non-associativity obstructs $g\,X\,g^{\dagger}$. Accordingly
$SL(2,\mathbb{O})$ is \emph{defined} as $\mathrm{Spin}(1,9)$ acting through
left/right octonion-multiplication generators chosen so that the determinant
\eqref{eq:app-lorentz:det} and the spinor bilinears $\psi^{\dagger}X\psi$ on
$\psi\in\mathbb{O}^{2}$ are covariant; the double cover
$\mathrm{Spin}(1,9)\to SO(1,9)$ is the statement that the spinor action is
two-valued on the vector action of \autoref{thm:app-lorentz:Mink}. These
isomorphisms are classical (Sudbery; Kugo--Townsend; Manogue--Schray; Baez).
\end{proof}

\begin{theorem}[The physical Lorentz four-plane; \texorpdfstring{$SL(2,\mathbb{C})\subset SL(2,\mathbb{O})$}{SL(2,C) in SL(2,O)}]
\label{thm:app-lorentz:SL2C}
The chain of composition-algebra inclusions
$\mathbb{R}\subset\mathbb{C}\subset\mathbb{H}\subset\mathbb{O}$ induces a chain of
spin inclusions; in particular
\begin{equation}
\label{eq:app-lorentz:c-in-o}
\mathrm{Spin}(1,3)=SL(2,\mathbb{C})\ \subset\ SL(2,\mathbb{O})=\mathrm{Spin}(1,9)
\end{equation}
via the sub-block $H_{2}(\mathbb{C})\subset H_{2}(\mathbb{O})$, where
$H_{2}(\mathbb{C})=\{[[a,z],[\bar{z},b]]:a,b\in\mathbb{R},\,z\in\mathbb{C}\}$ has
real dimension $4$, $\det=ab-|z|^{2}$ is the $(1,3)$ Minkowski form, and
$\det|_{H_{2}(\mathbb{C})}$ is $SL(2,\mathbb{C})$-covariant in the standard
four-dimensional way. The ambient $\mathbb{R}^{1,9}$ is the indefinite arena; the
\emph{physical} $(1,3)$ is a soldered four-plane inside it. \emph{[Theorem.]}
\end{theorem}

\begin{proof}
Restrict the change of variables of \autoref{thm:app-lorentz:Mink} to
$z\in\mathbb{C}\subset\mathbb{O}$: $\det=ab-|z|^{2}=(X^{0})^{2}-(X^{3})^{2}-(X^{1})^{2}-(X^{2})^{2}$
is the $(1,3)$ form. The covariance of the $SL(2,\mathbb{C})$ action on
$H_{2}(\mathbb{C})$ is the textbook 4D statement $X\mapsto gXg^{\dagger}$,
$\det(gXg^{\dagger})=|\det g|^{2}\det X=\det X$ for $g\in SL(2,\mathbb{C})$.
\end{proof}

\subsection{The cubic norm \texorpdfstring{$N$}{N} is genuinely indefinite}
\label{subsec:app-lorentz:cubic}

The Albert algebra $J_{3}(\mathbb{O})$ of $3\times3$ octonion-Hermitian matrices,
$\dim_{\mathbb{R}}=3+3\cdot8=27$, carries the quadratic trace form
$Q(X)=\mathrm{Tr}(X^{2})$ and the cubic norm $N(X)=\det X$
(see \autoref{thm:Invs} and \autoref{thm:OrbitClass}). On a diagonal Jordan frame
$X=\mathrm{diag}(\xi_{1},\xi_{2},\xi_{3})$ one has $Q=\sum_{i}\xi_{i}^{2}$ and
$N=\xi_{1}\xi_{2}\xi_{3}$.

\begin{theorem}[Separation of invariance groups]
\label{thm:app-lorentz:groups}
The automorphism group $F_{4}=\mathrm{Aut}(J_{3}(\mathbb{O}))$ (compact, $\dim 52$)
preserves the Jordan product, hence preserves \emph{both} $Q=\mathrm{Tr}(X^{2})$ and
$N$; since $Q$ is positive-definite, $F_{4}\subset SO(27)$. The reduced structure
group $E_{6(-26)}$ ($\dim 78$) preserves the cubic norm $N$ \emph{only up to a real
character} $\lambda(g)$, $N(gX)=\lambda(g)N(X)$, and does \emph{not} preserve $Q$.
\emph{[Theorem.]}
\end{theorem}

\begin{proof}
Automorphisms fix $\mathbb{1}$, commute with $\circ$, and hence fix every
coefficient of the characteristic polynomial, in particular $\mathrm{Tr}(X^{2})$ and
$\det=N$; Euclidean signature of $Q$ then forces $F_{4}\subset SO(27)$. The reduced
structure group is by definition the linear group fixing $N$ up to scalar. If it also
preserved $Q$ it would lie in $SO(27)$ and be compact, contradicting
$\dim E_{6(-26)}=78>52=\dim F_{4}$ together with the non-compactness of the real form
$E_{6(-26)}$ (its maximal compact is exactly $F_{4}$, so the $26$ non-compact
directions move $Q$). Hence $Q$ is not $E_{6}$-invariant; only $N$ is.
\end{proof}

\begin{remark}[Orbits are classified by rank and \texorpdfstring{$N$}{N}, not by \texorpdfstring{$(Q,N)$}{(Q,N)}]
\label{rmk:app-lorentz:orbits}
Real $E_{6}$-orbits on the $\mathbf{27}$ are classified by
$\mathrm{rank}\in\{0,1,2,3\}$ and the value of $N$, via the $E_{6}$-covariant
sharp map $X^{\#}$ ($\mathrm{rank}\,X\le1\Leftrightarrow X^{\#}=0$;
$\mathrm{rank}\,X\le2\Leftrightarrow N(X)=0$). It is \emph{not} the pair $(Q,N)$
that separates orbits, since $Q$ is $F_{4}$-invariant only. This corrects any
statement to the contrary (cf. \autoref{thm:OrbitClass}).
\end{remark}

\begin{theorem}[Indefiniteness, and the Minkowski block inside \texorpdfstring{$N$}{N}]
\label{thm:app-lorentz:indef}
The cubic norm $N$ takes strictly positive, strictly negative, and zero values on
$J_{3}(\mathbb{O})\cong\mathbb{R}^{27}$ (it is an odd-degree form,
$N(-X)=-N(X)$), so it is genuinely indefinite. Moreover, restricting $N$ to the
block obtained by fixing the third pivot idempotent, $\xi_{3}=1$, $z_{1}=z_{2}=0$,
and freezing $Y=\begin{psmallmatrix}\xi_{1}&z_{3}\\ \bar z_{3}&\xi_{2}\end{psmallmatrix}\in H_{2}(\mathbb{O})$,
\begin{equation}
\label{eq:app-lorentz:block}
N(X)\big|_{\mathrm{block}}=\xi_{1}\xi_{2}-n(z_{3})=\det Y,
\end{equation}
so the indefinite Minkowski form $-\det=-\eta$ of \autoref{thm:app-lorentz:Mink} is
the restriction of the cubic norm $N$ to an $H_{2}(\mathbb{O})\cong\mathbb{R}^{1,9}$
sub-block of $J_{3}(\mathbb{O})$. \emph{[Theorem.]}
\end{theorem}

\begin{proof}
$N(\mathrm{diag}(1,1,1))=1>0$, $N(\mathrm{diag}(-1,1,1))=-1<0$,
$N(\mathrm{diag}(0,1,1))=0$, and oddness is immediate from the cubic homogeneity;
\eqref{eq:app-lorentz:block} is the direct evaluation of $N$ on the stated block.
\end{proof}

\begin{theorem}[Consistency: definite \texorpdfstring{$K|_{\mathfrak{m}}$}{K|m} and indefinite \texorpdfstring{$N$}{N} live on different spaces]
\label{thm:app-lorentz:consistency}
The positive-definite coset/target metric $K|_{\mathfrak{m}}$ and the indefinite
spacetime form are mutually consistent because they are forms on \emph{different}
representation spaces. The coset metric is positive-definite on
$\mathfrak{m}=\mathbf{16}_{-3}\oplus\overline{\mathbf{16}}_{+3}$ (the adjoint
complement, the $32$-dimensional internal tangent of EIII), and is unique by Schur
($\mathfrak{m}$ is real-irreducible of complex type, commutant $\mathbb{C}$; see
\autoref{thm:uniqueMetric}). The cubic norm $N$ is indefinite and lives on the
$\mathbf{27}$, a \emph{different} module. Therefore ``the coset metric is Euclidean''
and ``the cubic norm is indefinite'' are statements about distinct bilinear data on
distinct modules; there is no tension. The Lorentzian signature used for spacetime is
imported from $N|_{H_{2}(\mathbb{O})}$ via the soldering map P3, and is \emph{never}
asked of $K|_{\mathfrak{m}}$. \emph{[Theorem.]}
\end{theorem}

\subsection{The soldering postulate (P3)}
\label{subsec:app-lorentz:p3}

We now state the soldering precisely. Throughout, P3 is a \emph{postulate}: added
structure, not a theorem of the coset geometry. The theorem-level inputs are: a
vacuum manifold $M=\mathrm{EIII}=E_{6}/(\mathrm{Spin}(10)\times U(1))$, $\dim_{\mathbb{R}}M=32$,
with reductive split $\mathfrak{e}_{6}=\mathfrak{h}\oplus\mathfrak{m}$,
$\mathfrak{h}=\mathfrak{so}(10)\oplus\mathfrak{u}(1)$,
$\mathfrak{m}=\mathbf{16}_{-3}\oplus\overline{\mathbf{16}}_{+3}$, and Maurer--Cartan
form $\theta=\theta_{\mathfrak{h}}+\theta_{\mathfrak{m}}$
(\autoref{cor:VacuumManifold}, \autoref{eq:MC}); together with the Minkowski-block
facts of \autoref{thm:app-lorentz:Mink}--\autoref{thm:app-lorentz:SL2C} and
\autoref{thm:app-lorentz:indef}. The $32$ clock fields are the coordinates along
$\mathfrak{m}$: $28$ are physical massive scalars (single common radiative mass,
Schur on the irreducible $\mathfrak{m}$) and $4$ are gauge/soldering modes eaten by
world-volume diffeomorphism plus local Lorentz (see \autoref{thm:MassSpectrum}).

\begin{postulate}[Soldering / frame identification, P3a]
\label{post:app-lorentz:P3a}
There exists a soldering map $\sigma$ identifying four distinguished clock-field
frame directions, written $e^{a}{}_{\mu}$ ($a=0,1,2,3$; $\mu$ a world index), with a
tangent four-plane $\Pi_{x}\subset T_{x}\mathcal{M}^{4}$ of an emergent four-manifold
$\mathcal{M}^{4}$. The vierbein is the pull-back of the Maurer--Cartan component
along these directions,
\begin{equation}
\label{eq:app-lorentz:vierbein}
e^{a}{}_{\mu}(x)=\big\langle\,\tau^{a},\,\theta_{\mathfrak{m}}(\partial_{\mu})\,\big\rangle,
\end{equation}
where $\{\tau^{a}\}_{a=0}^{3}$ is the image of the $H_{2}(\mathbb{C})\cong\mathbb{R}^{1,3}$
basis under the embedding of \autoref{post:app-lorentz:P3d}. The vierbein
\emph{emerges dynamically}, as the pull-back of $\theta_{\mathfrak{m}}$. \emph{[Postulate P3.]}
\end{postulate}

\begin{postulate}[Induced Minkowski form, P3b]
\label{post:app-lorentz:P3b}
The induced quadratic form on $\Pi_{x}$ is the Minkowski form inherited from the
cubic norm via the $H_{2}(\mathbb{C})$ block,
\begin{equation}
\label{eq:app-lorentz:metric}
g_{\mu\nu}(x)=e^{a}{}_{\mu}(x)\,e^{b}{}_{\nu}(x)\,\eta_{ab},
\qquad \eta_{ab}=\mathrm{diag}(-1,+1,+1,+1),
\end{equation}
with $\eta$ the restriction of $-\det=-N$ to the $H_{2}(\mathbb{C})$ sub-block
(\autoref{thm:app-lorentz:SL2C}, \autoref{thm:app-lorentz:indef}). \emph{[Postulate P3.]}
\end{postulate}

\begin{postulate}[Dimension and signature are added, P3c]
\label{post:app-lorentz:P3c}
The numbers $d=4$ and signature $(1,3)=(-,+,+,+)$ are \emph{part of P3}: added
structure. They are \emph{not} consequences of the Killing form, the coset metric, or
any rank lemma. \emph{[Postulate P3.]}
\end{postulate}

\begin{postulate}[Frame reduction by a vierbein condensate, P3d]
\label{post:app-lorentz:P3d}
A vierbein condensate (a VEV of the soldering bilinear) selects the
$H_{2}(\mathbb{C})\cong\mathbb{R}^{1,3}$ sub-block of
$H_{2}(\mathbb{O})\cong\mathbb{R}^{1,9}$ and breaks the ambient frame symmetry
\begin{equation}
\label{eq:app-lorentz:reduction}
\mathrm{Spin}(1,9)=SL(2,\mathbb{O})\ \longrightarrow\
SL(2,\mathbb{C})\times \mathrm{Spin}(6)=\mathrm{Spin}(1,3)\times\mathrm{Spin}(6),
\end{equation}
the orthogonal complement of $\mathbb{R}^{1,3}$ in $\mathbb{R}^{1,9}$ being the
six-dimensional spacelike $\mathbb{R}^{6}$ acted on by $\mathrm{Spin}(6)\cong SU(4)$
(internal). The residual $\mathrm{Spin}(1,3)=SL(2,\mathbb{C})$ is the local Lorentz
group of the soldered four-plane. \emph{[Postulate P3.]}
\end{postulate}

\begin{theorem}[Internal consistency of P3]
\label{thm:app-lorentz:P3consistency}
Given \autoref{post:app-lorentz:P3a}--\autoref{post:app-lorentz:P3d}, the tensor
$g_{\mu\nu}$ of \eqref{eq:app-lorentz:metric} is a non-degenerate symmetric rank-$2$
tensor of Lorentzian signature $(1,3)$ on $\mathcal{M}^{4}$, invariant under the
residual local Lorentz group $SO(1,3)$; the four frame modes $e^{a}{}_{\mu}$ carry
exactly the gauge degrees of freedom eaten by world-volume diffeomorphisms together
with local Lorentz, consistent with the count $32=28\ \text{physical}+4\ \text{gauge/soldering}$.
\emph{[Theorem (conditional on P3).]}
\end{theorem}

\begin{proof}
Non-degeneracy and signature are inherited from
$\eta_{ab}=\mathrm{diag}(-1,+1,+1,+1)$ (\autoref{thm:app-lorentz:SL2C}) whenever
$\det(e^{a}{}_{\mu})\ne0$, which holds for a generic condensate;
$SO(1,3)$-invariance is the defining relation
$\eta_{ab}=\Lambda^{c}{}_{a}\Lambda^{d}{}_{b}\eta_{cd}$ for
$\Lambda\in SO(1,3)$. The gauge count is the Ward identity
$\nabla_{\mu}\big(\delta\Gamma/\delta e^{a}{}_{\mu}\big)=0$, so the four frame
directions are protected by diffeomorphism plus local Lorentz, not by a Hessian zero.
\end{proof}

\begin{remark}[What P3 does not claim]
\label{rmk:app-lorentz:notclaim}
P3 does not claim that $K|_{\mathfrak{m}}$ is Lorentzian: it stays positive-definite,
the kinetic/target-space metric of the $\sigma$-model. P3 does not use a $(4,28)$
split of $\mathfrak{m}$ with Lorentzian signature: the $4$ are gauge soldering modes,
the $28$ are Euclidean target coordinates. P3 does not build the four-plane from
$\nabla I_{2},\nabla I_{3}$ singlets: $\mathfrak{m}$ contains no $SO(10)$ singlet, and
$\nabla Q|_{X_{0}}=2X_{0}\notin\mathfrak{m}$. The four-plane is the soldered image of
$H_{2}(\mathbb{C})\subset\mathbf{27}$, not a sub-module of $\mathfrak{m}$.
\end{remark}

\begin{remark}[Anti-theorem: there is no Killing-form derivation of \texorpdfstring{$(1,3)$}{(1,3)}]
\label{rmk:app-lorentz:anti}
There is \emph{no} derivation of $d=4$ or of the signature $(1,3)$ from the Killing
form of $E_{6}$, from the coset metric $K|_{\mathfrak{m}}$, or from any rank lemma.
The Killing form of compact $E_{6}$ is negative-definite; its restriction to
$\mathfrak{m}$ is (sign-flipped to) positive-definite of signature $(32,0)$ --- zero
Lorentzian content. Any claim of a $(4,28)$ or $(1,3)$ Killing/coset signature is
false. Lorentzian signature enters only through P3, rooted in the indefinite cubic
norm $N$ / $H_{2}(\mathbb{O})\cong\mathbb{R}^{1,9}$.
\end{remark}

\subsection{Emergent time: Page--Wootters clock and Osterwalder--Schrader reconstruction}
\label{subsec:app-lorentz:time}

Lorentzian time is recovered by a relational-clock plus reconstruction-theorem
mechanism; no ``multiply $e^{0}$ by $i$'' Wick trick is used. The soldered metric
$\eta^{(1,3)}$ is already Lorentzian \emph{before} any continuation; the role of
analytic continuation is confined to the correlation functions in the clock-time
argument, justified by reflection positivity.

\begin{postulate}[Timelike clock selection, P3e]
\label{post:app-lorentz:clock}
Among the $32$ clock fields, the soldering condensate
(\autoref{post:app-lorentz:P3d}) selects a distinguished timelike direction, whose
clock field we call $\phi^{0}$ --- concretely the trace/$X^{0}$ direction of the
soldered $H_{2}(\mathbb{C})\subset J_{3}(\mathbb{O})$ block (the unique
$\eta$-timelike soldering leg $e^{0}{}_{\mu}$). The selection of \emph{one} timelike
direction is precisely the $(1,3)$ (rather than $(2,8)$, etc.) content of P3.
\emph{[Postulate P3.]}
\end{postulate}

\begin{theorem}[Page--Wootters relational evolution]
\label{thm:app-lorentz:PW}
Promote $\phi^{0}$ to a Page--Wootters clock: enlarge the kinematical Hilbert space
to $\mathcal{H}=\mathcal{H}_{\mathrm{clock}}\otimes\mathcal{H}_{\mathrm{sys}}$ and
impose the induced-gravity Hamiltonian constraint
\begin{equation}
\label{eq:app-lorentz:constraint}
\hat{H}_{\mathrm{phys}}\,|\Psi\rangle\!\rangle=0,
\qquad
\hat{H}_{\mathrm{phys}}=\hat{H}_{\mathrm{clock}}+\hat{H}_{\mathrm{sys}}
\ (+\ \text{interaction}),
\end{equation}
on physical states. Conditioning on clock readings $\phi^{0}=\tau$,
$|\psi(\tau)\rangle_{\mathrm{sys}}=\langle\phi^{0}=\tau|\Psi\rangle\!\rangle$, yields
the Schr\"odinger equation in relational time,
\begin{equation}
\label{eq:app-lorentz:schrodinger}
i\,\partial_{\tau}|\psi(\tau)\rangle_{\mathrm{sys}}
 =\hat{H}_{\mathrm{sys}}|\psi(\tau)\rangle_{\mathrm{sys}},
\end{equation}
with the unitary group $U(\tau)=e^{-i\hat{H}_{\mathrm{sys}}\tau}$ recovered on
physical states in the semiclassical regime where the soldered $g_{\mu\nu}$ admits a
timelike Killing vector and $\hat{H}_{\mathrm{sys}}=\int T_{00}$ along it.
\emph{[Theorem (Page--Wootters), conditional on the constraint and a good clock; global hyperbolicity assumed.]}
\end{theorem}

\begin{proof}
Under \eqref{eq:app-lorentz:constraint} with a good clock
($\hat{H}_{\mathrm{clock}}=\hat{p}_{\phi^{0}}$ canonically conjugate to
$\hat{\phi}^{0}$), the conditioned state obeys \eqref{eq:app-lorentz:schrodinger};
this is the Page--Wootters theorem. The identification $t\leftrightarrow X^{0}$ is
$\hat{H}_{\mathrm{sys}}=\int T_{00}$ along the timelike Killing direction.
\end{proof}

\begin{construction}[Osterwalder--Schrader reconstruction along \texorpdfstring{$\phi^{0}$}{phi0}]
\label{con:app-lorentz:OS}
The Euclidean $\sigma$-model built on the positive-definite $K|_{\mathfrak{m}}$
defines a Euclidean functional measure with Schwinger functions
$S_{n}(x_{1},\dots,x_{n})$. Take $\tau=\phi^{0}$ as Osterwalder--Schrader time and
assume the axioms: (OS0) Euclidean invariance of the spatial slice; (OS1) reflection
positivity with respect to $\theta:\tau\mapsto-\tau$ about a $\phi^{0}$-slice,
\begin{equation}
\label{eq:app-lorentz:rp}
\sum_{i,j}\bar{c}_{i}\,c_{j}\,S\!\left(\theta F_{i}\cdot F_{j}\right)\ \ge\ 0,
\qquad \mathrm{supp}\,F_{i}\subset\{\tau>0\};
\end{equation}
(OS2) symmetry and clustering.
\end{construction}

\begin{theorem}[OS reconstruction \texorpdfstring{$\Rightarrow$}{=>} Lorentzian Wightman theory]
\label{thm:app-lorentz:OS}
Given (OS0)--(OS2), the Osterwalder--Schrader reconstruction theorem produces a
Hilbert space $\mathcal{H}$, a positive self-adjoint Hamiltonian $\hat{H}\ge0$
generating $\tau$-translations as the contraction semigroup $e^{-\hat{H}\tau}$, and,
by analytic continuation of the Schwinger functions in the time argument
$\tau\to it$, Lorentzian Wightman functions on $(\mathcal{M}^{4},\eta^{(1,3)})$ with
unitary evolution $e^{-i\hat{H}t}$ covariant under the reconstructed $SO(1,3)$. The
OS theorem itself is rigorous given its hypotheses \emph{[Theorem]}; that the UCFT
$\sigma$-model \emph{satisfies} reflection positivity \eqref{eq:app-lorentz:rp} along
$\phi^{0}$ is \emph{[Indicative]} --- the precise property to be checked, and the
rigorous replacement for the Wick trick.
\end{theorem}

\begin{proof}[Proof (theorem part)]
Reflection positivity \eqref{eq:app-lorentz:rp} induces a physical inner product on
the quotient $\{F:\tau>0\}/(\text{null})$; the generator of the resulting contraction
semigroup $e^{-\hat{H}\tau}$ is self-adjoint and positive. The Schwinger functions,
analytic in a tube, continue to Wightman distributions satisfying the Wightman axioms,
including the spectrum condition and Lorentz covariance under the reconstructed
$SO(1,3)$ (Osterwalder--Schrader, 1973, 1975).
\end{proof}

\begin{remark}[Why this is not the Wick trick]
\label{rmk:app-lorentz:notwick}
The continuation $\tau\to it$ acts on the correlation functions in the
\emph{clock-time argument}, justified by reflection positivity
\eqref{eq:app-lorentz:rp} and the resulting analyticity tube --- it is \emph{not} the
ad hoc replacement $e^{0}{}_{\mu}\mapsto i\,e^{0}{}_{\mu}$ of a frame field. The
metric $\eta^{(1,3)}$ is the soldered (P3) form, Lorentzian before any continuation;
OS reconstruction supplies the unitary dynamics $e^{-i\hat{H}t}$ and the
positive-energy spectrum, with $\phi^{0}$ (the Page--Wootters clock) serving as the
OS reflection time. The two mechanisms dovetail:
\autoref{thm:app-lorentz:PW} gives relational unitarity $U(\tau)$ on physical states,
and \autoref{thm:app-lorentz:OS} gives the positive-energy, Lorentz-covariant
Wightman theory.
\end{remark}

\subsection{Summary: theorem versus postulate}
\label{subsec:app-lorentz:summary}

The division of labour is as follows.
\begin{itemize}
\item \textbf{Theorems} (proved): $(H_{2}(\mathbb{O}),-\det)\cong\mathbb{R}^{1,9}$
  (\autoref{thm:app-lorentz:Mink}); $SL(2,\mathbb{O})\cong\mathrm{Spin}(1,9)$ and
  $SL(2,\mathbb{C})\cong\mathrm{Spin}(1,3)$
  (\autoref{thm:app-lorentz:SL2O}, \autoref{thm:app-lorentz:SL2C}); $F_{4}$ preserves
  $Q$, $E_{6}$ preserves only $N$ (\autoref{thm:app-lorentz:groups}); $N$ indefinite
  and the Minkowski block $\det Y\subset N$ (\autoref{thm:app-lorentz:indef});
  definite $K|_{\mathfrak{m}}$ and indefinite $N$ are consistent
  (\autoref{thm:app-lorentz:consistency}); internal consistency of $g_{\mu\nu}$ given
  P3 (\autoref{thm:app-lorentz:P3consistency}); the Page--Wootters and OS
  reconstruction theorems (\autoref{thm:app-lorentz:PW}, \autoref{thm:app-lorentz:OS}).
\item \textbf{Postulate} (P3, added structure, cf. P1--P6): the soldering map, the
  choice of a four-plane, $d=4$, signature $(1,3)$, the vierbein condensate and frame
  reduction, and the timelike clock selection
  (\autoref{post:app-lorentz:P3a}--\autoref{post:app-lorentz:P3d},
  \autoref{post:app-lorentz:clock}).
\item \textbf{Indicative}: that the UCFT $\sigma$-model satisfies reflection
  positivity along $\phi^{0}$ (\autoref{thm:app-lorentz:OS}).
\end{itemize}
In short, the Lorentzian arena is a \emph{conditional embedding}: the algebra
$H_{2}(\mathbb{O})\cong\mathbb{R}^{1,9}$ and $SL(2,\mathbb{O})\cong\mathrm{Spin}(1,9)$
are theorems, while the selection of the physical $(1,3)$ four-plane and the recovery
of unitary Lorentzian dynamics rest on the soldering postulate P3 plus an indicative
reflection-positivity assumption.
% --- end inlined: unit-app-lorentz.tex ---
% --- begin inlined: unit-app-spectrum.tex ---
%======================================================================%
\section{Branchings and the Coset Module}
\label{sec:app-spectrum:main}
%======================================================================%

This appendix collects the representation-theoretic facts on which the
dynamical claims of the manuscript rest: the branching of the
$\mathbf{78}$ (adjoint) and $\mathbf{27}$ (fundamental) of $E_6$ under
the maximal compact subgroup $H=\Spin(10)\times U(1)$, the
$U(1)$ charge normalization, and the structural fact that the coset
tangent module
\[
  \mathfrak m \;=\; \mathbf{16}_{-3}\oplus\overline{\mathbf{16}}_{+3}
\]
is \emph{real-irreducible of complex type}, with real commutant
$\End_H(\mathfrak m)\cong\mathbb C$.  These results fix (i) that the
invariant coset metric $K|_{\mathfrak m}$ is unique up to scale and
positive-definite (\autoref{thm:app-spectrum:metric}); and (ii) that one
radiative Coleman--Weinberg loop produces a \emph{single common mass} on
all $32$ clock fields by Schur's lemma, the input used in the mass-spectrum
and induced-gravity units (\autoref{thm:app-spectrum:schurmass}).

Throughout, $H=\Spin(10)\times U(1)$ is the (compact) stabilizer of a
rank-one primitive idempotent, so the relevant real form for the dynamics
is the compact $E_6$ (equivalently $E_{6(-14)}$, whose maximal compact is
exactly $H$); see \autoref{sec:realform:main}. The Killing form on this
real form is negative-definite on $\mathfrak e_6$, hence
\emph{positive-definite} on the coset $\mathfrak m$ once the standard
sign for symmetric spaces of compact type is taken. There is no
$(4,28)$ Killing signature and no Lorentzian content in $\mathfrak m$;
Lorentzian signature is supplied separately by the soldering postulate
P3 (\autoref{sec:soldering:main}).

%%%%%%%%%%%%%%%%%%%%%%%%%%%%%%%%%%%%%%%%%%%%%%%%%%%%%%%%%%%%%%%%%%%%%%
\subsection{The $U(1)$ charge normalization}
\label{sec:app-spectrum:u1}
%%%%%%%%%%%%%%%%%%%%%%%%%%%%%%%%%%%%%%%%%%%%%%%%%%%%%%%%%%%%%%%%%%%%%%

Fix the abelian generator $Y\in\mathfrak u(1)\subset\mathfrak h$ that
commutes with $\mathfrak{so}(10)$. We normalize $Y$ so that the
$\Spin(10)$ spinor $\mathbf{16}$ appearing in the \emph{fundamental}
$\mathbf{27}$ carries charge $+1$:
\begin{equation}
  q\bigl(\mathbf{16}\subset\mathbf{27}\bigr) \;=\; +1 .
  \label{eq:app-spectrum:norm}
\end{equation}
This is the smallest charge quantum appearing in the $\mathbf{27}$ and
is the convenient unit because it makes every charge in the
$\mathbf{27}$ and the $\mathbf{78}$ an integer. With this choice the
remaining charges are fixed by requiring (a) that the cubic invariant
$N$ (the $E_6$ singlet in $\Sym^3(\mathbf{27})$) be neutral, and (b)
that the adjoint decomposition close consistently. We record both
branchings, then verify these two conditions.

\begin{theorem}[Branchings under $H=\Spin(10)\times U(1)$]
\label{thm:app-spectrum:branchings}
With the normalization \eqref{eq:app-spectrum:norm}, the fundamental and
adjoint of $E_6$ decompose as
\begin{align}
  \mathbf{27}
    &= \mathbf{1}_{+4}\;\oplus\;\mathbf{10}_{-2}\;\oplus\;\mathbf{16}_{+1},
    \label{eq:app-spectrum:27}\\[2pt]
  \mathbf{78}
    &= \mathbf{45}_{0}\;\oplus\;\mathbf{1}_{0}\;\oplus\;
       \mathbf{16}_{-3}\;\oplus\;\overline{\mathbf{16}}_{+3}.
    \label{eq:app-spectrum:78}
\end{align}
The coset (broken) directions are
\begin{equation}
  \mathfrak h = \mathfrak{so}(10)\oplus\mathfrak u(1)
     = \mathbf{45}_{0}\oplus\mathbf{1}_{0},
  \qquad
  \mathfrak m = \mathbf{16}_{-3}\oplus\overline{\mathbf{16}}_{+3},
  \qquad \dim_{\mathbb R}\mathfrak m = 32 .
  \label{eq:app-spectrum:hm}
\end{equation}
\end{theorem}

\begin{proof}
\emph{Dimensions and $\Spin(10)$ content.}
$1+10+16=27$ and $45+1+16+16=78$, matching $\dim\mathbf{27}=27$ and
$\dim\mathbf{78}=78$. The $\Spin(10)$ pieces are forced by the chain
$E_6\supset\Spin(10)\times U(1)$: the adjoint $\mathbf{78}$ contains
$\mathfrak{so}(10)=\mathbf{45}$, the abelian generator $\mathbf{1}$, and
the off-diagonal coset which is the unique pair of chiral half-spinors
$\mathbf{16}\oplus\overline{\mathbf{16}}$ (a Hermitian symmetric pair,
since $E_6/H=\mathrm{EIII}$ is Hermitian symmetric of complex type). The
$\mathbf{27}$ is the smallest faithful rep; under $\Spin(10)$ it is the
sum of a singlet, a vector $\mathbf{10}$, and a spinor $\mathbf{16}$.

\emph{Charge assignment.} Set $q(\mathbf{16}\subset\mathbf{27})=+1$ by
\eqref{eq:app-spectrum:norm}. The cubic $E_6$ singlet
$N\in\Sym^3(\mathbf{27})$ contains a $\Spin(10)$-singlet coupling
$\mathbf{1}\cdot\mathbf{10}\cdot\mathbf{10}$ and a coupling
$\mathbf{10}\cdot\mathbf{16}\cdot\mathbf{16}$ (the unique
$\Spin(10)$-invariants in $\mathbf{10}\otimes\mathbf{16}\otimes\mathbf{16}$
and in $\mathbf{1}\otimes\mathbf{10}\otimes\mathbf{10}$). Charge
neutrality of $N$ requires
\[
  q(\mathbf 1)+2\,q(\mathbf{10})=0,
  \qquad
  q(\mathbf{10})+2\,q(\mathbf{16})=0 .
\]
With $q(\mathbf{16})=+1$ the second gives $q(\mathbf{10})=-2$, and the
first gives $q(\mathbf 1)=-2\,q(\mathbf{10})=+4$, yielding
\eqref{eq:app-spectrum:27}.

For the adjoint, the coset $\mathfrak m$ transforms in the
$\mathbf{16}\oplus\overline{\mathbf{16}}$ that maps
$\mathbf{27}\to\mathbf{27}$ raising/lowering the $\Spin(10)$ grading by
one spinor step. Acting on the graded $\mathbf{27}$
\eqref{eq:app-spectrum:27}, the raising operator carries the
$\mathbf{16}_{+1}$ piece into the $\mathbf{1}_{+4}\oplus\mathbf{10}_{-2}$
pieces; matching $U(1)$ charges across the three components forces the
coset spinor to carry charge $-3$ and its conjugate $+3$:
\[
  q(\mathbf{16}\subset\mathfrak m)
   = q(\mathbf{1}_{+4})-q(\mathbf{16}_{+1}) = 4-1 = 3
   \;\;(\text{or}\;\; q(\mathbf{10}_{-2})-q(\mathbf{16}_{+1})=-3),
\]
which are conjugate, so $\mathfrak m=\mathbf{16}_{-3}\oplus
\overline{\mathbf{16}}_{+3}$ with the diagonal $H$ at charge $0$. This is
\eqref{eq:app-spectrum:78} and \eqref{eq:app-spectrum:hm}.
\end{proof}

\begin{remark}\label{rem:app-spectrum:nosinglet}
The coset $\mathfrak m=\mathbf{16}_{-3}\oplus\overline{\mathbf{16}}_{+3}$
contains \emph{no} $\Spin(10)$ singlet, and indeed no $\Spin(10)$ vector.
Any decomposition of the form $2\times\mathbf{10}+\mathbf 4$,
$\mathbf{10}+\overline{\mathbf{10}}+\mathbf 4$, or one containing a
singlet-in-$\mathfrak m$ is therefore \emph{false}: those would neither
reproduce the chiral spinor content of $\mathrm{EIII}$ nor sum correctly
($32=16+16$, not $2\cdot10+4=24$). The absence of a singlet in
$\mathfrak m$ is what makes $\mathfrak m$ irreducible and forces the
single-common-mass conclusion below.
\end{remark}

%%%%%%%%%%%%%%%%%%%%%%%%%%%%%%%%%%%%%%%%%%%%%%%%%%%%%%%%%%%%%%%%%%%%%%
\subsection{Relevant tensor products of $\Spin(10)$ spinors}
\label{sec:app-spectrum:tensor}
%%%%%%%%%%%%%%%%%%%%%%%%%%%%%%%%%%%%%%%%%%%%%%%%%%%%%%%%%%%%%%%%%%%%%%

We record the $\Spin(10)$ tensor products used repeatedly in the gauge,
metric, and Yukawa arguments. Writing $\mathbf{16}$ for a chiral
half-spinor:
\begin{align}
  \mathbf{16}\otimes\mathbf{16}
    &= \mathbf{10}\;\oplus\;\mathbf{120}\;\oplus\;\mathbf{126},
    \label{eq:app-spectrum:1616}\\
  \mathbf{16}\otimes\overline{\mathbf{16}}
    &= \mathbf{1}\;\oplus\;\mathbf{45}\;\oplus\;\mathbf{210}.
    \label{eq:app-spectrum:16bar}
\end{align}
Dimension checks: $16\cdot16=256=10+120+126$ and
$16\cdot16=256=1+45+210$. In \eqref{eq:app-spectrum:1616} the
$\mathbf{10}$ is symmetric, the $\mathbf{126}$ is symmetric (self-dual
$5$-form), and the $\mathbf{120}$ ($3$-form) is antisymmetric:
\begin{equation}
  \Sym^2(\mathbf{16}) = \mathbf{10}\oplus\mathbf{126},
  \qquad
  \wedge^2(\mathbf{16}) = \mathbf{120}.
  \label{eq:app-spectrum:symanti}
\end{equation}
In \eqref{eq:app-spectrum:16bar} the singlet $\mathbf 1$ is the
$H$-invariant Hermitian pairing of a spinor with its conjugate.

\begin{remark}[Yukawa consequence]\label{rem:app-spectrum:yukawa}
The decompositions \eqref{eq:app-spectrum:1616}--\eqref{eq:app-spectrum:16bar}
imply there is \emph{no} $H$-singlet in
$\mathbf{16}\otimes\mathbf{16}\otimes\{\mathbf{16},\overline{\mathbf{16}}\}$:
a scalar valued in $\mathbf{16}$ or $\overline{\mathbf{16}}$ cannot form
an invariant Yukawa coupling $-y\,\bar\psi\,\phi\,\psi$. Since $H$ is
unbroken at the vacuum, fermion masses must come instead from the real
$\Spin(10)$ Higgs/Yukawa sector ($\mathbf{10}_H$ for Dirac masses via the
$\mathbf{10}\subset\mathbf{16}\otimes\mathbf{16}$, and
$\overline{\mathbf{126}}_H$ for Majorana masses via the
$\mathbf{126}\subset\mathbf{16}\otimes\mathbf{16}$), as assumed in
Postulate~P6 (\autoref{sec:sm:yukawa}). This is an embedding, not a
derivation.
\end{remark}

%%%%%%%%%%%%%%%%%%%%%%%%%%%%%%%%%%%%%%%%%%%%%%%%%%%%%%%%%%%%%%%%%%%%%%
\subsection{Reality type and commutant of the coset module}
\label{sec:app-spectrum:reality}
%%%%%%%%%%%%%%%%%%%%%%%%%%%%%%%%%%%%%%%%%%%%%%%%%%%%%%%%%%%%%%%%%%%%%%

We now establish the structural fact that controls both the coset metric
and the radiative mass spectrum.

\begin{theorem}[$\mathfrak m$ is real-irreducible of complex type]
\label{thm:app-spectrum:complextype}
The real $32$-dimensional $H$-module
$\mathfrak m=\mathbf{16}_{-3}\oplus\overline{\mathbf{16}}_{+3}$ is
\emph{irreducible over $\mathbb R$} and is of \emph{complex type}: its
real endomorphism algebra is
\[
  \End_H(\mathfrak m)\;\cong\;\mathbb C .
\]
\end{theorem}

\begin{proof}
The complex $\Spin(10)$ spinor $\mathbf{16}$ is a representation of
\emph{complex type} (it is not self-conjugate;
$\overline{\mathbf{16}}\not\cong\mathbf{16}$ for $\Spin(10)$, since
$5\equiv1\pmod 4$ places $\Spin(10)$ in the class with inequivalent chiral
spinors). The $U(1)$ factor multiplies the two summands by the opposite
characters $e^{-3i\theta}$ and $e^{+3i\theta}$. As a \emph{complex}
$H$-representation, $\mathfrak m_{\mathbb C}=\mathbf{16}_{-3}\oplus
\overline{\mathbf{16}}_{+3}$ is a sum of two \emph{inequivalent}
irreducibles (different $\Spin(10)$ type \emph{and} opposite $U(1)$
charge), so $\End_{H,\mathbb C}(\mathfrak m_{\mathbb C})\cong
\mathbb C\oplus\mathbb C$.

Now $\mathfrak m$ is the \emph{real} module obtained by viewing
$\mathbf{16}_{-3}\oplus\overline{\mathbf{16}}_{+3}$ as the realification
of the single complex irreducible $\mathbf{16}_{-3}$: complex conjugation
on the underlying real space is an antilinear $H$-map that exchanges the
two summands (it sends $\mathbf{16}_{-3}\mapsto\overline{\mathbf{16}}_{+3}$,
respecting both the conjugate $\Spin(10)$ type and the sign-reversed
$U(1)$ charge). Hence $\mathfrak m\cong(\mathbf{16}_{-3})_{\mathbb R}$,
the realification of one complex irreducible. The standard reality
trichotomy then gives, for a complex irreducible $W$ that is not
self-conjugate,
\[
  \End_H\bigl(W_{\mathbb R}\bigr)\cong\mathbb C ,
\]
i.e.\ complex type, and $W_{\mathbb R}$ is irreducible over $\mathbb R$
(its only real-invariant subspaces are $0$ and the whole space, because
any proper one would have to be one of the two complex summands, which are
not $\mathbb R$-invariant under the conjugation map). This proves
irreducibility over $\mathbb R$ and $\End_H(\mathfrak m)\cong\mathbb C$.
\end{proof}

The commutant being $\mathbb C$ (rather than $\mathbb R$ or
$\mathbb H$) is the precise sense in which $\mathfrak m$ is ``of complex
type.'' It supplies exactly one $\mathbb R$-linear complex structure
$J\in\End_H(\mathfrak m)$ with $J^2=-\mathbf 1$ commuting with $H$ --- the
complex structure of the Hermitian symmetric space $\mathrm{EIII}$.

%%%%%%%%%%%%%%%%%%%%%%%%%%%%%%%%%%%%%%%%%%%%%%%%%%%%%%%%%%%%%%%%%%%%%%
\subsection{Uniqueness and definiteness of the invariant metric}
\label{sec:app-spectrum:metric}
%%%%%%%%%%%%%%%%%%%%%%%%%%%%%%%%%%%%%%%%%%%%%%%%%%%%%%%%%%%%%%%%%%%%%%

\begin{theorem}[Unique definite invariant metric]
\label{thm:app-spectrum:metric}
The space of $H$-invariant real symmetric bilinear forms on
$\mathfrak m$ is \emph{one-dimensional}. Its generator is the restriction
of the Killing form of (compact) $E_6$ to $\mathfrak m$, and is
\emph{positive-definite}. Consequently the coset/target-space metric
\begin{equation}
  \mathfrak g_{\alpha\beta} \;=\; f^2\,K_{\alpha\beta}\big|_{\mathfrak m},
  \qquad f^2>0,
  \label{eq:app-spectrum:metric}
\end{equation}
is the unique (up to the positive scale $f^2$) admissible
$\sigma$-model kinetic metric, and it is Euclidean.
\end{theorem}

\begin{proof}
\emph{Counting invariants via the $U(1)$ charge.} An $H$-invariant
symmetric bilinear form on $\mathfrak m$ is an $H$-invariant element of
$\Sym^2(\mathfrak m^\ast)$. Decompose
\[
  \Sym^2\bigl(\mathbf{16}_{-3}\oplus\overline{\mathbf{16}}_{+3}\bigr)
   = \underbrace{\Sym^2(\mathbf{16}_{-3})}_{\text{charge }-6}
   \;\oplus\;
   \underbrace{\Sym^2(\overline{\mathbf{16}}_{+3})}_{\text{charge }+6}
   \;\oplus\;
   \underbrace{\mathbf{16}_{-3}\otimes\overline{\mathbf{16}}_{+3}}_{\text{charge }0}.
\]
The first two blocks carry net $U(1)$ charge $\mp6\neq0$, so they contain
\emph{no} $U(1)$-singlet and hence no $H$-invariant: the $U(1)$ charge
\emph{forbids} the symmetric $\mathbf{16}\times\mathbf{16}$ and
$\overline{\mathbf{16}}\times\overline{\mathbf{16}}$ invariants. The only
charge-$0$ block is the cross-pairing
$\mathbf{16}_{-3}\otimes\overline{\mathbf{16}}_{+3}$, whose $\Spin(10)$
content is $\mathbf 1\oplus\mathbf{45}\oplus\mathbf{210}$ by
\eqref{eq:app-spectrum:16bar}. Exactly one $\Spin(10)$ singlet $\mathbf 1$
appears, and it is symmetric under the exchange that defines a bilinear
\emph{form}. Hence there is a \emph{unique} $H$-invariant symmetric form,
the charge-$0$ cross-pairing.

\emph{Identification with the Killing form and definiteness.} The Killing
form of $E_6$ restricted to $\mathfrak m$ is an $H$-invariant symmetric
form; by the uniqueness just proved it must coincide (up to scale) with
the cross-pairing. On the \emph{compact} real form of $E_6$ the Killing
form is negative-definite on all of $\mathfrak e_6$, hence on the symmetric
complement $\mathfrak m$ the induced metric of the symmetric space of
compact type is positive-definite; equivalently the Hermitian pairing
$\langle v,v\rangle=\sum_i|v_i|^2$ furnished by the cross-pairing is
strictly positive. Schur's lemma (in the strong form
$\End_H(\mathfrak m)\cong\mathbb C$, \autoref{thm:app-spectrum:complextype})
guarantees there is no second, independent invariant form to spoil
definiteness. Positivity of the $\sigma$-model kinetic term then fixes the
scale to $f^2>0$, giving \eqref{eq:app-spectrum:metric}.
\end{proof}

\begin{remark}\label{rem:app-spectrum:nolorentz}
The metric \eqref{eq:app-spectrum:metric} has signature $(32,0)$ ---
Euclidean, with \emph{zero} Lorentzian content. There is no $(4,28)$
Killing/coset signature, and the physical Lorentzian signature $(1,3)$ in
the convention $(-,+,+,+)$ does \emph{not} fall out of the coset geometry.
Signature enters only through the soldering postulate~P3, soldered from the
cubic norm $N$ via $H_2(\mathbb O)\cong\mathbb R^{1,9}$
(\autoref{sec:soldering:main}).
\end{remark}

%%%%%%%%%%%%%%%%%%%%%%%%%%%%%%%%%%%%%%%%%%%%%%%%%%%%%%%%%%%%%%%%%%%%%%
\subsection{One common radiative mass (Schur)}
\label{sec:app-spectrum:mass}
%%%%%%%%%%%%%%%%%%%%%%%%%%%%%%%%%%%%%%%%%%%%%%%%%%%%%%%%%%%%%%%%%%%%%%

\begin{theorem}[Single common scalar mass]
\label{thm:app-spectrum:schurmass}
At tree level the $32$ clock fields are exact Goldstone modes,
$\Hess V_{\mathrm{tree}}|_{\mathfrak m}=0$. The one-loop Coleman--Weinberg
Hessian on $\mathfrak m$ is an $H$-invariant symmetric operator; by
irreducibility of $\mathfrak m$ over $\mathbb R$
(\autoref{thm:app-spectrum:complextype}) and Schur's lemma it is a
\emph{single scalar multiple of the identity}. Hence all clock fields
acquire one common radiative mass. After removing the four gauge/soldering
frame modes (eaten by diffeomorphism + local Lorentz, P3), the physical
spectrum is
\begin{equation}
  32\ \text{real}\;=\;\underbrace{28\ \text{massive scalars (single common mass)}}_{\text{use }28\text{ in all loop sums}}
  \;+\;4\ \text{gauge/soldering (eaten)} .
  \label{eq:app-spectrum:count}
\end{equation}
With $[f]=\mathrm{mass}$, the dimension-$4$ Hessian eigenvalue is
$f^2\mu^2$, so the canonically normalized physical mass is
\begin{equation}
  m^2=\mu^2 .
  \label{eq:app-spectrum:mass}
\end{equation}
\end{theorem}

\begin{proof}
Tree-level masslessness is Goldstone's theorem for the broken $E_6/H$
directions. The one-loop effective potential is built from $H$-invariant
data (the vacuum preserves $H$), so its Hessian
$\mathcal H\in\End_H(\mathfrak m)$ is an $H$-equivariant symmetric
operator. Over the \emph{real} commutant
$\End_H(\mathfrak m)\cong\mathbb C$ a symmetric (self-adjoint) operator is
real, hence lies in $\mathbb R\cdot\mathbf 1\subset\mathbb C$: it is a
real scalar times the identity. (Equivalently, the only $H$-invariant
symmetric forms are multiples of the unique metric of
\autoref{thm:app-spectrum:metric}, so $\mathcal H=m^2\,\mathfrak g$.) Thus
all $32$ modes share one eigenvalue.

The four frame directions are not Hessian zeros of an internal singlet
subspace --- $\mathfrak m$ has no $\Spin(10)$ singlet
(\autoref{rem:app-spectrum:nosinglet}) --- but are gauge modes protected
by the Ward identity $\nabla_\mu(\delta\Gamma/\delta e^a_\mu)=0$ and eaten
by the diffeomorphism $+$ local-Lorentz gauge symmetry. Removing them
leaves $28$ massive physical scalars, all of common mass, giving
\eqref{eq:app-spectrum:count}. Canonical normalization of the
kinetic term \eqref{eq:app-spectrum:metric} divides the dimension-$4$
Hessian eigenvalue $f^2\mu^2$ by $f^2$, yielding
\eqref{eq:app-spectrum:mass}.
\end{proof}

\begin{remark}\label{rem:app-spectrum:downstream}
The count $28$ in \eqref{eq:app-spectrum:count} and the mass
\eqref{eq:app-spectrum:mass} are the inputs to the induced Planck mass
$M_{\mathrm{Pl}}^2=(28/6)\,f^2\mu^2\log(\Lambda^2/\mu^2)/(16\pi^2)>0$ and
to the gauge $\beta$-function scalar count $n_S=28$
(\autoref{sec:gravity:main}). The $27\times27$ scalar block likewise
inherits the multiplicities of \eqref{eq:app-spectrum:27},
$M=\diag(\mu_1\,\mathbf 1_1,\ \mu_{10}\,\mathbf 1_{10},\
\mu_{16}\,\mathbf 1_{16})$, the residual symmetry being $H$, not $E_6$.
\end{remark}

%%%%%%%%%%%%%%%%%%%%%%%%%%%%%%%%%%%%%%%%%%%%%%%%%%%%%%%%%%%%%%%%%%%%%%
\subsection*{Status of the claims in this appendix}
\label{sec:app-spectrum:status}
%%%%%%%%%%%%%%%%%%%%%%%%%%%%%%%%%%%%%%%%%%%%%%%%%%%%%%%%%%%%%%%%%%%%%%

All statements of this appendix are \textbf{theorems}: the branchings
\eqref{eq:app-spectrum:27}--\eqref{eq:app-spectrum:78}, the tensor
products \eqref{eq:app-spectrum:1616}--\eqref{eq:app-spectrum:16bar}, the
complex-type/commutant result (\autoref{thm:app-spectrum:complextype}),
the unique definite metric (\autoref{thm:app-spectrum:metric}), and the
single-common-mass result (\autoref{thm:app-spectrum:schurmass}) follow
from the representation theory of $E_6\supset\Spin(10)\times U(1)$ on the
compact real form. The \emph{use} of these facts to assert the physical
$(1,3)$ signature, three generations, and the Yukawa sector relies on the
added postulates P3, P5, P6 and is flagged where invoked.
% --- end inlined: unit-app-spectrum.tex ---
% --- begin inlined: unit-app-ym.tex ---
%======================================================================%
\section{The Induced Yang--Mills Coefficient}
\label{sec:app-ym:main}
%======================================================================%

This appendix derives the one-loop coefficient of the composite
$\SOTen\times\UOne$ Yang--Mills term obtained by integrating out the
massive clock fields, and traces how that coefficient feeds the
asymptotic-freedom slope $b_{0}$ of the functional renormalization group.
The computation is presented at \emph{derivation grade}: the heat-kernel
machinery, the group-theory weights, and the resulting logarithm are all
theorems given the stated background; the embedding of the clock-field
$\sigma$-model into a four-dimensional Lorentzian arena is the soldering
\emph{postulate} P3, invoked only to provide the $X=M^{4}$ over which the
proper-time integral is taken.

Throughout we use the metric signature $(-,+,+,+)$, the decay constant
$f$ with $[f]=\text{mass}$, the breaking/mass scale $\mu$, and the UV
cutoff $\Lambda$. The hatted coupling $\ghat^{2}$ is dimensionless. The
relevant gauge algebra is $\mathfrak{so}_{10}\oplus\mathfrak u_{1}$, and
\emph{all} group-theory weights are taken in $\SOTen$ with the locked
Casimirs
\begin{equation}\label{eq:app-ym:casimirs}
  C_{A}(\SOTen)=8,\qquad
  T(\mathbf{16})=2,\qquad
  C_{2}(\mathbf{16})=\tfrac{45}{8}.
\end{equation}
The $E_{6}$ value $C_{A}(E_{6})=12$ is \emph{never} mixed into these
sums: it is the spinor index $T(\mathbf{16})=2$, not a count of sixteen
index-$1$ fundamentals, that enters.

%----------------------------------------------------------------------%
\subsection{The Maurer--Cartan curvature and the one-loop determinant}
\label{sec:app-ym:setup}
%----------------------------------------------------------------------%

Let $g(x)\in E_{6}$ (compact real form, the dynamical form fixed in
\autoref{sec:foundation:realform}) be a coset representative of the
vacuum map into $M=E_{6}/(\SOTen\times\UOne)=\mathrm{EIII}$. The
Maurer--Cartan form splits along the reductive decomposition
$\mathfrak e_{6}=\mathfrak h\oplus\mathfrak m$,
\begin{equation}\label{eq:app-ym:MC}
  g^{-1}\partial_{\mu}g
  = \mathcal A_{\mu}(x) + e_{\mu}(x),
  \qquad
  \mathcal A_{\mu}\in\mathfrak{so}_{10}\oplus\mathfrak u_{1},
  \qquad
  e_{\mu}\in\mathfrak m
  =\mathbf{16}_{-3}\oplus\overline{\mathbf{16}}_{+3},
\end{equation}
with $\mathfrak h$-part $\mathcal A_{\mu}$ the composite connection of
\autoref{sec:gauge:composite} and $\mathfrak m$-part $e_{\mu}$ the
clock-field vielbein. The $\mathfrak h$-curvature
\begin{equation}\label{eq:app-ym:curv}
  \mathcal F_{\mu\nu}
  =\partial_{\mu}\mathcal A_{\nu}-\partial_{\nu}\mathcal A_{\mu}
   +[\mathcal A_{\mu},\mathcal A_{\nu}]
  =-\,P_{\mathfrak h}\bigl([e_{\mu},e_{\nu}]\bigr)
\end{equation}
is, by the Maurer--Cartan structure equation, the projection onto
$\mathfrak h$ of the bracket of two coset directions. This is the gauge
field whose kinetic term we now generate radiatively.

\begin{remark}[Corrected massive multiplet]\label{rem:app-ym:multiplet}
The fluctuations that couple minimally to $\mathcal A_{\mu}$ are the
\emph{coset} directions $\mathfrak m=\mathbf{16}_{-3}\oplus
\overline{\mathbf{16}}_{+3}$, a single real-irreducible $H$-module of
complex type (commutant $\mathbb C$; see
\autoref{sec:foundation:tangent}). It contains \emph{no} $\SOTen$
singlet. Earlier drafts decomposed the massive sector as
$\mathbf{10}_{+2}\oplus\mathbf{10}_{-2}\oplus\mathbf 4_{0}$, with total
Dynkin index $\tfrac52$; that decomposition is incorrect (it places a
nonexistent singlet $\mathbf 4_{0}$ in $\mathfrak m$ and uses
fundamental rather than spinor reps). The corrected weight is computed
below from $\mathbf{16}\oplus\overline{\mathbf{16}}$.
\end{remark}

Expanding the kinetic action to quadratic order about the vacuum, the
$32$ real coset fluctuations $\xi$ acquire the common radiative mass
$m^{2}=f^{2}\mu^{2}$ (one common eigenvalue by Schur on the irreducible
$\mathfrak m$; \autoref{sec:gravity:spectrum}). Of these, $28$ are
physical massive scalars and $4$ are the gauge/soldering frame modes
eaten by diffeomorphism $+$ local Lorentz. The Euclidean quadratic
action in the composite background reads
\begin{equation}\label{eq:app-ym:Squad}
  S_{m}[\xi]
  =\frac12\int_{X}\xi^{\dagger}\bigl[-D^{2}+m^{2}\bigr]\xi\,d^{4}x,
  \qquad
  D_{\mu}=\partial_{\mu}+\mathcal A_{\mu},
\end{equation}
and integrating out $\xi$ produces the determinant
$[\det(-D^{2}+m^{2})]^{-1/2}$, i.e. the one-loop effective action
\begin{equation}\label{eq:app-ym:W}
  \Gamma_{\mathrm{1L}}
  =\frac12\Tr\log(-D^{2}+m^{2})
  =-\frac12\int_{\Lambda^{-2}}^{\infty}\frac{ds}{s}\,
    e^{-sm^{2}}\,\Tr\,e^{-s(-D^{2})}.
\end{equation}

%----------------------------------------------------------------------%
\subsection{Heat-kernel extraction of the gauge weight}
\label{sec:app-ym:heatkernel}
%----------------------------------------------------------------------%

For small proper time the heat trace expands as
\begin{equation}\label{eq:app-ym:HK}
  \Tr\,e^{-s(-D^{2})}
  =\frac{1}{(4\pi s)^{2}}\sum_{n\ge0}s^{n}\,a_{2n},
\end{equation}
and the $\mathcal F^{2}$ content sits in the Seeley--DeWitt coefficient
$a_{4}$. For a real scalar multiplet in an $H$-representation $R$ with
Dynkin index $T(R)$ normalized by
$\Tr_{R}(T^{a}T^{b})=T(R)\,\delta^{ab}$,
\begin{equation}\label{eq:app-ym:a4}
  a_{4}(R)\supset
  \frac{1}{12}\int_{X}\Tr_{R}\bigl(\mathcal F_{\mu\nu}
  \mathcal F^{\mu\nu}\bigr)\,d^{4}x ,
\end{equation}
the omitted terms being curvature-squared invariants on $X$ suppressed
by powers of $\mu/\Lambda$.

\begin{lemma}[Total gauge weight of the coset multiplet]
\label{lem:app-ym:weight}
The total Dynkin index carried by the massive coset multiplet
$\mathfrak m=\mathbf{16}_{-3}\oplus\overline{\mathbf{16}}_{+3}$ is
\begin{equation}\label{eq:app-ym:Ttot}
  T_{\mathfrak m}
  =T(\mathbf{16})+T(\overline{\mathbf{16}})
  =2+2=4 ,
  \qquad
  \dim\mathfrak m=16+16=32 .
\end{equation}
\end{lemma}

\begin{proof}
$\mathfrak m$ is a single real-irreducible $H$-module of complex type
realized as the chiral spinor pair $\mathbf{16}\oplus
\overline{\mathbf{16}}$ of $\SOTen$
(\autoref{sec:foundation:tangent}). The Dynkin index is additive over a
direct sum, and conjugation preserves it, so
$T(\overline{\mathbf{16}})=T(\mathbf{16})=2$ from
\eqref{eq:app-ym:casimirs}. Hence $T_{\mathfrak m}=4$ and
$\dim\mathfrak m=32$. (The $\UOne$ charges $\pm3$ do not enter the
$\SOTen$ trace normalization; they contribute only to the Abelian
factor, treated separately.) The $4$ soldering frame modes are eaten
gauge degrees of freedom carrying no net dynamical $\SOTen$ index, so the
gauge-coupling sum is the same whether weighted by the full $32$ or by
the $28$ physical scalars: the $\SOTen$ trace of any $H$-singlet eaten
combination vanishes. We therefore use the $H$-irreducible weight
$T_{\mathfrak m}=4$. This corrects the legacy value $\tfrac52$ noted in
\autoref{rem:app-ym:multiplet}.
\end{proof}

Inserting \eqref{eq:app-ym:a4} into \eqref{eq:app-ym:W} and using
\autoref{lem:app-ym:weight},
\begin{equation}\label{eq:app-ym:insert}
  \Gamma_{\mathrm{1L}}
  \supset
  -\frac12\int_{\Lambda^{-2}}^{\infty}\frac{ds}{s}\,e^{-sm^{2}}
   \frac{1}{(4\pi)^{2}}\,\frac{s^{2}}{12}\,T_{\mathfrak m}
   \int_{X}\Tr\bigl(\mathcal F_{\mu\nu}\mathcal F^{\mu\nu}\bigr)\,d^{4}x .
\end{equation}
The proper-time integral is logarithmically divergent: with
$\int_{\Lambda^{-2}}^{\infty}ds\,s\,e^{-sm^{2}}/s
=\int_{\Lambda^{-2}}^{\infty}ds\,e^{-sm^{2}}
\;\to\;m^{-2}$ at the upper end and the cutoff supplying
$\log(\Lambda^{2}/m^{2})$ in the $s^{2}\!\cdot s^{-1}$ measure, one
obtains
\begin{equation}\label{eq:app-ym:log}
  \Gamma_{\mathrm{1L}}
  \supset
  -\frac14\,\frac{T_{\mathfrak m}}{24\pi^{2}}\,
   \log\frac{\Lambda^{2}}{m^{2}}
   \int_{X}\Tr\bigl(\mathcal F_{\mu\nu}\mathcal F^{\mu\nu}\bigr)\,d^{4}x .
\end{equation}

%----------------------------------------------------------------------%
\subsection{The induced coupling}
\label{sec:app-ym:coupling}
%----------------------------------------------------------------------%

\begin{theorem}[Induced Yang--Mills coefficient]
\label{thm:app-ym:coupling}
Integrating out the massive clock fields induces the local Yang--Mills
term
\begin{equation}\label{eq:app-ym:YM}
  S_{\mathrm{YM}}
  =-\frac{1}{4\ghat^{2}}\int_{X}
    \Tr\bigl(\mathcal F_{\mu\nu}\mathcal F^{\mu\nu}\bigr)\,d^{4}x ,
\end{equation}
with the one-loop coefficient
\begin{equation}\label{eq:app-ym:ghat}
  \ghat^{-2}
  =\frac{T_{\mathfrak m}}{24\pi^{2}}\,f^{2}\,
   \log\frac{\Lambda^{2}}{\mu^{2}}
  =\frac{4}{24\pi^{2}}\,f^{2}\,\log\frac{\Lambda^{2}}{\mu^{2}}
  =\frac{f^{2}}{6\pi^{2}}\,\log\frac{\Lambda^{2}}{\mu^{2}}
  >0 ,
\end{equation}
where $m^{2}=f^{2}\mu^{2}$ has been restored and $f^{2}$ extracted by
the canonical normalization of the kinetic term.
\end{theorem}

\begin{proof}
Match \eqref{eq:app-ym:log} to \eqref{eq:app-ym:YM}: the coefficient of
$-\tfrac14\Tr(\mathcal F^{2})$ in $\Gamma_{\mathrm{1L}}$ is
$\ghat^{-2}=T_{\mathfrak m}\log(\Lambda^{2}/m^{2})/(24\pi^{2})$ in units
where the $\sigma$-model is canonically normalized; restoring the decay
constant $f$ that sets the overall scale of the Maurer--Cartan kinetic
term multiplies this by $f^{2}$, and $\log(\Lambda^{2}/m^{2})
=\log(\Lambda^{2}/\mu^{2})+\log(1/f^{2})\to\log(\Lambda^{2}/\mu^{2})$ to
leading log. With $T_{\mathfrak m}=4$ from \autoref{lem:app-ym:weight},
$\ghat^{-2}=4f^{2}\log(\Lambda^{2}/\mu^{2})/(24\pi^{2})
=f^{2}\log(\Lambda^{2}/\mu^{2})/(6\pi^{2})$. Positivity follows from
$\Lambda>\mu$, $f^{2}>0$.
\end{proof}

\begin{remark}[Comparison with the legacy coefficient]
\label{rem:app-ym:legacy}
The corrected prefactor $T_{\mathfrak m}/(24\pi^{2})=4/(24\pi^{2})
=1/(6\pi^{2})$ replaces the legacy
$\tfrac52\cdot 1/(24\pi^{2})=5/(48\pi^{2})$, which used the spurious
$\mathbf{10}_{\pm2}\oplus\mathbf 4_{0}$ decomposition with index
$\tfrac52$. The corrected value is larger by a factor
$T_{\mathfrak m}/(5/2)=4/(5/2)=8/5$. The physical content -- a
positive, logarithmically running induced gauge coupling -- is unchanged;
only the numerical prefactor is fixed to the spinor-pair value.
\end{remark}

\begin{remark}[Normalization convention]\label{rem:app-ym:norm}
$\Tr(\mathcal F_{\mu\nu}\mathcal F^{\mu\nu})$ is normalized in the
fundamental $\mathbf{10}$ of $\SOTen$, $\Tr_{\mathbf{10}}(T^{a}T^{b})
=T(\mathbf{10})\,\delta^{ab}$ with $T(\mathbf{10})=1$. Any rescaling of
this trace rescales $\ghat^{2}$ uniformly and drops out of all
physical (renormalization-group) statements below.
\end{remark}

%----------------------------------------------------------------------%
\subsection{Feeding the asymptotic-freedom slope $b_{0}$}
\label{sec:app-ym:b0}
%----------------------------------------------------------------------%

The coefficient \eqref{eq:app-ym:ghat} sets the boundary value of the
gauge coupling at the breaking scale; its scale dependence is governed
by the full one-loop $\beta$-function of the FRG sector
(\autoref{sec:frg:gauge}). With the canonical dimensionless coupling
$\ghat^{2}$ the locked form is
\begin{equation}\label{eq:app-ym:beta}
  \beta_{\ghat^{2}}
  =2\,\ghat^{2}+\frac{b_{0}}{48\pi^{2}}\,\ghat^{4},
  \qquad
  b_{0}=\frac{11}{3}\,C_{A}-\frac{4}{3}\,T\,n_{F}
        -\frac16\,T\,n_{S} ,
\end{equation}
where the canonical $+2\,\ghat^{2}$ term enforces a \emph{Gaussian} UV
fixed point $\ghat^{2}_{\star}=0$: the gauge sector is asymptotically
\emph{free}, and there is no interacting gauge fixed point.

\begin{theorem}[Slope from the corrected matter content]
\label{thm:app-ym:b0}
With the locked Casimirs \eqref{eq:app-ym:casimirs}, the slope is
$b_{0}=26$ for one $\SOTen$ generation ($n_{F}=1$, $n_{S}=2$) and
$b_{0}=12$ for the three-generation theory of \autoref{sec:sm:generations}
($n_{F}=3$, $n_{S}=28$).
\end{theorem}

\begin{proof}
\emph{One generation.} Insert $C_{A}=8$, $T=T(\mathbf{16})=2$, $n_{F}=1$,
and $n_{S}=2$ (minimal GUT scalars at the matching scale):
\begin{equation}\label{eq:app-ym:b0one}
  b_{0}(1)
  =\frac{11}{3}\cdot 8-\frac43(2)(1)-\frac16(2)(2)
  =\frac{78}{3}=26 .
\end{equation}
\emph{Three generations.} With three chiral $\mathbf{16}$ representations
($n_{F}=3$, not $3\times16$) and $n_{S}=28$ massive coset scalars,
\begin{equation}\label{eq:app-ym:b0lin}
  b_{0}(3)
  =\frac{11}{3}(8)-\frac43(2)(3)-\frac16(2)(28)
  =\frac{88-24-28}{3}=12 .
\end{equation}
Hence $b_{0}(3\,\mathrm{gen})=12$, the value used everywhere downstream.
\end{proof}

\begin{remark}[Why the slope does not change the UV behavior]
\label{rem:app-ym:freedom}
The magnitude of $b_{0}$ controls the \emph{rate} of running but not the
existence of the UV fixed point: the canonical $+2\,\ghat^{2}$ term in
\eqref{eq:app-ym:beta} produces $\ghat^{2}_{\star}=0$. In either case the
gauge sector is asymptotically free, and the induced coefficient
\eqref{eq:app-ym:ghat} supplies a finite, positive boundary value at the
breaking scale.
\end{remark}

This closes the chain from the Maurer--Cartan curvature
\eqref{eq:app-ym:curv}, through the corrected coset weight
$T_{\mathfrak m}=4$ (\autoref{lem:app-ym:weight}), to the induced
coupling \eqref{eq:app-ym:ghat} and the asymptotic-freedom slope
$b_{0}=12$ for three generations (\autoref{thm:app-ym:b0}) carried into
the FRG analysis of \autoref{sec:frg:gauge}.
% --- end inlined: unit-app-ym.tex ---

\section{Deferred Computational Appendices}
\label{app:gravity:main}
\label{app:frg}
\label{app:ws:main}
\label{app:cy:main}
\label{app:anomaly:main}
\label{eq:MC}
\label{eq:sigma}
The heat-kernel, functional-renormalization-group, worldsheet, Calabi--Yau,
and full anomaly-arithmetic details referenced in the main text are summarized
in the sections and algebraic appendices above. This section preserves stable
cross-reference anchors for those computations.

\begin{thebibliography}{99}

\bibitem{BAMLP:2026}
B.~Belna,
Barrier-Aware Metric--Laguerre Partitioning,
manuscript, 2026.

\bibitem{Albert:1934}
A.~A.~Albert,
On a certain algebra of quantum mechanics,
\emph{Annals of Mathematics} \textbf{35} (1934).

\bibitem{Baez:2002}
J.~C.~Baez,
The octonions,
\emph{Bulletin of the American Mathematical Society} \textbf{39} (2002).

\bibitem{Freudenthal:1985}
H.~Freudenthal,
Oktaven, Ausnahmegruppen und Oktavengeometrie,
selected works on exceptional Jordan structures.

\bibitem{Georgi:1974}
H.~Georgi and S.~L.~Glashow,
Unity of all elementary-particle forces,
\emph{Physical Review Letters} \textbf{32} (1974).

\bibitem{GreenSchwarz:1984}
M.~B.~Green and J.~H.~Schwarz,
Anomaly cancellation in supersymmetric $D=10$ gauge theory,
\emph{Physics Letters B} \textbf{149} (1984).

\bibitem{GlimmJaffe:1987}
J.~Glimm and A.~Jaffe,
\emph{Quantum Physics: A Functional Integral Point of View},
Springer, 1987.

\bibitem{Helgason:1978}
S.~Helgason,
\emph{Differential Geometry, Lie Groups, and Symmetric Spaces},
Academic Press, 1978.

\bibitem{JNW:1934}
P.~Jordan, J.~von~Neumann, and E.~Wigner,
On an algebraic generalization of the quantum mechanical formalism,
\emph{Annals of Mathematics} \textbf{35} (1934).

\bibitem{Jacobson:1968}
N.~Jacobson,
\emph{Structure and Representations of Jordan Algebras},
American Mathematical Society, 1968.

\bibitem{Krutelevich:2007}
S.~Krutelevich,
Jordan algebras, exceptional groups, and higher composition laws,
\emph{Journal of Algebra} \textbf{314} (2007).

\bibitem{McCrimmon:2004}
K.~McCrimmon,
\emph{A Taste of Jordan Algebras},
Springer, 2004.

\bibitem{OsterwalderSchrader:1973}
K.~Osterwalder and R.~Schrader,
Axioms for Euclidean Green's functions,
\emph{Communications in Mathematical Physics} \textbf{31} (1973).

\bibitem{Reuter:1998}
M.~Reuter,
Nonperturbative evolution equation for quantum gravity,
\emph{Physical Review D} \textbf{57} (1998).

\bibitem{Sakharov:1967pk}
A.~D.~Sakharov,
Vacuum quantum fluctuations in curved space and the theory of gravitation,
\emph{Soviet Physics Doklady} \textbf{12} (1968).

\bibitem{Schafer:1966}
R.~D.~Schafer,
\emph{An Introduction to Nonassociative Algebras},
Academic Press, 1966.

\bibitem{Slansky:1981}
R.~Slansky,
Group theory for unified model building,
\emph{Physics Reports} \textbf{79} (1981).

\bibitem{Springer:2000}
T.~A.~Springer and F.~D.~Veldkamp,
\emph{Octonions, Jordan Algebras and Exceptional Groups},
Springer, 2000.

\bibitem{Sudbery:1984}
A.~Sudbery,
Division algebras, (pseudo)orthogonal groups and spinors,
\emph{Journal of Physics A} \textbf{17} (1984).

\bibitem{Wetterich:1993}
C.~Wetterich,
Exact evolution equation for the effective potential,
\emph{Physics Letters B} \textbf{301} (1993).

\bibitem{Yokota:2009}
I.~Yokota,
\emph{Exceptional Lie Groups},
arXiv/lecture notes, 2009.

\end{thebibliography}

\end{document}
